% =====================================================================
%  NVIDIA GPU Architectures after Volta: Lecture Notes
%  Companion notes to the 32-slide deck "NVIDIA GPU Architectures after Volta".
%  One subsection per slide: what the slide shows, the details behind
%  it, definitions of every term, and the intuition.
%
%  Compile:  pdflatex, makeindex, pdflatex, pdflatex
%  Needs only standard packages (mathptmx, tcolorbox, imakeidx, ...).
% =====================================================================
\documentclass[11pt,letterpaper]{article}
\usepackage[T1]{fontenc}
\usepackage{mathptmx}
\usepackage[scaled=0.9]{helvet}
\usepackage{courier}
\usepackage[margin=1in,headheight=14pt]{geometry}
\usepackage{amsmath,amssymb,bm}
\usepackage{booktabs,array}
\usepackage{xcolor}
\usepackage{microtype}
\usepackage{enumitem}
\usepackage{titlesec}
\usepackage[most]{tcolorbox}
\usepackage{imakeidx}
\usepackage{fancyhdr}
\PassOptionsToPackage{hyphens}{url}
\usepackage[hidelinks,bookmarksnumbered]{hyperref}

% ---------- palette (same as the deck) --------------------------------
\definecolor{ink}{HTML}{1F2937}
\definecolor{muted}{HTML}{6B7280}
\definecolor{tok}{HTML}{5F6F85}
\definecolor{emb}{HTML}{0CA678}
\definecolor{posc}{HTML}{7048E8}
\definecolor{attn}{HTML}{E8590C}
\definecolor{mlp}{HTML}{1C7ED6}
\definecolor{outc}{HTML}{D6336C}
\definecolor{trn}{HTML}{9C6644}
\definecolor{vc}{HTML}{1098AD}
\colorlet{stagec}{ink}
\newcommand{\setstage}[1]{\colorlet{stagec}{#1}}

% ---------- headings ---------------------------------------------------
\titleformat{\section}{\sffamily\bfseries\Large\color{ink}}{\thesection}{0.8em}{}[{\color{stagec}\titlerule[1.2pt]}]
\titleformat{\subsection}{\sffamily\bfseries\large\color{ink}}{\thesubsection}{0.7em}{}
\titleformat{\subsubsection}{\sffamily\bfseries\color{ink}}{\thesubsubsection}{0.6em}{}
\titlespacing*{\section}{0pt}{20pt}{10pt}
\titlespacing*{\subsection}{0pt}{16pt}{6pt}

\pagestyle{fancy}
\fancyhf{}
\fancyhead[L]{\sffamily\small\color{muted}\nouppercase{\leftmark}}
\fancyhead[R]{\sffamily\small\color{muted}\thepage}
\renewcommand{\headrulewidth}{0pt}

\makeindex[title=Index of terms,intoc,columns=2]

% ---------- slide-note machinery -------------------------------------
% \slide{number}{title}{subtitle}
\newcommand{\slide}[3]{%
  \subsection[Slide #1: #2]{Slide #1: #2}%
  {\color{muted}\itshape #3\par}\medskip}
\newcommand{\para}[1]{\medskip\noindent\textbf{#1}\quad}

\tcbset{
  termsbox/.style={breakable,enhanced,colback=stagec!4,colframe=stagec!60,
    boxrule=0.5pt,arc=2pt,left=6pt,right=6pt,top=4pt,bottom=4pt,
    title=Terms,fonttitle=\sffamily\bfseries\small,coltitle=white,colbacktitle=stagec!80,
    attach boxed title to top left={yshift=-2mm,xshift=4mm},boxed title style={arc=1.5pt,boxrule=0pt}},
  intuitionbox/.style={breakable,enhanced,colback=white,colframe=white,
    borderline west={3pt}{0pt}{stagec},left=8pt,right=2pt,top=2pt,bottom=2pt,arc=0pt,boxrule=0pt,
    title=Intuition,fonttitle=\sffamily\bfseries\small,coltitle=stagec,colbacktitle=white,
    attach boxed title to top left={yshift=-1mm,xshift=0mm},boxed title style={boxrule=0pt,left=8pt}},
  formulabox/.style={enhanced,colback=ink!4,colframe=ink!30,boxrule=0.4pt,arc=2pt,
    left=6pt,right=6pt,top=4pt,bottom=4pt}
}
\newenvironment{terms}{\begin{tcolorbox}[termsbox]\begin{description}[leftmargin=1.1em,itemsep=3pt,topsep=2pt,font=\normalfont]}{\end{description}\end{tcolorbox}}
\newcommand{\term}[2]{\item[\textbf{#1}]\index{#1} #2}
\newenvironment{intuition}{\begin{tcolorbox}[intuitionbox]}{\end{tcolorbox}}
\newenvironment{formulas}{\begin{tcolorbox}[formulabox]}{\end{tcolorbox}}

\setlist[itemize]{leftmargin=1.4em,itemsep=2pt}
\setlength{\emergencystretch}{1.5em}
\setlist[enumerate]{leftmargin=1.6em,itemsep=2pt}
\newcommand{\paper}[1]{\textit{#1}}
\newcommand{\passk}{\mathrm{pass}@k}
\newcommand{\E}{\mathbb{E}}
\DeclareMathOperator{\clip}{clip}






\begin{document}

\begin{titlepage}
\thispagestyle{empty}
\vspace*{2.5cm}
{\color{attn}\rule{3cm}{3pt}}\par\vspace{0.6cm}
{\sffamily\bfseries\color{muted}\large LECTURE NOTES \ $\cdot$ \ HARDWARE, 2026\par}
\vspace{0.4cm}
{\sffamily\bfseries\fontsize{30}{34}\selectfont\color{ink} NVIDIA GPU Architectures\\after Volta\par}
\vspace{0.3cm}
{\sffamily\fontsize{18}{22}\selectfont\color{ink} Companion notes to the slides\par}
\vspace{1.2cm}
{\large\color{muted}\parbox{0.82\textwidth}{Turing to Vera Rubin, 2018--2026: what each generation added to the streaming multiprocessor and the tensor core, the number formats from FP16 to FP4 and the block scaling that makes 4 bits usable, the memory system from HBM2 to HBM4, NVLink and NVSwitch from the 8-GPU board to the 72-GPU rack, the Grace and Vera CPUs, the fate of FP64, power and cooling, the systems built from these parts, the software that drives them, and the gap between what was announced and what shipped. One section per slide, every term defined, and a reading list for each part.}\par}
\vspace{1.5cm}
{\color{muted}For advanced undergraduates and beginning graduate students\par}
\vspace{0.3cm}
{\color{muted}November 2026\par}
\vfill
{\small\color{muted}\parbox{0.9\textwidth}{These notes follow the deck \textit{NVIDIA GPU Architectures after Volta} (32 slides) slide by slide. Numbers quoted from an architecture whitepaper, a datasheet or a paper were checked against that document; vendor peaks are quoted dense unless marked sparse, and the datasheet headline is usually the 2:4-sparse figure; numbers from product pages, press releases and press coverage are marked $\dagger$; numbers marked ``derived'' were computed from quoted ones and the formulas are given. The state of the products is as of November 8, 2026: Rubin's figures are from NVIDIA's specification page as revised in September 2026, and items NVIDIA has not disclosed are marked ``not published''.}\par}
\end{titlepage}

\tableofcontents
\clearpage

\section*{How to use these notes}
\addcontentsline{toc}{section}{How to use these notes}
Each section of these notes corresponds to one section of the deck, and each subsection to one slide. A subsection has four parts:
\begin{itemize}
  \item \textbf{What the slide shows}: a description of the figure or table, so the notes can be read without the slides.
  \item \textbf{Details}: the material the slide compresses, with the arithmetic written out and the numbers given their sources.
  \item A \textbf{Terms} box defining every technical term the slide uses, the first time it appears. All terms are collected in the index at the end.
  \item An \textbf{Intuition} box: the picture to keep in mind, and the question to ask when the picture seems to fail.
\end{itemize}
Conventions. A \emph{GPU} here is one NVIDIA data-center or graphics processor from Turing (2018) to Rubin (2026); where a package holds two dies (Blackwell, Rubin) the package is the GPU, and a \emph{superchip} is one CPU and one or two GPUs on a board. An \emph{SM} is a streaming multiprocessor; a \emph{lane} (NVIDIA: ``CUDA core'') is one FP32 pipeline in it; a \emph{tensor core} is the matrix unit, four per SM since Ampere. Peak FLOP/s counts a fused multiply-add as two operations; tensor peaks are \emph{dense} unless marked sparse (2:4 structured sparsity doubles the figure). Bandwidths use decimal prefixes (1 TB/s $=10^{12}$ B/s). A dagger ($\dagger$) marks any number whose source is a datasheet, product page, press release, benchmark list or press coverage rather than an architecture whitepaper or a paper; ``n.p.'' means not published by NVIDIA as of November 2026.

A reading plan for a single lecture: Sections 1--3 (the overview, Turing, Ampere) take about a quarter of the session; Sections 4--5 (Hopper, Ada) a quarter; Sections 6--7 (Blackwell, Rubin) a third; Sections 8--10 (systems, software, wrap-up) the rest. The exercises in Appendix~C are arithmetic on the deck's own numbers and can be set as homework. Appendix~D is the reading list, section by section, and matches the paper kit that accompanies these notes.

% =====================================================================
\setstage{ink}
\section{Overview: nine years after Volta (slides 1--3)}

\slide{1}{NVIDIA GPU Architectures after Volta}{Turing to Vera Rubin, 2018--2026}

\para{What the slide shows} The cover carries the title, a one-paragraph description of the lecture, and a schematic drawing of a GB200 NVL72 rack: eighteen compute trays drawn as orange bars, each with four small blocks for its four Blackwell GPUs, arranged in two groups of nine above and below a block of nine purple NVLink switch trays. The caption identifies the drawing and gives the rack's defining property: 72 GPUs that form a single NVLink domain, meaning that any GPU in the rack can read and write the memory of any other at the full link rate. The description paragraph lists the threads the lecture follows: what each generation added to the streaming multiprocessor and the tensor core, the number formats from FP16 to FP4, the memory system from HBM2 to HBM4, NVLink from the eight-GPU DGX-1 to the 72-GPU rack, the Grace and Vera CPUs, the systems built from these parts, and what shipped against what was announced.

\para{Details} The deck that these notes accompany picks up where \emph{Introduction to GPU Architectures} left off. That lecture ended with Volta (2017): a GV100 die of 21.1 billion transistors on 815 mm$^2$, 80 streaming multiprocessors of 64 FP32 lanes, 640 tensor cores performing a $4\times4\times4$ FP16 matrix multiply-add per clock, 16 or 32 GB of HBM2 at 900 GB/s, six NVLink 2 links for 300 GB/s per GPU, and a board power of 300 W. The V100 was the GPU of Summit, the first machine in which 97\% of the floating-point capacity sat on GPUs, and it was the GPU on which the first large transformer models were trained. Everything in this lecture is a change relative to that starting point, which is why the first content slide (slide 2) compares the V100 to the newest shipping part.

The lecture is organized in ten sections that follow the footer of every slide: Overview (slides 1--3), Turing (4--5), Ampere (6--10), Hopper (11--15), Ada (16), Blackwell (17--22), Rubin (23--26), Systems (27--28), Software (29) and Wrap-up (30--32). Eight NVIDIA architectures appear. Six are data-center compute parts in the direct line of the V100: Ampere's GA100 (the A100, 2020), Hopper's GH100 (the H100 and H200, 2022--24), Blackwell (the B200 and GB200, 2024), Blackwell Ultra (the B300 and GB300, 2025), Rubin (2026), and the announced Rubin Ultra and Feynman. Two are graphics dies that matter to the story because they carry the same tensor cores into inference and because they show what a GPU looks like when FP64, HBM and NVLink are left out: Turing (2018) and Ada Lovelace (2022). The Ampere generation had both kinds, GA100 and GA102, and slide 9 uses that pair to explain the fork.

The cover drawing is the lecture's thesis in one picture. In 2017 the unit of GPU computing was a board with one GPU; by 2024 it was a rack of 72. A GB200 NVL72 rack holds 18 compute trays, each with two Grace CPUs and four Blackwell GPUs, and 9 switch trays, each with two NVLink Switch chips. The 18 switch chips connect every GPU to every other at 1.8 TB/s, for a total of 130 TB/s of NVLink bandwidth; the rack is liquid cooled and draws about 120 kW$^\dagger$. Slide 19 returns to this rack in detail, and slide 25 compares it with its two successors, the GB300 NVL72 (2025) and the Vera Rubin NVL72 (2026). The point to carry through the lecture is that the GPU's \emph{boundary} moved: the SM stayed the unit of execution, the die stayed the unit of manufacturing, but the unit of purchase, of software scheduling and of thermal design became the rack.

A note on sources, which the conventions in the front matter spell out. NVIDIA publishes an architecture whitepaper for each compute generation through Hopper (Turing 2018, Ampere 2020, Hopper 2022, Ada 2022) and these are the primary references for SM organization, cache sizes and peak rates. For Blackwell, NVIDIA published a shorter technical brief and a developer-blog article on Blackwell Ultra; SM counts and the size of the tensor memory come from independent microbenchmark papers (Jarmusch et al., 2025). For Rubin, the sources are keynote presentations (GTC March 2025, CES January 2026, GTC March 2026), the Hot Chips 38 talks of August 2026, and NVIDIA's product specification page as revised in September 2026; several Rubin figures changed between these, and the deck gives the history where it matters (slides 3, 23 and 25). Vendor peak figures are quoted dense unless marked sparse; NVIDIA's headline numbers usually assume 2:4 structured sparsity, which doubles them (slide 8).

\begin{terms}
\term{streaming multiprocessor (SM)}{The unit of execution in an NVIDIA GPU: a block of arithmetic lanes, tensor cores, registers, shared memory and schedulers that runs a set of thread blocks. Every generation in this lecture is built from SMs of the same general shape as Volta's.}
\term{tensor core}{The matrix multiply-add unit inside an SM, introduced with Volta. It multiplies small matrix tiles at a much higher rate than the FP32 lanes and is the source of every ``AI'' peak figure.}
\term{NVLink}{NVIDIA's point-to-point interconnect between GPUs (and, from Pascal and POWER onward, between GPUs and CPUs). Its generations are numbered: NVLink 2 (Volta, 300 GB/s per GPU), 3 (Ampere, 600), 4 (Hopper, 900), 5 (Blackwell, 1,800), 6 (Rubin, 3,000 GB/s).}
\term{NVLink domain}{The set of GPUs that can reach one another over NVLink at full rate, through NVSwitch chips; 8 in a DGX, 72 in an NVL72 rack.}
\term{NVL72}{NVIDIA's rack-scale system of 72 GPUs and 36 CPUs on one NVLink fabric, shipped in GB200 (2024), GB300 (2025) and Vera Rubin (2026) versions.}
\term{compute tray}{One chassis in an NVL72 rack holding two superchips: two CPUs and four GPUs.}
\term{switch tray}{One chassis in an NVL72 rack holding two NVLink Switch chips; nine trays connect the 72 GPUs.}
\term{superchip}{NVIDIA's name for a board carrying one of its CPUs and one or two of its GPUs joined by NVLink-C2C: GH200 (Grace plus one H100), GB200 (Grace plus two Blackwell), Vera Rubin (Vera plus two Rubin).}
\term{compute die}{A GPU die designed for data-center training and HPC, with FP64 units, HBM and NVLink: GA100, GH100, Blackwell, Rubin.}
\term{graphics die}{A GPU die designed for gaming and rendering, with ray-tracing cores, GDDR memory and display engines, and only token FP64: TU102, GA102, AD102.}
\term{dense and sparse peaks}{A dense peak is the matrix throughput on ordinary operands; a sparse peak assumes 2:4 structured sparsity (two of every four weights zero) and is twice the dense figure. NVIDIA's headline figures are usually sparse.}
\end{terms}

\begin{intuition}
Read this lecture as a sequel. The first lecture explained why a GPU looks the way it does, up to Volta; this one asks what NVIDIA changed once the customer was no longer a gamer or a national laboratory but an organization training and serving neural networks by the hundred thousand GPUs. The answer is: the tensor core, the number formats, the memory, the links and the rack, in roughly that order of importance, with the SM itself nearly unchanged.
\end{intuition}

% ---------------------------------------------------------------------
\slide{2}{Nine years after Volta: what moved, and what did not}{Per-GPU tensor peaks grew about 280$\times$ by cutting the precision; FP64 grew 4$\times$ and dipped; bandwidth grew 21$\times$.}

\para{What the slide shows} A log-scale plot of peak TFLOP/s per GPU against the year of first shipment, from 2017 to 2027, with two series. The solid line is each generation's peak at its fastest tensor-core precision, dense: V100 at 125 TFLOP/s (FP16), A100 at 312 (BF16), H100 at 1,979 (FP8), GB200 at 10,000 (FP4), GB300 at 15,000 (FP4) and Rubin at 35,000 (FP4, the training figure). The dashed line is FP64 on the CUDA cores: 7.8, 9.7, 34, 40, 1.4 and 33 TFLOP/s for the same six parts, with the dip at Blackwell Ultra (B300) labeled. Two stat tiles give the headline ratios, ``280$\times$'' for tensor-peak growth and ``21 $\to$ 336 B'' for transistors per GPU, and a panel explains what the curves say.

\para{Details} The two curves summarize the lecture. Take the solid one first. From the V100's 125 TFLOP/s to Rubin's 35 PFLOP/s is a factor of 280 in nine years, about $2^{8}$, or a doubling every 13 months. That rate is far faster than transistor density improved: the V100 had 21.1 billion transistors on one die, Rubin has 336 billion on two (CES 2026 keynote slide$^\dagger$), a factor of 16. Where did the other 17$\times$ come from? Part of it is the number format. Each halving of the operand width lets a tensor core of the same area perform roughly twice as many multiply-adds per clock, because a multiplier's area grows with the square of the mantissa width and the accumulator and routing scale roughly linearly; going from FP16 to FP4 is two halvings, so the format alone accounts for about 4$\times$ (slide 10 explains why 4-bit operands are usable at all). The remaining factor of about 4 is the tensor cores themselves: Ampere doubled the matrix rate per SM per clock relative to Volta (slide 7), Hopper doubled it again (slide 11), Blackwell doubled the per-SM rate at FP4 by making FP4 a native rate, and SM counts rose from 80 to 148--160 while clocks rose from about 1.5 to about 1.9 GHz. The decomposition $280 \approx 4 \text{ (format)} \times 4 \text{ (tensor core per SM per clock)} \times 2 \text{ (SMs)} \times 1.3 \text{ (clock)} \times 3 \text{ (Rubin's undisclosed internals)}$ is approximate, because NVIDIA has not published Rubin's SM count or clock; the deck labels the factorization ``derived'' for that reason. If the 2:4-sparse figures were used instead, every point on the curve would double, and the ratio would be unchanged, because Volta had no sparsity feature and its 125 is already the only number it has.

The dashed curve is the lecture's second theme. FP64 on the CUDA cores rose from 7.8 TFLOP/s (V100) to 9.7 (A100) to 34 (H100); on the tensor cores, which Ampere and Hopper gave an FP64 mode, it reached 19.5 and 67 TFLOP/s respectively (slide 7, slide 11). Blackwell dropped the tensor-core FP64 speedup and shipped 40 TFLOP/s on the GB200 (37 on the HGX B200), slightly above Hopper's vector rate and well below its tensor rate. Blackwell Ultra then removed most of the FP64 units: 1.2--1.4 TFLOP/s per GPU, 100 TFLOP/s for a 72-GPU rack$^\dagger$. Rubin restored it to 33 TFLOP/s, with no separate tensor rate. The plot shows a curve that quadruples over nine years and then dips by a factor of 30 for one product. Slide 21 is devoted to why this happened and what HPC centers do about it; the short version is that the data-center GPU is now designed around inference first, and an inference refresh does not need double precision.

The third quantity, memory bandwidth, is in the panel rather than the plot: 0.9 TB/s on V100 (HBM2), 1.555 TB/s on the 40 GB A100 and 2.0 on the 80 GB part (HBM2 and HBM2e), 3.35 on H100 (HBM3), 8.0 on Blackwell and Blackwell Ultra (HBM3e), 19.2 on Rubin (HBM4, the revised figure; slide 23). That is a factor of 21, against 280 for arithmetic, so the ratio of the two, the ridge point of the roofline, moved by a factor of 13: from 125,000 GFLOP/s $\div$ 900 GB/s $=$ 139 FLOP per byte on V100 to 35,000,000 $\div$ 19,200 $\approx$ 1,800 FLOP per byte on Rubin (derived); the ridge in fact peaks at the GB300 (15,000,000 $\div$ 8,000 $=$ 1,875) and falls slightly on Rubin, whose HBM4 bandwidth grew faster than its FP4 rate (slide 15). Slide 15 tabulates this and draws the consequence: only large, dense matrix multiplications at low precision can use the arithmetic, and everything else a program does is bound by memory. That is why the headline feature of several products in this lecture is memory, not FLOP/s (the H200, the B300 and the HBM4 of Rubin), and why inference GPUs are now compared in tokens per second rather than TFLOP/s.

Two cautions about the plot. First, the $x$ coordinate is the year of first shipment, not of announcement; Hopper was announced in March 2022 and shipped in the fourth quarter, Blackwell was announced in March 2024 and shipped from the fourth quarter, Rubin was announced with figures in March 2025 and shipped first revenue racks in August 2026. Second, every number is a vendor peak at boost clock, which no real program sustains; the plot is a plot of what the hardware is built to do, not of what it does, and slide 29 gives the fractions of peak that measured kernels reach.

\begin{formulas}
Tensor-peak ratio, V100 to Rubin: $35{,}000 / 125 = 280$. Transistor ratio: $336 / 21.1 \approx 16$. Bandwidth ratio: $19.2 / 0.9 \approx 21$. Ridge point $I^\ast = P_{\text{peak}} / B$: V100 $125 \times 10^{12} / (0.9 \times 10^{12}) = 139$ FLOP/B; Rubin $35 \times 10^{15} / (19.2 \times 10^{12}) \approx 1{,}820$ FLOP/B. The ridge moved by $1{,}820 / 139 \approx 13$, which equals $280 / 21$.
\end{formulas}

\begin{terms}
\term{peak FLOP/s}{The maximum rate of floating-point operations a chip can perform if every arithmetic unit completes an operation every clock: units $\times$ operations per unit per clock $\times$ clock frequency. A fused multiply-add counts as two operations.}
\term{fastest tensor precision}{The narrowest number format a generation's tensor cores support at full rate: FP16 on Volta, BF16/FP16 on Ampere (INT8 and INT4 excluded as integer formats), FP8 on Hopper, FP4 on Blackwell and Rubin.}
\term{FP64}{Double-precision floating point, 64 bits with an 11-bit exponent and 52-bit mantissa; the format of most scientific computing. On NVIDIA GPUs it runs on dedicated FP64 lanes and, on Ampere and Hopper only, on the tensor cores.}
\term{HBM}{High-bandwidth memory: DRAM dies stacked vertically and connected through a very wide interface (1,024 bits per stack) on a silicon interposer next to the GPU. Generations HBM2, HBM2e, HBM3, HBM3e and HBM4 appear in this lecture.}
\term{ridge point}{In the roofline model, the arithmetic intensity (FLOP per byte of memory traffic) at which a kernel stops being memory-bound and becomes compute-bound: peak FLOP/s divided by memory bandwidth.}
\term{arithmetic intensity}{The number of floating-point operations a kernel performs per byte it moves to or from memory.}
\term{boost clock}{The highest clock frequency a GPU reaches under light load; vendor peak figures are computed at it, and sustained clocks under a power limit are lower.}
\term{year of first shipment}{The year a product reached customers in volume, as opposed to the year it was announced; this lecture dates products by shipment where the two differ.}
\end{terms}

\begin{intuition}
Picture two lines on a log plot. The top one rises steeply because the industry kept finding a shorter number it could tolerate; the bottom one rises slowly and then falls off a cliff for one product because the buyers of that product did not need it. Every other slide is a footnote to one of these two lines.
\end{intuition}

% ---------------------------------------------------------------------
\slide{3}{The cadence: every two years, then every year}{Nine products in nine years; since 2024 the announced numbers have moved before shipment.}

\para{What the slide shows} A horizontal timeline from 2017 to 2027--28 with nine markers, each labeled with the architecture, its announcement month, its first-shipment timing where that differs, and its defining features: Volta V100 (May 2017), Turing (August 2018), Ampere A100 (May 2020), Hopper H100 (March 2022, shipping in the fourth quarter), Ada (September 2022), Blackwell (March 2024, shipping in the fourth quarter), Blackwell Ultra (March 2025, shipping in July), Rubin and Vera (CES, January 2026, first racks in August) and the announced Rubin Ultra, Kyber rack and Feynman for 2027--28. Two panels follow: one on the change from a two-year to a one-year cadence, and one on the gap between announced and shipped specifications.

\para{Details} For most of the history covered by the first lecture, NVIDIA introduced a new compute architecture every two to three years: Fermi (2010), Kepler (2012), Maxwell (2014), Pascal (2016), Volta (2017). Ampere followed in 2020 and Hopper in 2022, keeping to that rhythm. NVIDIA first signaled a move to a yearly cadence to investors in October 2023 and laid the roadmap out at Computex in June 2024$^\dagger$: Blackwell in 2024, Blackwell Ultra in 2025, Rubin in 2026, Rubin Ultra in 2027, and later a new architecture named Feynman for 2028. The pattern is a new architecture in even years and an ``Ultra'' refresh in odd years. The refresh keeps the same compute dies and changes what surrounds them: taller HBM stacks for more capacity (B300: 288 GB against 186), a rebalanced tensor core (B300: 1.5$\times$ the dense FP4 rate, with FP64 and INT8 reduced), a new NIC generation and a revised rack (slide 20). A yearly cadence is possible because the refresh does not require a new die, and because the rack-scale system (slide 19) gives NVIDIA other places to improve than the GPU itself: the switch, the NIC, the CPU, the power delivery and cooling.

Each step on the timeline is now a platform rather than a chip. The Vera Rubin generation was announced at CES on January 5, 2026 as six chips designed together$^\dagger$: the Vera CPU, the Rubin GPU, the NVLink 6 switch, the ConnectX-9 network interface, the BlueField-4 data processing unit and the Spectrum-6 Ethernet switch with co-packaged optics (slide 24). A customer buys a rack or a row of racks that contains all six; the GPU's specifications are only one column of the purchase. This matters for reading the rest of the deck: several of the quantities that changed most between generations (NVLink bandwidth, rack power, scale-out bandwidth) belong to the platform, not the GPU.

The second panel raises the lecture's methodological theme. Since 2022, the numbers NVIDIA presented at announcement have differed from those on the datasheet at shipment, in both directions. Hopper's GTC 2022 whitepaper gave preliminary peaks of 30 TFLOP/s FP64, 60 FP64 on the tensor cores, 60 FP32 and 1,000 dense FP16; the shipping H100 SXM5 datasheet gives 34, 67, 67 and 989. The March 2024 Blackwell customer deck gave the GB200 45 TFLOP/s FP64, 192 GB of HBM3e at 8 TB/s and PCIe Gen 6; the shipping GB200 has 40 TFLOP/s and 186 GB, the HGX B200 37 TFLOP/s, 180 GB and 7.7 TB/s, and both use PCIe Gen 5$^\dagger$. Rubin's HBM4 bandwidth was announced as 13 TB/s per GPU at GTC in March 2025, raised to 22 TB/s at CES in January 2026, and lowered to 19.2 TB/s on NVIDIA's specification page, edited in early September 2026 (archived copies date the change to September 1--11; it was reported on October 26), where the NVLink 6 figure also fell from 3.6 to 3.0 TB/s per GPU and the rack total from 260 to 216 TB/s$^\dagger$. NVIDIA's statement of October 29, 2026 described the earlier figures as ``preliminary architectural targets'' and the new ones as what ``the system sustains at qualified, volume production''$^\dagger$. Rack-level compute peaks were not changed. Slide 25 gives the full before-and-after list.

Two lessons follow for anyone who reads GPU specifications. First, the quantities that move are the ones set late in the product cycle by what the manufacturing process yields: memory interface speeds (HBM pin rates), high-speed serial links (NVLink SerDes rates), power limits and PCIe generation. The quantities that do not move are set by the die: SM counts, tensor-core rates per clock, cache sizes. A FLOP/s figure is therefore more trustworthy at announcement than a bandwidth figure. Second, a roadmap entry two years out is a plan. In this lecture, one announced product (Rubin CPX, slide 26) disappeared from the roadmap within six months of its announcement, one rack design (Kyber) was redesigned at half its density, and one shipment date (Rubin) held but was preceded by reports of delays. The deck's convention is to use datasheet values for everything that has shipped, the September 2026 specification page for Rubin, and to mark vendor material with a dagger throughout; the notes follow the same rule.

A remark on what the timeline leaves out. Between the markers, NVIDIA shipped memory upgrades of existing dies (the 80 GB A100 in November 2020, the H200 with 141 GB of HBM3e in 2024), PCIe and dual-card variants (H100 PCIe, H100 NVL), and the graphics dies (GA102 in 2020, AD102 in 2022, GB202 in 2025) that share tensor cores with the compute dies but are sold for different purposes. The timeline shows the architectural steps; the detailed slides mention the variants where they matter.

\begin{terms}
\term{cadence}{The interval between successive architecture releases: two to three years for NVIDIA from Fermi to Hopper, one year from Blackwell onward.}
\term{Ultra refresh}{A product that keeps the previous year's die design with new memory stacks, a rebalanced tensor core and an updated rack: Blackwell Ultra (B300, 2025) reuses the Blackwell dies; the announced Rubin Ultra (2027) puts four Rubin-class dies in one package.}
\term{platform}{The set of chips NVIDIA designs and sells together for one generation: GPU, CPU, NVLink switch, NIC, DPU and Ethernet switch.}
\term{GTC}{NVIDIA's GPU Technology Conference, held each March in San Jose, where new architectures and their first specifications are usually announced.}
\term{CES}{The Consumer Electronics Show in Las Vegas each January; NVIDIA used its January 2026 keynote to declare Rubin in production.}
\term{Computex}{The Taipei trade show each June; NVIDIA laid out its yearly cadence there in 2024 (first signaled to investors in late 2023) and announced Rubin's full production there in 2026.}
\term{Hot Chips}{The annual August conference at which chip designers present microarchitecture; NVIDIA presented A100 (2020), Blackwell (2024), Blackwell Ultra (2025), and Rubin and Vera (2026) there.}
\term{preliminary specification}{A figure given at announcement before the manufacturing process is characterized; NVIDIA's October 2026 term for such figures was ``preliminary architectural targets''.}
\term{datasheet}{The vendor document stating a shipping product's specifications; this lecture's source of record for products that have shipped.}
\term{SerDes}{Serializer/deserializer, the high-speed transmitter and receiver circuits behind NVLink, PCIe and network links; their achievable data rate is one of the quantities settled late in a product cycle.}
\end{terms}

\begin{intuition}
The timeline reads like a metronome that doubled its tempo in 2024. The useful habit is to attach a date to every number: the same product can have three values for the same quantity, one from the keynote, one from the revised web page and one from the datasheet, and only the last is what a program will see.
\end{intuition}

% =====================================================================
\setstage{tok}
\section{Turing (slides 4--5)}

\slide{4}{Turing TU102 (2018): ray tracing and integer tensor cores}{A graphics die that carried Volta's SM into consumer GPUs and added two fixed-function units beside it.}

\para{What the slide shows} A block diagram of one Turing streaming multiprocessor, drawn in the same style as the Volta SM of the first lecture. Under an L1 instruction cache sit four partitions, each with a warp scheduler, a 64 KB register file, sixteen FP32 and sixteen INT32 lanes, two tensor cores and four load/store units. Below the partitions, an RT core sits beside the 96 KB L1 data cache and shared memory, and a strip for the four texture units closes the SM. The caption counts: 72 SMs on the TU102 die, 68 enabled on the GeForce RTX 2080 Ti, which gives 4,352 FP32 lanes, 544 tensor cores and 68 RT cores. Two stat tiles give 18.6 billion transistors on TSMC's 12 nm FFN process and a 754 mm$^2$ die with a 6 MB L2 and GDDR6 memory. The panel lists what Turing added.

\para{Details} Turing was announced at SIGGRAPH on August 13, 2018 as a professional product (the Quadro RTX line) and a week later, on August 20, as the GeForce RTX 20 series$^\dagger$; the RTX 2080 Ti shipped in September 2018$^\dagger$. It is a graphics architecture: TU102 is the largest of three Turing dies, and no Turing part has HBM, a full-rate FP64 unit or more than two NVLink links. It belongs in this lecture for three reasons. It was the first architecture after Volta and carried Volta's SM organization (slide 21 of the first lecture) into every consumer and workstation GPU; it added the integer modes to the tensor core that made inference a tensor-core workload; and it introduced the RT core, the only entirely new fixed-function unit NVIDIA has added to the SM since the texture unit.

The SM is Volta's with one change of balance and two additions. Volta's SM had four partitions of 16 FP32 lanes, 16 INT32 lanes, 8 FP64 lanes and 2 tensor cores; Turing's partitions keep the 16 FP32 and 16 INT32 lanes and the 2 tensor cores but reduce FP64 to a token two lanes per SM (a 1/32 rate, enough for correctness, useless for HPC). The separate INT32 datapath, which Volta introduced, lets address arithmetic and loop counters issue concurrently with floating-point work; NVIDIA's whitepaper reports that about 36 integer instructions are issued per 100 floating-point instructions in a sample of shader workloads$^\dagger$, which is the argument for spending lanes on it. The L1 data cache and shared memory are a unified 96 KB array per SM (Volta had 128 KB), configurable as 64 KB shared plus 32 KB L1 or 32 plus 64; the register file is 256 KB per SM, as on every generation from Pascal onward. The full TU102 die has 72 SMs in 6 GPCs, 4,608 lanes, 576 tensor cores, 72 RT cores and a 6 MB L2; the RTX 2080 Ti enables 68 SMs and 5.5 MB of L2, the Quadro RTX 6000 and 8000 and the Titan RTX enable all 72$^\dagger$.

The tensor cores are the second generation. Volta's tensor core computed a $4\times4\times4$ FP16 matrix multiply-add per clock, 64 fused multiply-adds, with FP32 or FP16 accumulation. Turing's does the same for FP16 and adds INT8 at twice the rate and INT4 at four times, with INT32 accumulation, plus an experimental INT1 (binary) mode at sixteen times$^\dagger$. The per-SM rate for FP16 is unchanged from Volta: 8 tensor cores $\times$ 64 FMA $=$ 512 FMA per clock. On the RTX 2080 Ti at its 1,545 MHz boost clock$^\dagger$ that gives $68 \times 512 \times 2 \times 1.545 \times 10^9 = 107.6$ TFLOP/s; NVIDIA's whitepaper quotes 113.8 TFLOP/s for the Founders Edition at its higher 1,635 MHz boost, with FP16 accumulation, and 56.9 with FP32 accumulation, since the consumer parts run FP32 accumulation at half rate$^\dagger$. The deck rounds to 114 TFLOP/s FP16 and 228 TOPS INT8 at the Founders Edition boost clock (the whitepaper gives 227.7 TOPS INT8 and 455.4 TOPS INT4)$^\dagger$. The integer modes exist for inference. A network trained in FP32 or FP16 can be quantized to INT8 weights and activations with a calibration pass that picks a scale per tensor, usually with less than one percent loss of accuracy; TensorRT 5, released alongside Turing, automated this$^\dagger$. INT4 is harder to use and saw little adoption; its successor idea, 4-bit floating point with block scaling, is Blackwell's (slide 18).

The RT core is a unit that accelerates the two inner loops of ray tracing: traversing a bounding-volume hierarchy (a tree of nested boxes enclosing the scene's triangles) and testing a ray against a box or a triangle. In a rasterization pipeline, visibility is decided per pixel by the depth test; in ray tracing, each pixel casts a ray into the scene and asks what it hits first, and reflections, shadows and global illumination cast further rays from the hit points. Without hardware, the traversal costs thousands of instructions per ray on the SM's lanes, with irregular memory access and divergence. The RT core performs the traversal and intersection tests as a fixed-function operation that the SM launches and waits on, freeing the lanes for shading; NVIDIA claimed about 10 billion rays per second for the RTX 2080 Ti against about 1 billion for a Pascal GPU doing the work in software$^\dagger$. Microsoft's DirectX Raytracing API (announced March 2018) and Vulkan's extensions expose it; it arrived in games in late 2018. The RT core has no role in AI or HPC, and the compute dies (GA100, GH100, Blackwell) omit it; it appears here because it is the clearest example of what a graphics die spends its area on that a compute die does not (slide 16).

Turing's memory is GDDR6 on a 352-bit bus on the 2080 Ti (384-bit on the full die), 11 GB at 616 GB/s$^\dagger$. Two x8 NVLink 2 links, 100 GB/s total, connect two consumer cards, the first NVLink on a GeForce product$^\dagger$. The die is large, 754 mm$^2$, close to Volta's 815, on the same TSMC 12 nm FFN process; a graphics die this large had not been made before, and the RTX 20 series was expensive for it. Graphics features that were new with Turing and that appear in later generations are mesh shaders (a programmable replacement for the vertex and geometry stages), variable-rate shading (shading at lower resolution where the eye will not notice) and DLSS, the first use of tensor cores in graphics, which the next slide takes up.

\begin{formulas}
Tensor FP16 peak, RTX 2080 Ti (FP16 accumulate): $68 \text{ SMs} \times 8 \text{ TC} \times 64 \text{ FMA} \times 2 \text{ FLOP} \times f$. At $f = 1.545$ GHz: 107.6 TFLOP/s; at the Founders Edition's 1.635 GHz: 113.8 TFLOP/s (the whitepaper's figure). INT8: twice; INT4: four times. FP32 peak: $4{,}352 \times 2 \times 1.635 = 14.2$ TFLOP/s.
\end{formulas}

\begin{terms}
\term{Turing}{NVIDIA's 2018 GPU architecture (dies TU102, TU104, TU106 and smaller), the first after Volta; a graphics architecture with RT cores and integer tensor cores.}
\term{TU102}{The largest Turing die: 18.6 billion transistors, 754 mm$^2$, 72 SMs, 6 MB L2, a 384-bit GDDR6 interface; used in the RTX 2080 Ti, Titan RTX and Quadro RTX 6000/8000.}
\term{RT core}{A fixed-function unit, one per SM from Turing on, that traverses bounding-volume hierarchies and tests ray--box and ray--triangle intersections for ray tracing.}
\term{bounding-volume hierarchy (BVH)}{A tree of nested boxes enclosing a scene's triangles, used to find which triangles a ray might hit without testing all of them.}
\term{ray tracing}{Rendering by casting rays from the camera into the scene and from hit points toward lights and other surfaces; physically faithful but expensive compared with rasterization.}
\term{INT8, INT4}{8-bit and 4-bit integer operand formats for the tensor cores, introduced with Turing at two and four times the FP16 rate, used for quantized inference.}
\term{quantization}{Converting a trained network's weights and activations from floating point to a narrower format, usually integers with a per-tensor scale, for faster inference.}
\term{TensorRT}{NVIDIA's inference compiler and runtime, which fuses layers, selects kernels and performs INT8 (later FP8 and FP4) quantization with calibration.}
\term{GDDR6}{Graphics DDR memory, the sixth generation: individual chips on the board connected over 32-bit channels; 14 Gb/s per pin on Turing, 616 GB/s over 352 bits on the RTX 2080 Ti.}
\term{12 nm FFN}{TSMC's 12 nm FinFET process variant used for Volta and Turing, a refinement of 16 nm rather than a full node shrink.}
\term{mesh shader}{A programmable pipeline stage introduced with Turing that generates geometry in groups, replacing the fixed vertex and geometry stages.}
\term{variable-rate shading}{Shading some screen regions at lower resolution than others, trading quality where the eye is insensitive for speed.}
\term{Quadro RTX, Titan RTX}{The workstation and prosumer Turing products with the full 72-SM TU102 die and 24 or 48 GB of GDDR6.}
\end{terms}

\begin{intuition}
Turing is Volta's SM wearing a gaming jersey: the FP64 units are gone, two new fixed-function helpers (the RT core and the integer tensor modes) stand beside it, and GDDR6 replaces HBM. Keep it in mind as the template for every graphics die that follows, and as the moment the tensor core found a second customer, inference, that would soon outgrow the first.
\end{intuition}

% ---------------------------------------------------------------------
\slide{5}{The T4: inference at 70 watts, and tensor cores in every game}{Turing's data-center part was a low-profile PCIe card; its consumer part ran a neural network on every frame.}

\para{What the slide shows} A table comparing three NVIDIA inference cards: the Tesla P4 of 2016 (Pascal GP104), the Tesla T4 of 2018 (Turing TU104) and the V100 PCIe of 2017 (GV100) by FP32 lanes, tensor cores, FP32 and FP16-tensor TFLOP/s, INT8 and INT4 TOPS, memory type and size, bandwidth and board power. The T4 has 2,560 lanes and 320 tensor cores, 8.1 TFLOP/s FP32, 65 TFLOP/s FP16, 130 TOPS INT8 and 260 TOPS INT4, 16 GB of GDDR6 at 300 GB/s and a 70 W power limit; the P4 has 22 TOPS of INT8 at 50--75 W and no tensor cores; the V100 PCIe has 112 TFLOP/s FP16 but no INT8 tensor mode, 900 GB/s of HBM2 and 250 W. A panel below the table draws the lesson; a panel to the right covers DLSS, the T4 microbenchmarks and the split between compute and graphics dies.

\para{Details} The Tesla T4 was announced on September 12, 2018 at GTC Japan$^\dagger$. It is the data-center product of the Turing generation, built on TU104, the second-largest Turing die (13.6 billion transistors, about 545 mm$^2$ according to secondary sources; NVIDIA's whitepaper does not give these figures for TU104$^\dagger$). The card enables 40 of TU104's 48 SMs: 2,560 lanes and 320 tensor cores, with 16 GB of GDDR6 and a 4 MB L2$^\dagger$. Its defining specification is the power limit: 70 W, delivered entirely through the PCIe slot, in a half-height, half-length, passively cooled form factor. Every server with a free PCIe slot could take one; most could take several. NVIDIA's datasheet quotes 8.1 TFLOP/s FP32, 65 TFLOP/s FP16 on the tensor cores, 130 TOPS INT8 and 260 TOPS INT4, and 300 GB/s of memory bandwidth$^\dagger$ (some NVIDIA web pages give 320 GB/s; the datasheet figure is used here).

The table's comparison explains why the T4 mattered. Against the V100 PCIe, the data-center flagship of the time, the T4 has half the lanes and half the tensor cores, a third of the memory bandwidth and less than a third of the power. On FP16 tensor arithmetic it reaches about 58\% of the V100's rate (65 against 112 TFLOP/s); on INT8 it has a mode the V100 lacks, since Volta's tensor cores are FP16-only (the V100 can do INT8 dot products on its CUDA cores through the DP4A instruction, at a far lower rate). Against the P4, its predecessor in the same form factor, it has six times the INT8 throughput at the same power. For inference the comparison that matters is not peak against peak but throughput at the batch sizes a serving system actually runs. A request arriving at a server is processed alone or in a small batch, so the matrix multiplications are thin, the arithmetic intensity is low, and the kernel is bound by memory bandwidth and latency rather than by peak FLOP/s; a card with a third of the bandwidth loses less than a third of its throughput, and a card with a fraction of the price and power wins on throughput per dollar and per watt. TensorRT 5 added INT8 calibration that converted FP32 models with little accuracy loss$^\dagger$, so the 130 TOPS was usable. The cloud providers adopted the T4 widely (Amazon's G4 instances in 2019, Google Cloud's T4 instances) and video-transcoding fleets used it for its encoders$^\dagger$; it became the standard inference accelerator of 2019--2021 until the A10 and L4, its Ampere and Ada successors, replaced it.

The T4 also served as a laboratory. Jia, Maggioni, Smith and Scarpazza, who had microbenchmarked the V100 in 2018, repeated their measurements on the T4 in 2019 (arXiv:1903.07486): instruction latencies, cache sizes and latencies for the unified 96 KB L1 and the 4 MB L2, the throughput and numerical behavior of the INT8 and INT4 tensor modes, and the clock's behavior under the 70 W limit, which a sustained tensor-core workload exceeds, so the card throttles below its boost clock. The paper is the reference for anyone who wants to know what the Turing SM does, as opposed to what its datasheet says, and it is the second in a series the deck cites for every generation (Abdelkhalik et al. for Ampere, Luo et al. for Hopper, Jarmusch et al. for Blackwell; slide 15).

The right panel turns to the consumer side. DLSS (deep learning super sampling) renders a frame at reduced resolution and uses a convolutional neural network, running on the tensor cores, to upscale it to the display resolution. The first version (February 2019) trained a network per game and was judged blurry$^\dagger$; DLSS 2.0 (March 2020) replaced it with a single temporal network that takes the low-resolution frame, the motion vectors the game engine already computes, and the previous output frame, and produces an image that in many cases matches or exceeds native rendering$^\dagger$. The significance for this lecture is architectural rather than graphical: a unit introduced in Volta for training neural networks in data centers became, within three years, a standard stage of the real-time graphics pipeline, executed for every frame on every RTX card. Its successors run on later graphics dies: frame generation on Ada (slide 16), ray reconstruction (2023), and multi-frame generation on Blackwell's consumer parts. The tensor core, in other words, found three markets in sequence, training, inference and graphics, which is why every NVIDIA die made since 2018 has one, and why the format and rate of the tensor core is the quantity that changes most from generation to generation.

The panel's last point sets up slide 9 and slide 16. From Turing onward NVIDIA builds two kinds of die per generation. The compute die (GA100, GH100, Blackwell, Rubin) has full-rate FP64, HBM, many NVLink links, no RT cores and no display outputs, and is sold for training and HPC. The graphics die (TU102, GA102, AD102, GB202) has RT cores, GDDR memory, display and video engines, token FP64 and little or no NVLink, and is sold for gaming, rendering, visualization and inference. Both have the generation's tensor cores. Inference, the workload that cares about throughput per dollar at modest memory sizes, buys both: the L40S (slide 16) and the L4, graphics dies in server form, serve the same market the T4 opened.

\begin{terms}
\term{Tesla T4}{Turing's data-center inference card (2018): TU104 with 40 SMs, 320 tensor cores, 16 GB GDDR6 at 300 GB/s, 70 W, half-height PCIe.}
\term{Tesla P4}{Pascal's inference card (2016): GP104 with 2,560 lanes, 22 TOPS INT8 via DP4A, 8 GB GDDR5, 50--75 W.}
\term{TU104}{The second-largest Turing die (48 SMs), used in the T4, the RTX 2080 and the Quadro RTX 5000.}
\term{inference}{Running a trained neural network on new inputs, as opposed to training it; typically small batches, bound by memory bandwidth and latency.}
\term{batch size}{The number of inputs processed together in one pass; large batches make matrix multiplications fat and compute-bound, small ones make them thin and memory-bound.}
\term{DP4A}{A Pascal-era CUDA-core instruction computing a four-element INT8 dot product with INT32 accumulation; the pre-tensor-core path to INT8 inference.}
\term{calibration}{The step in post-training quantization that chooses the scale factors mapping floating-point tensors to INT8, by running representative inputs through the network.}
\term{DLSS}{Deep learning super sampling: rendering at reduced resolution and upscaling with a neural network on the tensor cores. Version 1 (2019) was per-game; version 2 (2020) is a single temporal network using motion vectors.}
\term{temporal upscaling}{Reconstructing a high-resolution frame from the current low-resolution frame and previous frames, aligned using motion vectors.}
\term{microbenchmark}{A small program that measures one hardware property (a latency, a throughput, a cache size) by timing a controlled instruction sequence, revealing what datasheets omit.}
\term{power limit}{The board power a GPU's firmware enforces by lowering clocks; sustained tensor-core workloads on the T4 run below boost clock for this reason.}
\term{half-height PCIe card}{A low-profile expansion card that fits in 1U and 2U servers; the T4's form factor, powered from the slot alone.}
\end{terms}

\begin{intuition}
The T4 is the proof that inference is a different problem from training: a card with a third of the bandwidth and a quarter of the power of the flagship could serve requests nearly as well, because requests arrive one at a time and the flagship's arithmetic sits idle. DLSS is the proof that a matrix unit, once present, finds uses its designers did not plan. Both lessons recur through the rest of the deck.
\end{intuition}

% =====================================================================
\setstage{emb}
\section{Ampere (slides 6--10)}

\slide{6}{Ampere GA100 (2020): the A100 die}{May 14, 2020: 54.2B transistors on TSMC 7 nm, 108 of 128 SMs enabled, 40 MB of L2, HBM2 at 1.6 TB/s.}

\para{What the slide shows} A schematic floorplan of the GA100 die. Eight GPC blocks, each containing sixteen small bars for its SMs, are arranged four above and four below two large blocks labeled ``L2 (20 MB)''; six HBM2 stacks sit at the left and right edges, one of them drawn dashed and labeled ``disabled''; a line along the top names the NVLink 3 interface (12 links, 600 GB/s), PCIe 4.0 and the GigaThread engine. The caption explains that each bar is one SM, 128 on the die and 108 enabled on the A100, that the L2 is two 20 MB halves joined by a crossbar, and that five of the six HBM2 stacks are active for 40 GB (80 GB of HBM2e later). Two stat tiles give 54.2 billion transistors on TSMC N7 and an 826 mm$^2$ die; a panel summarizes what is new.

\para{Details} The A100 was announced on May 14, 2020, in a GTC keynote recorded in Jensen Huang's kitchen during the pandemic, and shipped immediately in the DGX A100 (slide 9)$^\dagger$. It is the first compute die in this lecture and the most consequential: for the three years until the H100 shipped in volume it was the GPU on which nearly every large language model was trained, and its features (TF32, BF16, structured sparsity, MIG, the third-generation tensor core, NVLink 3) set the pattern for everything after it.

The die. GA100 is manufactured on TSMC's N7 process, the first 7 nm NVIDIA part, and carries 54.2 billion transistors on 826 mm$^2$, slightly larger than Volta's 815 mm$^2$ and again near the reticle limit of about 850 mm$^2$ that a single exposure can print. It is organized as 8 GPCs (graphics processing clusters), each of 8 TPCs (texture processing clusters) of 2 SMs: 128 SMs on the full die. The A100 product enables 108 of them, 7 GPCs' worth plus some, a yield margin of 16\%; every data-center die from Volta onward ships with SMs disabled for this reason (Volta 80 of 84, Hopper 132 of 144). Each SM has 64 FP32 lanes, 32 FP64 lanes, 64 INT32 lanes and four tensor cores (slide 7), so the A100 has $108 \times 64 = 6{,}912$ FP32 lanes and $108 \times 4 = 432$ tensor cores. The L2 cache is 40 MB, 6.7 times Volta's 6 MB, split into two halves each serving the GPCs and memory controllers on its side of the die, with a crossbar between the halves; a memory request from an SM is served by the near half when the address maps there and crosses to the far half otherwise. NVIDIA quotes 2.3 times the L2 read bandwidth of the V100$^\dagger$. The memory interface has six HBM2 controllers for six stacks; the A100 activates five, giving 40 GB (five 8 GB stacks) at 1,555 GB/s, and the sixth position is a yield spare that also lets the package hold a sixth stack mechanically. The 80 GB A100, announced in November 2020$^\dagger$, uses HBM2e stacks of 16 GB at a higher pin rate for 2,039 GB/s; it kept five active stacks.

The interfaces. NVLink 3 doubles the per-link rate to 50 GB/s bidirectional (25 GB/s each way, as on Volta) by halving the lane count per link while doubling the signaling rate, and the A100 has 12 links for 600 GB/s per GPU, twice Volta's 300$^\dagger$. The host interface is PCIe 4.0 x16 at 32 GB/s each way, twice Volta's PCIe 3.0. The SXM4 module is rated at 400 W, up from Volta's 300; a PCIe card version is rated at 250 W and runs lower clocks$^\dagger$.

The new features, which the next three slides take in turn, are of three kinds. In the tensor core: TF32 and BF16 input formats, an FP64 mode, and support for 2:4 structured sparsity at twice the dense rate (slides 7 and 8). In the memory system: the large split L2 with residency controls, an asynchronous copy instruction from global to shared memory, hardware compression of sparse data, and a 192 KB L1/shared memory per SM (slide 8). In partitioning: multi-instance GPU, which divides the A100 into up to seven isolated GPUs with their own SMs, L2 slices and memory (slide 8). The announcement also introduced the HGX A100 baseboard, eight SXM4 modules and six NVSwitches, which server vendors integrate; the DGX A100 is NVIDIA's own server built on it (slide 9).

Two sources. NVIDIA's Ampere architecture whitepaper (2020) is the primary reference, with the SM diagram, the peak-rate table (its Table 4) and the feature descriptions; Choquette and Gandhi presented the design at Hot Chips 32 in August 2020 and in \paper{IEEE Micro} in 2021. Abdelkhalik et al.\ (2022) microbenchmarked the A100 and report instruction latencies and the measured behavior of the tensor cores; the deck cites them on slide 15.

\begin{terms}
\term{Ampere}{NVIDIA's 2020 architecture; the compute die is GA100 (A100), the graphics dies GA102 and smaller (GeForce RTX 30 series).}
\term{GA100}{The Ampere compute die: 54.2 billion transistors, 826 mm$^2$, TSMC N7, 128 SMs (108 enabled on the A100), 40 MB L2, six HBM2 controllers, 12 NVLink 3 links.}
\term{A100}{The product built on GA100: 40 GB HBM2 (1,555 GB/s) or 80 GB HBM2e (2,039 GB/s), SXM4 at 400 W or PCIe at 250 W.}
\term{GPC, TPC}{Graphics processing cluster and texture processing cluster: the die's hierarchy above the SM. On GA100, a TPC is two SMs and a GPC is eight TPCs.}
\term{reticle limit}{The largest die a lithography stepper can expose in one shot, about 850 mm$^2$ on current tools; GV100, GA100 and GH100 are all near it, and Blackwell and Rubin use two such dies.}
\term{yield margin}{The fraction of a die's units left disabled so that dies with manufacturing defects can still be sold: 16 of 128 SMs and one of six HBM stacks on the A100.}
\term{split L2}{A level-2 cache divided into two halves, each close to half the SMs and memory controllers, joined by a crossbar; introduced with GA100 at 40 MB.}
\term{HBM2e}{An extended HBM2 generation with 16 GB stacks and higher pin rates; the 80 GB A100 uses five stacks for 2,039 GB/s.}
\term{SXM4}{The mezzanine module form factor of the A100, mounted on an HGX baseboard, rated at 400 W; successors are SXM5 (H100) and the Blackwell modules.}
\term{HGX}{NVIDIA's eight-GPU baseboard with NVSwitches, sold to server makers; HGX A100, HGX H100, HGX B200.}
\term{PCIe 4.0}{The fourth generation of the host bus, 16 GT/s per lane, 32 GB/s each way over 16 lanes; the A100's host link.}
\term{NVLink 3}{Ampere's NVLink generation: 50 GB/s per link, 12 links, 600 GB/s per GPU.}
\end{terms}

\begin{intuition}
GA100 is Volta scaled by the rules of the first lecture (more SMs on a new process, a much larger L2, faster links) plus one idea: make the tensor core speak more languages. Hold the floorplan in mind as the standard compute-die layout; Hopper's is the same picture with different numbers, and Blackwell's is two of them side by side.
\end{intuition}

% ---------------------------------------------------------------------
\slide{7}{The A100 SM and the third-generation tensor core}{Twice Volta's matrix rate per SM per clock, with new input formats: TF32 for free, BF16, and FP64 on the tensor cores.}

\para{What the slide shows} A block diagram of one A100 SM, after Choquette and Gandhi's Hot Chips 2020 presentation. Below the L1 instruction cache are four partitions, each with a warp scheduler, a 64 KB register file, sixteen FP32 and sixteen INT32 lanes, eight FP64 lanes, one tensor core and eight load/store units; below the partitions sits the 192 KB L1 data cache and shared memory (configurable to 164 KB of shared memory) with four texture units. The caption totals the SM: 64 FP32, 64 INT32 and 32 FP64 lanes, four tensor cores of 256 dense FP16 fused multiply-adds per clock (1,024 per SM, twice GV100's 512), 256 KB of registers, 2,048 resident threads. The panel derives the peak rates and explains the new formats.

\para{Details} The A100 SM keeps Volta's four-partition structure and Volta's lane counts; what changes is the tensor core. Volta and Turing had two tensor cores per partition, each computing a $4\times4\times4$ FP16 product per clock (64 FMA). Ampere has one tensor core per partition that computes an $8\times4\times8$ product per clock, 256 FMA: four tensor cores per SM give 1,024 FP16 FMA per clock, twice Volta's 512. The larger tile matters as much as the rate. A matrix unit's efficiency comes from reusing each operand element across many products; an $8\times8$ output tile accumulated over an inner dimension of 4 reads $8\times4 + 4\times8 = 64$ input elements and performs 256 FMAs, four FMAs per input element, against two for Volta's $4\times4\times4$ tile. Halving the operand traffic per FMA is what lets the rate double without doubling the register bandwidth; the same logic drives Hopper's warpgroup-wide instructions (slide 13) and Blackwell's tensor memory (slide 18).

The peak rates follow from the counts and the 1,410 MHz boost clock. FP32 on the lanes: $108 \times 64 \times 2 \times 1.41 \times 10^9 = 19.5$ TFLOP/s. FP64 on the lanes: $108 \times 32 \times 2 \times 1.41 = 9.7$ TFLOP/s, half the FP32 rate, the traditional HPC ratio. FP16 or BF16 on the tensor cores: $108 \times 1{,}024 \times 2 \times 1.41 = 312$ TFLOP/s, 2.5 times Volta's 125, with FP32 accumulation at full rate (unlike the consumer Turing parts). INT8 at twice that, 624 TOPS, and INT4 at 1,248. Each of these doubles with 2:4 sparsity (slide 8), which is where the datasheet's 624 TFLOP/s FP16 and 1,248 TOPS INT8 come from. TF32 runs at half the FP16 rate, 156 TFLOP/s dense, because a TF32 multiplier is about twice the size of an FP16 one. FP64 on the tensor cores runs at 64 FMA per clock per SM, 19.5 TFLOP/s, exactly twice the FP64 lanes; this is the first time an NVIDIA GPU's double-precision peak came from the matrix unit, and the DGEMM that HPL runs uses it, which is why the A100's HPL efficiency is high.

TF32 is the format that made the A100's tensor cores useful to programs that never asked for them. It has FP32's 8-bit exponent and FP16's 10-bit mantissa, 19 bits in all, and exists only inside the tensor core: the inputs in memory and registers are ordinary FP32, the tensor core rounds them to TF32 as it reads them, multiplies, and accumulates in full FP32. A program whose FP32 matrix multiplies are routed to TF32 therefore gets the tensor cores without any change to its source, at 156 TFLOP/s against 19.5 for the FP32 lanes, an 8$\times$ increase. Who routes them matters: cuDNN (from version 8.0) uses TF32 by default for FP32 convolutions, and the deep-learning frameworks and NVIDIA's NGC containers (PyTorch 1.7, TensorFlow 2.4) turned TF32 on by default for their matrix multiplies as well, while cuBLAS itself keeps FP32 as its default math mode and uses TF32 only when the caller opts in through \texttt{cublasSetMathMode} or the \texttt{NVIDIA\_TF32\_OVERRIDE} environment variable$^\dagger$. So most framework-based FP32 training sped up on the day the A100 shipped (PyTorch exposes a switch to turn it off), while an HPC code calling cuBLAS directly did not change unless it asked to. The cost is that the products are computed from 10-bit mantissas, the same as FP16; for neural-network training this is harmless, since FP16 training with FP32 accumulation was already standard, but it is not FP32 arithmetic, and a program that relies on single-precision accuracy in a matrix multiply must disable it. HPC codes that run in FP32 were caught out by this more than once.

BF16, the ``brain floating point'' format Google introduced for TPUs, has FP32's 8-bit exponent and a 7-bit mantissa. It runs at the FP16 rate on the A100's tensor cores. Its advantage over FP16 is range: FP16's 5-bit exponent spans about $6 \times 10^{-5}$ to $65{,}504$, so gradients in training underflow unless the loss is multiplied by a scale factor before backpropagation and divided out afterward (loss scaling, from Micikevicius et al.\ 2017, the first lecture's slide 23); BF16 spans the FP32 range and needs no loss scaling, at the cost of three fewer mantissa bits, which training tolerates. Within two years BF16 displaced FP16 as the training format on NVIDIA hardware, and FP8 (slide 12) was then added as a matrix-multiply format underneath it.

The slide's last point is the first appearance of the lecture's second theme. With FP64 on the tensor cores, the A100's HPC peak is a matrix-unit number, 19.5 TFLOP/s, and codes whose inner loops are dense linear algebra (HPL, many solvers) get it, while codes that are not (stencils, particle methods, sparse solvers) get the 9.7 TFLOP/s of the lanes at best. Hopper kept the arrangement and doubled both numbers (34 and 67); Blackwell gave up the 2$\times$: NVIDIA still lists an FP64 tensor-core figure for the B200 and GB200, but it equals the vector figure (37 and 40 TFLOP/s), the new \texttt{tcgen05} instructions have no FP64 mode and the legacy path remains; Blackwell Ultra then removed most of the lanes. Slide 21 tells that story, and this slide is where to notice that NVIDIA had, in 2020, chosen to serve HPC through the matrix unit.

\begin{formulas}
Peak at $f = 1.41$ GHz with 108 SMs:
FP32 lanes $108 \times 64 \times 2 f = 19.5$ TFLOP/s; FP64 lanes $108 \times 32 \times 2 f = 9.7$;
FP16/BF16 tensor $108 \times 1{,}024 \times 2 f = 312$ (624 sparse); TF32 tensor $108 \times 512 \times 2 f = 156$ (312 sparse);
FP64 tensor $108 \times 64 \times 2 f = 19.5$; INT8 $624$ TOPS (1,248 sparse); INT4 $1{,}248$ (2,496 sparse).
Operand reuse of a tensor tile $M\times N\times K$: FMAs per input element $= MNK / (MK + KN)$; $4\times4\times4$ gives 2, $8\times4\times8$ gives 4.
\end{formulas}

\begin{terms}
\term{third-generation tensor core}{Ampere's matrix unit: one per SM partition, 256 FP16 FMA per clock ($8\times4\times8$), with TF32, BF16, FP64, INT8, INT4 and binary modes and 2:4 sparsity support.}
\term{TF32}{TensorFloat-32: a 19-bit format with an 8-bit exponent and 10-bit mantissa used only inside the tensor core; FP32 inputs are rounded to it on the fly and products accumulate in FP32. The default matrix path for single-precision deep-learning code in cuDNN and the frameworks from Ampere on; opt-in in cuBLAS.}
\term{BF16}{Brain floating point: 16 bits with an 8-bit exponent and 7-bit mantissa; FP32's range at FP16's cost, so no loss scaling is needed. The standard training format from 2020.}
\term{FP16}{Half precision: 1 sign, 5 exponent, 10 mantissa bits; range about $6\times10^{-5}$ to 65,504.}
\term{loss scaling}{Multiplying the training loss by a constant so that gradients stay within FP16's range, then dividing the weight updates by it; needed for FP16 training, not for BF16.}
\term{FP32 accumulation}{Summing a matrix product's partial results in 32-bit floating point even when the inputs are 16- or 8-bit; standard on the data-center tensor cores and the reason reduced-precision training converges.}
\term{FP64 tensor core}{The mode, on Ampere and Hopper only, in which the tensor core performs double-precision matrix multiply-add at twice the rate of the FP64 lanes: 19.5 TFLOP/s on the A100, 67 on the H100.}
\term{operand reuse}{The number of multiply-adds performed per input element fetched; larger matrix tiles raise it and lower the register and memory bandwidth per FLOP.}
\term{cuBLAS, cuDNN}{NVIDIA's libraries for dense linear algebra and for neural-network primitives; the places where tensor-core modes become defaults.}
\term{DGEMM}{Double-precision general matrix multiply; the kernel behind the HPL benchmark and the one that uses the FP64 tensor cores.}
\end{terms}

\begin{intuition}
Think of the third-generation tensor core as a matrix unit that learned to accept more input types without asking the programmer: hand it FP32 and it quietly computes in TF32 at eight times the speed; hand it BF16 and it trains without loss scaling; hand it FP64 and it doubles the HPC peak. Each of those was a one-line change in a library and a large change in what people did with the hardware.
\end{intuition}

% ---------------------------------------------------------------------
\slide{8}{Sparsity, multi-instance GPU, and the A100 memory system}{Three features for the data center: a 2$\times$ from pruned weights, seven GPUs in one, and a bigger, smarter L2.}

\para{What the slide shows} Two diagrams and two panels. The upper diagram shows 2:4 structured sparsity: a row of twelve weight cells of which six are shaded (nonzero) and six are nearly white (zero), two zeros in every group of four, with an arrow down to the compressed form of six value cells plus six small index cells, each index two bits. The lower diagram shows multi-instance GPU: seven equal slices of an A100, each with 14 SMs, a slice of the L2 and 5 GB of memory. The left panel says why these three features exist; the right panel gives the details of sparsity, MIG, the memory system and the HBM2 configuration.

\para{Details} Structured sparsity is a hardware answer to an observation about neural networks: after training, many weights can be set to zero with little effect on accuracy if the network is then fine-tuned briefly. Unstructured pruning, zeroing any 50\% of weights, is hard to exploit in hardware because the pattern is irregular. Ampere's tensor cores support one specific pattern, 2:4: in every group of four consecutive weights along the inner dimension of the matrix product, exactly two are zero. A pruned weight matrix is stored compressed, as the two nonzero values of each group plus a 2-bit index saying which two positions they came from, so the weights take half the memory plus a 12.5\% overhead for indices. The tensor core reads the compressed weights and the dense activations, uses the indices to select the matching activation elements, and skips the multiplications by zero, completing the group in half the clocks: the matrix rate doubles. NVIDIA's whitepaper reports that a prune-and-retrain recipe (train dense, prune to 2:4, fine-tune with the sparsity mask) recovered accuracy on dozens of vision and language networks$^\dagger$. The feature is exposed through the cuSPARSELt library and later through TensorRT.

The convention the slide stresses is this: every ``sparse'' peak on an NVIDIA datasheet since 2020 means 2:4 structured sparsity, and is exactly twice the dense figure. The A100 datasheet reads ``312 TFLOP/s $|$ 624 TFLOP/s'' for FP16; the H100's ``3,958 TFLOP/s FP8'' headline is the sparse figure for a dense 1,979; the GB200 NVL72's ``1.44 exaflops'' is sparse FP4 for a dense 720 petaflops. The sparse rate is achievable only on a network pruned to the pattern; dense training, and most inference, runs at the dense rate. The deck quotes dense figures throughout and marks the exceptions.

Multi-instance GPU addresses a different customer. A cloud provider renting GPU time, or an inference service running many small models, cannot use a 400 W A100 for one job at a time; before MIG, sharing a GPU between processes was possible through time slicing or the multi-process service, but one tenant's memory traffic or faults could slow another's. MIG partitions the A100 at the hardware level into up to seven instances. Each instance receives its own SMs (14 for the smallest instance, so seven instances use 98 of the 108 SMs), its own slices of the L2 cache and its own memory controllers and DRAM (5 GB per smallest instance on the 40 GB part, 10 GB on the 80 GB part), its own address space and its own fault isolation; instances can also be combined (two, three, four or seven slices). The result is that an instance's throughput does not depend on what the others are doing, which is the quality-of-service property a cloud needs to sell fractions of a GPU, and the reason the feature is called an ``instance'' rather than a partition. Hopper kept MIG and added confidential-computing isolation to it (slide 11).

The memory system changed in three ways besides the 40 MB L2 of slide 6. First, L2 residency controls let a program ask that a range of addresses be kept in L2 (``persisting'' accesses) while streaming accesses pass through, so a working set such as a neural network's weights can be pinned while activations flow by. Second, hardware compression of sparse and repeated data in L2 and DRAM raises effective bandwidth for suitable data$^\dagger$. Third, and most important for programming, the asynchronous copy instruction (\texttt{cp.async} in PTX, exposed as \texttt{memcpy\_async} in CUDA 11) copies data from global memory directly into shared memory without staging it through registers. On Volta, a thread that wanted to stage a tile of a matrix into shared memory had to load it into registers and then store it to shared memory, occupying registers and lanes for the duration; with \texttt{cp.async} the load proceeds in the background while the lanes compute on the previous tile, and a completion mechanism tells the thread when the data has arrived. This is the first step of a sequence the lecture follows: Hopper's TMA moves the address generation into hardware too (slide 13), and Blackwell's tensor memory removes the accumulators from the register file (slide 18). Each step takes work away from the threads so the tensor cores can be fed.

The HBM2 details close the slide. The five active stacks of 8 GB run at 1,215 MHz (2.43 Gb/s per pin) over 5,120 bits for 1,555 GB/s, 1.7 times the V100's 900 GB/s on the same HBM2 technology at higher clocks and one more stack. The sixth controller and stack position are a manufacturing spare. The 80 GB part replaced the stacks with 16 GB HBM2e at 3.2 Gb/s per pin for 2,039 GB/s; NVIDIA rounds both figures to 1.6 and 2.0 TB/s. The 192 KB combined L1 data cache and shared memory per SM (up to 164 KB shared) is 1.5 times Volta's 128 KB.

\begin{formulas}
2:4 sparsity storage: $N$ weights become $N/2$ values plus $N/2$ two-bit indices, $= 0.5 N$ values $+ 0.125 N$ bytes of indices for 16-bit values, about 56\% of the dense size; compute: $2\times$ the dense rate.
A100 HBM2 bandwidth: $5 \text{ stacks} \times 1{,}024 \text{ bits} \times 2.43 \text{ Gb/s} / 8 = 1{,}555$ GB/s; HBM2e at 3.2 Gb/s: $2{,}048$ GB/s (datasheet 2,039).
MIG: 7 instances $\times$ 14 SMs $= 98$ of 108 SMs; 7 $\times$ 5 GB $= 35$ of 40 GB.
\end{formulas}

\begin{terms}
\term{structured sparsity (2:4)}{A weight pattern in which exactly two of every four consecutive weights are zero, stored as two values plus two-bit indices; Ampere's tensor cores run such matrices at twice the dense rate.}
\term{pruning}{Setting a fraction of a trained network's weights to zero, followed by fine-tuning to recover accuracy.}
\term{sparse peak}{A datasheet rate that assumes 2:4 structured sparsity; twice the dense rate, achievable only on pruned networks.}
\term{cuSPARSELt}{NVIDIA's library for matrix multiplication with 2:4 structured-sparse operands.}
\term{multi-instance GPU (MIG)}{Hardware partitioning of an A100 or later data-center GPU into up to seven isolated instances, each with its own SMs, L2 slices, memory and address space.}
\term{quality of service}{The guarantee that one tenant's work does not slow another's; the property MIG provides and time-slicing does not.}
\term{L2 residency control}{A CUDA 11 feature letting a program mark address ranges as persisting in L2 while other accesses stream through.}
\term{asynchronous copy (cp.async)}{An Ampere instruction that copies data from global memory into shared memory without passing through registers, overlapping with computation; exposed as memcpy\_async.}
\term{shared memory}{The programmer-managed on-chip memory of an SM, shared by a thread block; 164 KB of the 192 KB L1/shared array on the A100.}
\term{compute data compression}{Ampere's hardware compression of sparse or repeated values in the L2 and memory path, raising effective bandwidth.}
\term{pin rate}{The data rate of one memory interface signal, in Gb/s; HBM bandwidth is pin rate $\times$ bits per stack $\times$ stacks.}
\end{terms}

\begin{intuition}
Three customers, three features. The inference buyer wanted more throughput for the same chip and got sparsity; the cloud wanted to sell a GPU seven times over and got MIG; the kernel writer wanted the tensor cores fed without spending threads on it and got \texttt{cp.async}. The last is the beginning of a line that runs through Hopper's TMA to Blackwell's tensor memory.
\end{intuition}

% ---------------------------------------------------------------------
\slide{9}{Ampere deployed: DGX A100, early systems, and GA102}{Eight A100s behind six NVSwitches became the unit of AI infrastructure; the graphics die went its own way.}

\para{What the slide shows} A diagram of the DGX A100 server: eight A100 boxes in a row, each connected by thin lines to six NVSwitch boxes beneath, with two AMD EPYC 7742 CPU boxes and a note on PCIe 4.0 and eight 200 Gb/s InfiniBand adapters at the bottom. The caption explains that each GPU has two NVLink 3 links to each switch, that every GPU reaches every other at the full 600 GB/s, that the system holds 320 or 640 GB of HBM, that the advertised ``5 PFLOP/s'' is eight times the 624 TFLOP/s sparse FP16 figure, and that it draws 6.5 kW. A panel below names the first A100 systems; the panel on the right describes GA102, the Ampere graphics die, and the fork between compute and graphics dies.

\para{Details} The DGX A100 was announced with the A100 on May 14, 2020 and shipped at once, the first time an NVIDIA architecture launched with a complete system$^\dagger$. Its design follows the DGX-2 of 2018 (the first lecture's slide 25) rather than the DGX-1: all GPU-to-GPU traffic goes through NVSwitch chips, so that any pair of GPUs communicates at the full NVLink rate regardless of position. The HGX A100 baseboard carries eight SXM4 modules and six second-generation NVSwitch chips; each GPU's twelve NVLink 3 links are divided two per switch, and each switch has 36 ports, so the six switches together provide 4.8 TB/s of bisection bandwidth, enough for all eight GPUs to send at 600 GB/s simultaneously$^\dagger$. The host side is two 64-core AMD EPYC 7742 processors (NVIDIA's first use of AMD CPUs, chosen for PCIe 4.0 support that Intel then lacked), 1 TB of system memory, PCIe 4.0 switches, and eight Mellanox ConnectX-6 adapters at 200 Gb/s HDR InfiniBand, one per GPU, so that GPUs on different servers can communicate directly through GPUDirect RDMA. The system draws up to 6.5 kW and cost \$199,000 at launch$^\dagger$. Its headline ``5 petaflops of AI performance'' is $8 \times 624$ TFLOP/s, the 2:4-sparse FP16 figure; the dense figure is 2.5 PFLOP/s, and the FP64 figure is $8 \times 19.5 = 156$ TFLOP/s on the tensor cores or 78 on the lanes.

The DGX A100 and the HGX A100 baseboard became the unit in which AI infrastructure was bought and the module in which the first large clusters were built. NVIDIA's own Selene, assembled in 2020 as a DGX SuperPOD (the reference design of 140 DGX A100 systems on a rail-optimized InfiniBand fabric, in which the $k$-th adapter of every server connects to the same leaf switch so that GPU $k$ of every node is one hop from GPU $k$ of every other node), reached about 63 PFLOP/s on HPL and fifth place on the November 2020 TOP500 list$^\dagger$. National centers followed: Perlmutter at NERSC, whose GPU partition of A100 nodes entered the list in fifth place in June 2021, and Leonardo at CINECA in Italy, fourth in November 2022, both with four A100s per node connected by NVLink and HDR InfiniBand between nodes$^\dagger$. The hyperscalers' clusters of this generation, Microsoft's for OpenAI among them, were built from the same HGX baseboards in larger numbers but were not submitted to the TOP500; the pattern of a public HPC list and a larger private AI fleet begins here (slide 27).

The right panel turns to the other Ampere. GA102 is the die of the GeForce RTX 3090 and 3080 (announced September 1, 2020) and of the professional RTX A6000 and the data-center A40 and A10$^\dagger$. It differs from GA100 in nearly every respect except the name. It is made by Samsung on its 8 nm process, not by TSMC; it has 28.3 billion transistors on 628 mm$^2$. Its SM has the Turing arrangement of four partitions with 16 FP32 and 16 INT32 lanes, but with a change: one of the two 16-wide datapaths in each partition can execute either INT32 or FP32 instructions, so a partition can issue 32 FP32 operations per clock when there is no integer work, and NVIDIA counts 128 ``CUDA cores'' per SM. The RTX 3090 enables 82 of 84 SMs for 10,496 lanes and 35.6 TFLOP/s FP32 at its 1.7 GHz boost clock, 1.8 times the A100's FP32 rate on a smaller, cheaper die. Its tensor cores are third-generation like GA100's, with sparsity, but there are 328 of them (four per SM), and FP32 accumulation runs at half rate; the RT cores are second-generation. What it lacks tells the story: FP64 at 1/64 of the FP32 rate (two lanes per SM, for correctness only), GDDR6X memory instead of HBM (24 GB on a 384-bit bus at 936 GB/s, less than half the A100's bandwidth and without ECC), a 6 MB L2 instead of 40 MB, no MIG, and NVLink 3 only on the 3090 and only as a two-card bridge. It has, instead, display outputs, video encoders and decoders and RT cores, and a 350 W board power.

The fork that the slide names runs through the rest of the lecture. From Ampere on, NVIDIA designs each generation twice: a compute die for data centers (GA100, GH100, Blackwell B100/B200, Rubin) and a graphics die for gaming and workstations (GA102, AD102, GB202). The compute die spends its area on FP64 units, HBM controllers, a large L2 and many NVLink links; the graphics die spends it on RT cores, display and video engines, and (from Ada) a very large L2 to compensate for GDDR's narrow bus. Both carry the generation's tensor cores, which is why the graphics die is also an inference product: the A10 and A40, and later the L40S (slide 16), serve the market the T4 opened. The two lines also run on different foundries and processes at times (Samsung 8 nm for GA102, TSMC N7 for GA100), and are announced months apart.

\begin{terms}
\term{DGX A100}{NVIDIA's eight-A100 server (May 2020): HGX A100 baseboard with six NVSwitches, two AMD EPYC 7742 CPUs, eight 200 Gb/s InfiniBand adapters, 6.5 kW.}
\term{NVSwitch (second generation)}{The Ampere-generation NVLink crossbar chip with 36 ports of NVLink 3; six of them give eight A100s full-bandwidth all-to-all connectivity.}
\term{bisection bandwidth}{The bandwidth across a cut dividing a network in half; 4.8 TB/s for the DGX A100's NVSwitch fabric.}
\term{EPYC 7742}{AMD's 64-core server processor (Zen 2, 2019) used in the DGX A100 for its PCIe 4.0 support.}
\term{InfiniBand HDR}{The 200 Gb/s generation of InfiniBand (2018--19); the DGX A100 has one adapter per GPU.}
\term{GPUDirect RDMA}{Transfers between a network adapter and GPU memory that bypass the CPU and system memory.}
\term{DGX SuperPOD}{NVIDIA's reference design for a cluster of DGX systems on a rail-optimized InfiniBand fabric; 140 DGX A100 in the 2020 design.}
\term{rail-optimized fabric}{A network topology in which adapter $k$ of every node connects to the same leaf switch, so that the $k$-th GPUs of all nodes are one hop apart; used for data-parallel collectives.}
\term{Selene}{NVIDIA's in-house DGX A100 SuperPOD (2020), about 63 PFLOP/s on HPL, fifth on the November 2020 TOP500.}
\term{Perlmutter, Leonardo}{A100-based national systems at NERSC (Berkeley, 2021) and CINECA (Bologna, 2022).}
\term{GA102}{The Ampere graphics die: Samsung 8 nm, 28.3 billion transistors, 628 mm$^2$, 84 SMs of 128 FP32 lanes, 6 MB L2, GDDR6X; used in the RTX 3090/3080, RTX A6000, A40 and A10.}
\term{dual FP32 datapath}{The GA102 SM partition's arrangement in which one 16-lane path is FP32-only and the other executes FP32 or INT32, giving 32 FP32 operations per clock per partition when no integer work is present.}
\term{GDDR6X}{Micron's graphics memory using PAM4 signaling for higher pin rates; 936 GB/s over 384 bits on the RTX 3090.}
\term{compute/graphics fork}{NVIDIA's practice since Ampere of designing a data-center die and a graphics die per generation, differing in FP64, memory, L2, NVLink and RT cores.}
\end{terms}

\begin{intuition}
Two pictures to keep. One: eight GPUs behind a row of switches, every pair at full speed; copy this box a hundred times and connect the copies with InfiniBand and you have the 2020--22 AI cluster. Two: GA100 and GA102 side by side, same tensor cores, different everything else; the compute die buys bandwidth and precision, the graphics die buys ray tracing and price. Both pictures reappear, scaled up, in every later generation.
\end{intuition}

% ---------------------------------------------------------------------
\slide{10}{The precision ladder: FP32 to FP4}{Every format trades mantissa bits for throughput; below 8 bits a shared scale per block of values keeps the range.}

\para{What the slide shows} Eight bit-layout bars, one per number format, each divided into a sign bit, exponent bits and mantissa bits (drawn to a common scale down to 16 bits; FP32's 23-bit mantissa is truncated), with the exponent and mantissa widths and the generation that first ran the format on its tensor cores: FP32 (8 exponent, 23 mantissa; the CUDA cores of every generation), TF32 (8, 10; Ampere), BF16 (8, 7; Ampere), FP16 (5, 10; Pascal at double rate, Volta's tensor cores), FP8 E4M3 (4, 3; Hopper), FP8 E5M2 (5, 2; Hopper), FP6 E3M2 (3, 2; Blackwell) and FP4 E2M1 (2, 1; Blackwell, with NVFP4 block scaling). A note adds that INT8 and INT4 (Turing) are the integer rungs and that FP64 on tensor cores exists only on Ampere and Hopper. The panel explains how to read the ladder.

\para{Details} A floating-point format with $e$ exponent bits and $m$ mantissa bits represents numbers of the form $\pm 1.f \times 2^{x}$, where $f$ has $m$ bits and $x$ ranges over about $2^{e}$ values centered on zero (with the extremes reserved for zero, subnormals, infinities and NaN in IEEE-style formats). The exponent width sets the dynamic range: 8 bits give about $10^{\pm 38}$, 5 bits give $6 \times 10^{-5}$ to $65{,}504$, 4 bits (E4M3) give a largest value of 448, 2 bits (E2M1) give a largest value of 6. The mantissa width sets the relative precision: $m$ bits give a relative rounding error of $2^{-(m+1)}$, about $6 \times 10^{-8}$ for FP32, $5 \times 10^{-4}$ for FP16 and TF32, $4 \times 10^{-3}$ for BF16, $6 \times 10^{-2}$ for E4M3 and 25\% for E2M1. The ladder in the slide is a sequence of decisions about which of these to give up.

The 16- and 19-bit formats keep enough of both to serve as storage formats. BF16 and TF32 keep FP32's exponent and cut the mantissa, which suits neural-network training, where the values span many orders of magnitude but need little relative precision because the stochastic gradient descent process is itself noisy; FP16 keeps a 10-bit mantissa and cuts the exponent, which gives better precision but requires loss scaling to keep gradients in range. The 8-bit formats are matrix-multiply formats rather than storage formats: Hopper's Transformer Engine (slide 12) keeps the model in BF16 or FP32 and casts tensors to FP8 immediately before a matrix multiply, with a scale factor chosen per tensor so the values fill the representable range. Two FP8 encodings exist because weights and activations want precision (E4M3, three mantissa bits, range 448) while gradients want range (E5M2, two mantissa bits, range 57,344); the choice was standardized in the paper by Micikevicius et al.\ (2022) from NVIDIA, Arm and Intel.

Each rung down the ladder roughly doubles the tensor-core rate. The reason is the multiplier: the area of a floating-point multiplier is dominated by the $m \times m$ partial-product array, so halving the mantissa quarters that array, and the exponent adder and the accumulator (kept in FP32 in all cases) scale more slowly; within the same area budget a tensor core can place about twice as many FP8 multipliers as FP16 ones and about twice as many FP4 as FP8. Hopper runs FP8 at twice its FP16 rate (1,979 against 989 TFLOP/s dense on the H100); Blackwell runs FP4 at twice its FP8 rate (10 against 5 PFLOP/s dense on the GB200); FP6 runs at the FP8 rate on Blackwell because the datapath is the FP8 one. TF32 runs at half the FP16 rate, since its 10-bit mantissa needs the same multiplier as FP16 but the inputs are 32-bit and the tile is narrower. The accumulation is in FP32 throughout (FP16 or BF16 accumulation is offered on some parts for inference), so the sums are exact to about $10^{-7}$ relative even when the products are coarse.

Below 8 bits a value by itself carries too little information to be useful. E2M1 has sixteen codes: $0, \pm 0.5, \pm 1, \pm 1.5, \pm 2, \pm 3, \pm 4, \pm 6$ (the sixteenth code is a second zero). A tensor whose entries range over $10^{4}$ cannot be represented by scaling once per tensor, since most entries would round to zero or saturate. The remedy, called microscaling or block scaling, attaches a scale factor to each small block of values rather than to the whole tensor. The Open Compute Project's MX specification (2023; Rouhani et al.) defines MXFP8, MXFP6 and MXFP4 with blocks of 32 elements sharing one E8M0 scale, an 8-bit power of two, so each block costs 8/32 = 0.25 extra bits per element. NVIDIA's NVFP4 (2025) uses blocks of 16 elements with an E4M3 scale, which is not restricted to powers of two and so halves the quantization error of the scale, plus one FP32 scale for the whole tensor, costing 8/16 = 0.5 extra bits per element, 4.5 bits per value in all. Slide 18 covers the hardware that applies these scales inside the matrix unit and the accuracy results.

The slide's last line records who uses which rung as of 2026. Training keeps its master weights and optimizer state in FP32, its activations and gradients in BF16, and its matrix multiplies in FP8 (Hopper and later), with NVFP4 training demonstrated at scale in 2025 (slide 18). Inference stores weights in FP8 and increasingly in FP4, since the weights are read once per token and halving them halves the dominant memory traffic. HPC stays in FP64, with mixed-precision techniques (iterative refinement, which solves in low precision and corrects in FP64) where the problem allows; the TOP500's HPL-MxP benchmark measures exactly this. The integer formats INT8 and INT4 (Turing) remain in use for inference on graphics dies and for emulating FP64 arithmetic (slide 21), but on the data-center parts they have been overtaken by FP8 and FP4, and Blackwell Ultra reduced its INT8 rate sharply (slide 20).

\begin{formulas}
Largest finite value of a format with $e$ exponent bits and $m$ mantissa bits (IEEE-style bias $2^{e-1}-1$): $(2 - 2^{-m}) \times 2^{2^{e-1}-1}$; FP16: 65,504; E5M2: 57,344; E4M3 (NVIDIA's variant, which reclaims the infinity codes): 448; E2M1: 6.
Relative rounding error: $2^{-(m+1)}$. Bits per value with block scaling: $b + s/B$ for $b$-bit elements, an $s$-bit scale and block size $B$: MXFP4 $4 + 8/32 = 4.25$; NVFP4 $4 + 8/16 = 4.5$ (plus one 32-bit scale per tensor).
\end{formulas}

\begin{terms}
\term{exponent, mantissa}{The two fields of a floating-point number: the exponent sets the power of two (the range), the mantissa the significant digits (the relative precision).}
\term{dynamic range}{The ratio of the largest to the smallest positive representable value; set by the exponent width.}
\term{relative precision}{The largest relative rounding error of a format, $2^{-(m+1)}$ for an $m$-bit mantissa.}
\term{FP8 E4M3, E5M2}{The two 8-bit floating-point formats introduced with Hopper: four exponent and three mantissa bits (range 448, for weights and activations), or five and two (range 57,344, for gradients).}
\term{FP6}{A 6-bit floating-point family on Blackwell (E3M2 and E2M3, both accepted by the fifth-generation tensor cores), running at the FP8 rate.}
\term{FP4 E2M1}{The 4-bit floating-point format of Blackwell and Rubin: one sign, two exponent and one mantissa bit; sixteen codes, largest value 6; usable only with block scaling.}
\term{microscaling (MX)}{The Open Compute Project's 2023 family of block-scaled formats (MXFP8, MXFP6, MXFP4, MXINT8) in which 32 elements share one 8-bit power-of-two (E8M0) scale.}
\term{NVFP4}{NVIDIA's 4-bit format (2025): blocks of 16 E2M1 values with an E4M3 scale per block and an FP32 scale per tensor; 4.5 bits per value.}
\term{block scaling}{Attaching a scale factor to each small block of values so that narrow formats can represent tensors with a wide range of magnitudes.}
\term{per-tensor scaling}{One scale factor for a whole tensor, as in FP8 training on Hopper; sufficient at 8 bits, insufficient at 4.}
\term{subnormal}{A floating-point value below the smallest normal exponent, represented with a leading zero; extends the range toward zero at reduced precision.}
\term{iterative refinement}{Solving a linear system in low precision and correcting the result with residuals computed in FP64 until it reaches FP64 accuracy; the basis of the HPL-MxP benchmark.}
\term{HPL-MxP}{The TOP500's mixed-precision benchmark, which solves the HPL problem with iterative refinement and reports an FP64-equivalent rate.}
\end{terms}

\begin{intuition}
Each rung of the ladder is the same bargain: give up half the bits, get twice the multiplies, and find a trick that makes the lost bits not matter. For 16 bits the trick was FP32 accumulation and loss scaling; for 8 bits it was a scale per tensor chosen on the fly; for 4 bits it is a scale per sixteen values. When someone announces a 2-bit format, ask what the new trick is.
\end{intuition}

% =====================================================================
\setstage{attn}
\section{Hopper (slides 11--15)}

\slide{11}{Hopper GH100 (2022): the H100}{March 2022, shipping Q4 2022: 80B transistors on TSMC 4N, 132 of 144 SMs, HBM3 at 3.35 TB/s, 700 W.}

\para{What the slide shows} A schematic floorplan of the GH100 die fills the left half: eight graphics processing clusters (GPCs) drawn as eight blocks of eighteen small bars, one bar per SM, arranged four above and four below a pair of blue rectangles labeled ``L2 (25 MB)'' that together make the 50 MB level-two cache; three HBM3 stacks stand along each edge of the die, one of the six drawn dashed and labeled ``disabled''; a line across the top names the external interfaces, NVLink 4 (18 links, 900 GB/s), PCIe 5.0 and the GigaThread engine. The caption counts the SMs (144 on the die, 132 enabled on the SXM5 module, 114 on the PCIe card) and the cache (60 MB on the full die, 50 MB enabled). Two stat tiles give 80 billion transistors on TSMC's 4N process and an 814 mm$^2$ die with 16{,}896 FP32 lanes. The panel ``What is new on the die'' lists the SM changes, the memory system and the programming-model additions that the next slides expand.

\para{Details} Hopper is the data-center architecture NVIDIA announced at GTC on March 22, 2022, declared in ``full production'' on September 20, 2022, and shipped in volume from the fourth quarter of that year$^\dagger$. The compute die, GH100, is built on TSMC's 4N process, a customized 5 nm-class node, and holds 80 billion transistors on 814 mm$^2$$^\dagger$. Like GA100 before it (slide 6), it is a near-reticle-limit monolithic die: the 826 mm$^2$ of GA100 and the 814 mm$^2$ of GH100 are both within a few percent of the largest rectangle a lithography stepper can expose, which is why the next generation had to go to two dies (slide 17).

The die is organized as 8 GPCs of 9 texture processing clusters (TPCs), each TPC holding two SMs, for 144 SMs; the H100 SXM5 product enables 132 of them and the H100 PCIe card 114$^\dagger$. The difference is yield management: a die with a defect in one TPC is sold as a product that does not need that TPC, and the enabled count is chosen so that enough dies pass. With 128 FP32 lanes per SM the SXM5 part has $132\times128=16{,}896$ lanes, and with four tensor cores per SM it has 528 tensor cores. Each SM also has 64 FP64 lanes, twice Ampere's 32, so that FP64 throughput per SM per clock doubled along with FP32.

The SM's storage grew in one place: the combined L1 data cache and shared memory is 256 KB per SM, of which up to 228 KB can be configured as shared memory, against 192 KB (164 KB shared) on A100 and 128 KB (96 KB) on V100. The register file stayed at 256 KB per SM (64 KB per partition, 65{,}536 32-bit registers), as it has since Pascal; slide 15 makes the point that per-SM storage barely moved while the matrix units grew eight-fold. The L2 is 50 MB enabled (60 MB on the full die), split into two partitions with a crossbar between them as on A100; a thread's accesses are served by the near partition when the data is there and go across otherwise, which is one of the reasons the microbenchmark papers measure two L2 latencies on these parts.

Memory: the SXM5 module carries 80 GB of HBM3 in five active stacks (a sixth stack position is a yield spare, as on A100) at 3.35 TB/s; the PCIe card uses HBM2e at 2.0 TB/s$^\dagger$. Links: NVLink 4 provides 18 links at 50 GB/s each, 900 GB/s per GPU, 1.5 times A100's 600 GB/s; the host interface is PCIe 5.0 x16 at 128 GB/s bidirectional. The SXM5 module is rated at 700 W, the PCIe card at 350 W$^\dagger$.

The peaks are worth quoting with care, because the numbers moved between announcement and shipment and because NVIDIA's headline figures are the sparse ones. The GTC 2022 whitepaper gave preliminary peaks of 30 TFLOP/s FP64, 60 TFLOP/s FP64 on the tensor cores, 60 TFLOP/s FP32 and 1{,}000 TFLOP/s dense FP16; the shipping datasheet gives 34, 67, 67 and 989$^\dagger$. Dense tensor peaks for the SXM5 part are TF32 494, FP16/BF16 989, FP8 1{,}979 TFLOP/s and INT8 1{,}979 TOPS; with 2:4 structured sparsity each doubles (989, 1{,}979, 3{,}958, 3{,}958), and it is the doubled figures that appear in most marketing material$^\dagger$. The PCIe card scales these down (26 TFLOP/s FP64, 51 FP32, 756 dense FP16, 1{,}513 dense FP8) with its fewer SMs and lower clock$^\dagger$.

Two derivations show how the peaks relate to the die. The FP32 peak is $16{,}896 \text{ lanes} \times 2 \text{ FLOP} \times f$; with $f = 1.98$ GHz, the SXM5 boost clock, this is $66.9$ TFLOP/s, the datasheet's 67. The dense FP16 tensor peak, if each SM did 2{,}048 FMA per clock (twice A100's 1{,}024, as the whitepaper states), would be $132\times2{,}048\times2\times1.98\,\text{GHz} = 1{,}070$ TFLOP/s; the datasheet's 989 corresponds to the same rate at 1.83 GHz. The two headline numbers therefore do not sit at one clock: either the tensor peak is quoted at a lower sustained clock than the FP32 peak, or the per-SM matrix rate is slightly under 2{,}048 FMA per clock. The whitepaper does not resolve this, and it is a good example of why a peak should always be quoted together with the clock it assumes (derived).

The programming-model additions are the subject of slides 12 and 13: a Transformer Engine that runs matrix multiplies in FP8 with automatic scaling; thread block clusters, distributed shared memory and the tensor memory accelerator; and asynchronous transaction barriers. Two smaller additions are also on the die. DPX instructions accelerate the inner loops of dynamic-programming algorithms (minimum, maximum and add on pairs of values with fused comparisons), which NVIDIA says speed Smith--Waterman sequence alignment and Floyd--Warshall shortest paths by up to 7$\times$ over A100$^\dagger$. Confidential computing lets a virtual machine run on an H100 with its memory encrypted and its state attested, so that a cloud tenant's model and data are protected from the host; it is the first GPU to offer this. The second-generation multi-instance GPU partitioning (slide 8) continues, now with confidential computing per instance.

\begin{terms}
\term{Hopper}{NVIDIA's data-center GPU architecture of 2022, named after Grace Hopper; the compute die is GH100 and the products are the H100 (SXM5 and PCIe), H200 and the GH200 superchip.}
\term{GH100}{The Hopper compute die: 80 billion transistors, 814 square millimeters, TSMC 4N, 144 SMs on the die.}
\term{H100 SXM5}{The full-power H100 module: 132 SMs, 80 GB HBM3 at 3.35 TB/s, 700 W, NVLink 4; the part whose peaks are usually quoted.}
\term{TSMC 4N}{The customized 5 nm-class TSMC process on which GH100 and AD102 are built; Blackwell's 4NP is a further tuning of it, not a node shrink.}
\term{GPC}{Graphics processing cluster: the group of SMs (18 on GH100) that shares a raster engine and, on Hopper, hosts a thread block cluster.}
\term{TPC}{Texture processing cluster: a pair of SMs sharing texture units; the granularity at which SMs are disabled for yield.}
\term{reticle limit}{The largest die a lithography stepper can expose in one shot, about 830 square millimeters; GA100 and GH100 are within a few percent of it.}
\term{HBM3}{The third generation of stacked high-bandwidth memory; five stacks give the H100 SXM5 80 GB at 3.35 TB/s.}
\term{NVLink 4}{Hopper's GPU-to-GPU link: 18 links at 50 GB/s each, 900 GB/s per GPU.}
\term{DPX instructions}{Hopper instructions for dynamic programming (fused minimum, maximum and add on value pairs), used in sequence alignment and shortest-path algorithms.}
\term{confidential computing}{Running a workload with its memory encrypted and its state attested so that neither the host operating system nor the cloud operator can read it; first on a GPU with Hopper.}
\term{dense and sparse peaks}{A dense peak counts every multiply-add; a sparse peak assumes 2:4 structured sparsity in one operand and is twice the dense figure. NVIDIA headlines are usually sparse.}
\term{boost clock}{The highest clock a product runs at under its power and thermal limits; datasheet peaks assume it, sustained workloads usually run below it.}
\end{terms}

\begin{intuition}
Hopper is Ampere with every arithmetic unit doubled per SM, a third more SMs, twice the link bandwidth, and a new way to feed the matrix units. Hold the floorplan in mind as two halves of L2 with SMs around them and HBM on the edges, because every later NVIDIA die, including the two-die Blackwell, is that picture repeated.
\end{intuition}

% ---------------------------------------------------------------------
\slide{12}{Transformer Engine and FP8: the first 8-bit training format}{Two FP8 encodings, per-tensor scales from recent magnitudes, FP32 accumulation: twice the FP16 rate.}

\para{What the slide shows} A flow diagram across the top left traces one matrix multiply through a Transformer Engine layer: inputs arrive in BF16 or FP32, a scale factor is chosen from the recent history of maximum magnitudes, the operands are cast to FP8 (E4M3 or E5M2), the FP8 tensor-core GEMM runs, the result is accumulated in FP32, unscaled and written out, and the next layer receives BF16. A caption gives the ranges of the two encodings (448 for E4M3, 57{,}344 for E5M2) and notes that the per-tensor scale is updated every step. Below it, a panel ``What the numbers say'' gives the H100's dense and sparse FP16 and FP8 peaks, the 6.3$\times$ ratio of H100 FP8 to A100 BF16, and the result of Micikevicius et al.\ (2022). The right panel explains why 8 bits work with scaling, the two encodings, the library that does the work, and where FP8 went afterwards.

\para{Details} An 8-bit floating-point number has one sign bit and seven bits to divide between exponent and mantissa. Hopper's tensor cores accept two divisions. E4M3 has a 4-bit exponent and a 3-bit mantissa; its largest finite value is 448 and its smallest subnormal is $2^{-9}\approx 0.002$, and the format gives up the IEEE convention of reserving an exponent code for infinities in order to extend the range by one binade (only NaN is reserved). E5M2 has a 5-bit exponent and a 2-bit mantissa; it is FP16 with the low eight mantissa bits removed, so its range is FP16's ($\pm 57{,}344$ as the largest finite value) and it keeps infinities and NaNs. Three mantissa bits mean a relative precision of about $2^{-4}$, or 6\%, per value; two bits mean 12.5\%. Compared with FP16 (10 mantissa bits, relative precision $2^{-11}$) or BF16 (7 bits), these are coarse numbers, and the surprise of 2022 was that a neural network could be trained with them.

Why the precision suffices is a statistical fact about training: the matrix multiplies sum many products, the errors of individually rounded products are approximately independent and average out in the FP32 accumulator, and stochastic gradient descent tolerates noise in gradients. Why the range is the real problem is a fact about the distributions: the activations in one layer may span $10^{-3}$ to $10^{2}$ and the gradients $10^{-8}$ to $10^{0}$, and a format whose range is 448 (E4M3) cannot hold either without rescaling. FP16, with a range to 65{,}504 and subnormals down to $6\times10^{-8}$, already needed loss scaling (slide 23 of the previous deck): the loss is multiplied by a constant before the backward pass so that small gradients stay above the underflow threshold, and the weight update divides it back out. FP8 needs a scale per tensor, not one per training run.

The Transformer Engine's recipe is delayed scaling. For every tensor that will be cast to FP8 (the weights, the activations entering a matrix multiply, the gradients leaving one), the library records the maximum absolute value (``amax'') observed in each of the last several iterations, takes the largest, and computes a scale factor that maps that maximum to the top of the FP8 range. The cast multiplies by the scale, rounds to FP8, and the GEMM runs on the scaled values; the FP32 result is divided by the product of the two operand scales before it leaves the tensor core. The scale is ``delayed'' because it comes from previous iterations rather than from a pass over the current tensor, which would cost an extra read of the whole tensor; the history makes the estimate robust to a single unusually large step. Weights and activations in the forward pass use E4M3, because their precision matters more than their range once scaled; gradients in the backward pass use E5M2, because they span more orders of magnitude within one tensor. Accumulation is in FP32 inside the tensor core. The master copy of the weights, the optimizer state and the residual stream stay in FP32 or BF16; FP8 is a storage and multiplication format for the matrix multiplies, not the precision of the model.

The hardware side is the fourth-generation tensor core's FP8 mode. At the same clock an SM performs twice as many FP8 multiply-adds as FP16 ones, because the multipliers are a quarter the size, so the dense FP8 peak of the H100 SXM5 is 1{,}979 TFLOP/s against 989 for FP16; with 2:4 sparsity the datasheet quotes 3{,}958, which is the figure most often repeated$^\dagger$. Against the A100's 312 TFLOP/s dense BF16, the H100's dense FP8 is $1{,}979/312 = 6.3\times$ per GPU (derived), and the FP8 operands are half the bytes of BF16, so the bandwidth needed per FLOP also halves. Micikevicius, Stosic, Burgess and colleagues at NVIDIA, Arm and Intel described the two formats and the training results in ``FP8 formats for deep learning'' (2022): image classification, machine translation and language models up to 175 billion parameters trained in FP8 matched their BF16 baselines within the run-to-run noise. The paper also proposed the two encodings as a cross-vendor standard, and AMD's MI300X and Intel's Gaudi 3 implement the same pair.

The library did the engineering. Transformer Engine, released with Hopper in 2022, provides drop-in layers (linear, layer norm, attention) for PyTorch and JAX that keep the amax histories, choose the formats, perform the casts and call the FP8 kernels; a user converts a BF16 model by wrapping its layers and sees it run faster. That is the pattern slide 29 generalizes: every hardware feature since Volta arrived inside a library that hides it.

Where it went: FP8 was used first for inference (2023), where the weights are quantized once and the activations scaled dynamically; FP8 pretraining of large models at scale followed, and DeepSeek-V3 (December 2024), trained on H800 GPUs with FP8 matrix multiplies and a fine-grained scaling scheme, made the case publicly. The block-scaled formats followed in the order the deck gives. NVFP4, with 16-value blocks and an FP8 scale per block (slide 18), was used in a submitted training run for the first time in MLPerf Training v5.1 (November 12, 2025), where NVIDIA trained the Llama 3.1 405B benchmark on GB300 with NVFP4 matrix multiplies, which NVIDIA described as a first in the history of MLPerf Training$^\dagger$. In MLPerf Training 6.0 (June 2026) the DeepSeek-V3 run on 8{,}192 GB300 GPUs used MXFP8, the 32-value-block FP8 format (slide 10), for the attention batched matrix multiplies and the final logit GEMM, with the NVFP4 recipe for the rest (slide 20)$^\dagger$. The chronology matters: 4-bit training did not wait for a block-scaled FP8 to be proven; both arrived with Blackwell's hardware block scaling and were adopted together.

\begin{formulas}
FP8 casts and unscaling for one GEMM $C = AB$, with per-tensor scales $s_A, s_B$ chosen from the amax histories:
\[
\tilde A = \mathrm{fp8}(s_A A),\qquad \tilde B = \mathrm{fp8}(s_B B),\qquad
C = \frac{1}{s_A s_B}\,\mathrm{acc}_{32}\bigl(\tilde A \tilde B\bigr),\qquad
s_X = \frac{\mathrm{max\_fp8}}{\max_{k\in\text{history}} \mathrm{amax}_k(X)} .
\]
With $\mathrm{max\_fp8}=448$ for E4M3 and $57{,}344$ for E5M2.
\end{formulas}

\begin{terms}
\term{FP8}{An 8-bit floating-point format; Hopper supports two, E4M3 and E5M2.}
\term{E4M3}{FP8 with a 4-bit exponent and 3-bit mantissa; largest finite value 448, no infinities; used for weights and activations.}
\term{E5M2}{FP8 with a 5-bit exponent and 2-bit mantissa; the range of FP16 (largest finite value 57{,}344); used for gradients.}
\term{Transformer Engine}{NVIDIA's library (2022) of transformer layers that run matrix multiplies in FP8 on Hopper and later GPUs, choosing formats and scales automatically.}
\term{amax}{The maximum absolute value in a tensor, the statistic from which an FP8 scale factor is chosen.}
\term{delayed scaling}{Choosing a tensor's FP8 scale from the amax values of previous iterations rather than from a pass over the current tensor.}
\term{per-tensor scale}{One multiplicative factor applied to a whole tensor before casting to a narrow format and divided out after the multiply.}
\term{loss scaling}{Multiplying the training loss by a constant so that small gradients stay representable in a narrow format; the FP16 predecessor of per-tensor scaling.}
\term{FP32 accumulation}{Summing the products of a matrix multiply in 32-bit floating point even when the operands are 8 or 16 bits, so that rounding errors do not accumulate.}
\term{MXFP8}{The block-scaled FP8 of the OCP microscaling standard: every 32 values share one 8-bit power-of-two scale.}
\end{terms}

\begin{intuition}
Eight bits cannot hold a tensor's values, but they can hold the ratios of each value to the tensor's largest, which is all a matrix multiply needs if the result is accumulated in FP32 and the scale is multiplied back afterwards. The Transformer Engine is the bookkeeping that keeps track of the largest value; the tensor core is the part that runs twice as fast because the multipliers are a quarter the size.
\end{intuition}

% ---------------------------------------------------------------------
\slide{13}{Clusters, distributed shared memory, and the TMA}{Hopper added a level between the block and the grid, and a copy engine that moves tiles while the SM computes.}

\para{What the slide shows} The diagram at the top left shows eight SMs of one GPC as a row of boxes, each with a small ``shared'' memory block at its bottom; four of them, SMs 2 to 5, are highlighted and labeled block 0 to block 3, and double-headed arrows between their shared memories are labeled ``SM-to-SM network: distributed shared memory (DSMEM)''. A bar beneath the SMs represents the L2 cache and the tensor memory accelerator's asynchronous tile copies into shared memory. The panel ``The new hierarchy'' places the cluster between the block and the grid; the right panel, ``Three mechanisms'', explains the TMA, the warpgroup matrix instruction and DSMEM, and the pipeline they make possible.

\para{Details} The CUDA model of slide 9 in the previous deck has threads, blocks and grids: a block's threads run on one SM, share its shared memory and synchronize with a barrier; blocks of a grid are independent and may run anywhere, in any order. The two levels match the hardware of 2007 to 2020. Hopper inserts a level between them. A thread block cluster is a group of blocks that the hardware guarantees to schedule concurrently on the SMs of one GPC; CUDA guarantees clusters of up to 8 blocks as portable, and Hopper allows 16 with an opt-in attribute. Blocks in a cluster can address each other's shared memory, which NVIDIA calls distributed shared memory (DSMEM): a thread in block 2 can load from, store to and perform atomics on an address in block 3's shared memory, over a network inside the GPC that connects the SMs' shared-memory arrays, and the cluster has a barrier that all its threads can wait on. Luo, Fan, Li, Du, Wang and Chu, in ``Benchmarking and dissecting the Nvidia Hopper GPU architecture'' (2024), measure DSMEM access latency between the local shared memory's and the L2's, and bandwidth well above the L2's for the traffic patterns a cluster generates.

What the cluster buys is the removal of a round trip. In a tiled matrix multiply or a stencil, neighboring blocks need overlapping data: the same tile of one operand, or a halo of cells at a tile boundary. Before Hopper each block loaded its copy from L2 (or DRAM), and a reduction across blocks (split-$K$ in a GEMM, a histogram, a sum over the grid) went through global memory with a second kernel or atomics. With a cluster, one block loads a tile and the others read it from DSMEM; partial sums are combined in shared memory within the cluster; the L2 and DRAM traffic falls by the cluster size for the shared operand. The halo exchange of a stencil becomes an on-chip copy.

The second mechanism is the tensor memory accelerator (TMA). Through Ampere, moving a tile from global memory into shared memory took every thread of the block: each computed an address, issued a load (or, from Ampere, a \texttt{cp.async} copy that bypassed the registers), and the block waited at a barrier. The address arithmetic and the issue slots cost instructions that compete with the math, and with Hopper's tensor cores doing 2{,}048 multiply-adds per clock per SM, the instruction budget for feeding them became the limit. The TMA is a unit in the SM that takes a descriptor, the tensor map, describing a tensor in global memory (base address, element type, dimensions, strides) and a tile shape, and copies a tile between global and shared memory on its own. It computes the addresses, clamps or zero-fills at the tensor's boundaries so that edge tiles need no special code, applies a swizzle pattern to the shared-memory layout so that the subsequent reads by the tensor cores have no bank conflicts, and signals completion by incrementing an asynchronous transaction barrier, a new barrier that counts bytes arrived as well as threads. One thread issues the copy; the other threads of the block are free. The TMA also handles the reverse direction (shared to global) and multicast, in which one copy delivers a tile to the shared memories of several blocks in a cluster at once.

The third mechanism is the warpgroup matrix instruction, \texttt{wgmma} in PTX. On Volta through Ampere, a tensor-core operation was issued by one warp on operands in that warp's registers (\texttt{mma.sync}); on Hopper, four consecutive warps (128 threads, a warpgroup) issue one asynchronous instruction whose $A$ operand may come from registers or shared memory and whose $B$ operand comes from shared memory, with the accumulator in registers. Because the instruction is asynchronous and the operands are already in shared memory in the TMA's swizzled layout, the SM can run a software pipeline: the TMA fetches tile $k{+}1$ (and $k{+}2$) while the tensor cores consume tile $k$, with the transaction barriers sequencing producer and consumer warps. NVIDIA's CUTLASS 3 library and its CuTe layout algebra (2023) are built on this pipeline, and FlashAttention-3 (Shah, Bhatia, Dao et al., 2024) applied it to attention, overlapping the GEMMs with the softmax on the multi-function units; both reach roughly 70--75\% of the dense FP16 peak on H100, where the same algorithms on Ampere-style code reached 35--50\%$^\dagger$.

Together the three mechanisms change what a CUDA programmer does. The warps of a block specialize: a producer warp issues TMA copies, consumer warpgroups issue \texttt{wgmma}, and barriers connect them; the data moves in tiles described by shapes rather than by per-thread addresses; and the unit of cooperation is the cluster rather than the block. This is the programming model that Blackwell's tensor memory and \texttt{tcgen05} instructions (slide 18) extend, and it is why the memory-hierarchy table of slide 15 can show per-SM storage flat across three generations: the operands no longer pass through the registers at all.

\begin{terms}
\term{thread block cluster}{A group of up to 8 (16 non-portable) thread blocks guaranteed to run concurrently on the SMs of one GPC and able to access each other's shared memory; new in Hopper.}
\term{distributed shared memory}{DSMEM: the shared memories of the blocks in a cluster, addressable from any block in the cluster over an SM-to-SM network.}
\term{tensor memory accelerator}{TMA: a Hopper SM unit that copies multi-dimensional tiles between global and shared memory from a descriptor, handling addresses, boundaries and swizzling, and signals an asynchronous barrier when done.}
\term{tensor map}{The descriptor a TMA copy uses: base address, element type, dimensions, strides and tile shape of a tensor in global memory.}
\term{swizzle}{A permutation of the addresses within a shared-memory tile chosen so that the access pattern of a matrix instruction hits distinct banks.}
\term{asynchronous transaction barrier}{A Hopper barrier that completes when a count of threads has arrived and a count of bytes has been written, used to sequence TMA copies and the warps that consume them.}
\term{warpgroup}{Four consecutive warps, 128 threads, the unit that issues a wgmma instruction.}
\term{wgmma}{Warpgroup matrix multiply-accumulate: Hopper's asynchronous tensor-core instruction with operands in shared memory.}
\term{split-K}{Dividing a matrix multiply's inner dimension among blocks so that each computes a partial sum, which must then be reduced.}
\term{halo}{The cells at the border of a tile that a stencil computation needs from the neighboring tile.}
\term{CUTLASS}{NVIDIA's open-source template library of matrix-multiply kernels; version 3 (2023) is built on TMA and wgmma.}
\term{FlashAttention-3}{The 2024 attention kernel for Hopper that overlaps tensor-core GEMMs with the softmax using warp specialization and TMA.}
\end{terms}

\begin{intuition}
Hopper separated the two jobs an SM had always done with the same threads: moving tiles and multiplying them. A copy engine moves the tiles by description, a matrix instruction multiplies them from shared memory, and barriers connect the two; the threads become supervisors rather than couriers. The cluster is the same separation applied one level up, so that neighboring blocks pass tiles to each other instead of each fetching its own.
\end{intuition}

% ---------------------------------------------------------------------
\slide{14}{NVLink 4, NVSwitch 3, and the Grace Hopper superchip}{900 GB/s between GPUs, reductions in the switch, and a CPU joined to the GPU at memory speed.}

\para{What the slide shows} Two diagrams on the left. The upper one is a GH200 Grace Hopper superchip: a box for the Grace CPU (72 Neoverse V2 cores, 480 GB of LPDDR5X at up to 512 GB/s) and a box for the H100 GPU (132 SMs, 528 tensor cores, 96 GB of HBM3 or 144 GB of HBM3e), joined by a thick double arrow labeled ``C2C 900 GB/s'', with a caption on the single address space, the external links and the 450 to 1{,}000 W module power. The lower one is a DGX H100 node: eight H100 boxes connected to four third-generation NVSwitch chips, with a caption on the switch (64 ports, 13.6 Tb/s, SHARP reductions) and the NVLink Switch System that could join 256 GPUs. The right panel covers NVLink 4 and the HGX H100 board, what Grace changes, the H200 memory upgrade, and the systems built from Hopper.

\para{Details} NVLink 4, Hopper's GPU-to-GPU link, has 18 links per GPU at 50 GB/s each (25 GB/s per direction), for 900 GB/s of total bandwidth per GPU, against 600 GB/s on A100 and 300 on V100$^\dagger$. The standard way to use it in a node is the HGX H100 baseboard: eight SXM5 modules and four third-generation NVSwitch chips, every GPU connected to every switch, so that any two GPUs communicate at the full 900 GB/s and the eight behave as one memory pool for collective operations. The DGX H100 is NVIDIA's own server built on that board, with two Xeon CPUs, eight ConnectX-7 network adapters at 400 Gb/s (one per GPU), 640 GB of HBM3 across the GPUs and a peak of 32 PFLOP/s of sparse FP8 ($8\times3{,}958$), in about 10 kW$^\dagger$. It is the eight-GPU unit that the cloud providers bought in the hundreds of thousands in 2023 and 2024, and the HGX board (sold to server makers) is how most H100s reached data centers.

The third-generation NVSwitch is more than a crossbar. It has 64 NVLink 4 ports; NVIDIA quotes 13.6 Tb/s of aggregate bandwidth for the chip, and 64 ports at 25 GB/s per direction is 1.6 TB/s each way (derived)$^\dagger$. New in this generation is SHARP (scalable hierarchical aggregation and reduction protocol), brought over from Mellanox InfiniBand switches: the switch can perform the arithmetic of a reduction itself, so that an all-reduce of gradients across eight GPUs sends each GPU's contribution once to the switch, which adds them and returns the sum, instead of the ring or tree of pairwise exchanges that NCCL otherwise runs. This halves the data each GPU moves for an all-reduce and removes the latency of the intermediate hops; NVIDIA's NVLink 6 switch extends the idea to 130 TFLOP/s of in-network arithmetic (slide 25). NVIDIA also defined an NVLink Switch System in which a second tier of external NVSwitch chips joins up to 32 nodes, 256 GPUs, into one NVLink domain with 57.6 TB/s of bisection bandwidth$^\dagger$; it was offered with the DGX H100 SuperPOD, but no public deployment of a 256-GPU NVLink domain is known and NVIDIA never announced one; the cost of the optical links was high, and the 72-GPU rack of the next generation (slide 19) made the same point with copper.

Grace is NVIDIA's first server CPU, announced at GTC 2021 and shipped with Hopper. It has 72 Arm Neoverse V2 cores, up to 480 GB of LPDDR5X memory with error correction at up to 512 GB/s, and is built to pair with GPUs rather than to compete with x86 servers on general workloads$^\dagger$. LPDDR5X is the memory of phones and laptops; soldered to the module rather than socketed, it gives server-class capacity at a fraction of the power of DDR5 DIMMs, and the choice signals what the CPU is for: holding data for the GPU. The connection is NVLink-C2C, a chip-to-chip version of NVLink at 900 GB/s, coherent, so that the GPU and the CPU share one address space and the GPU can read CPU memory at a bandwidth seven times that of PCIe 5 x16 (128 GB/s). A model or a data structure larger than the GPU's HBM can live partly in the CPU's 480 GB and be paged or prefetched at NVLink speed; the ``superchip'', one CPU and one GPU on a module, is the unit of design.

The GH200 Grace Hopper superchip pairs Grace with an H100 that carries 96 GB of HBM3 at 4.0 TB/s (announced 2022, in production from mid-2023) or, in the 2023 revision, 144 GB of HBM3e at 4.9 TB/s$^\dagger$. The module is rated from 450 to 1{,}000 W depending on configuration$^\dagger$. Two GH200s can be joined with NVLink into a 288 GB HBM pair, and AWS announced a 32-superchip GH200 NVL32 configuration in November 2023 (Project Ceiba)$^\dagger$. The superchip is the ancestor of the GB200 (one Grace, two Blackwell) and the Vera Rubin superchip (slide 24).

The H200, announced on November 13, 2023 and available from the second quarter of 2024, is the GH100 die with 141 GB of HBM3e at 4.8 TB/s, 1.76 times the H100's capacity and 1.4 times its bandwidth, with the same compute$^\dagger$. That a memory-only upgrade was worth a product name is the lesson of slide 15: large-language-model inference at a given latency is bound by how fast the weights and the key--value cache can be read, so bandwidth and capacity, not FLOP/s, set the tokens per second. NVIDIA also sold an H100 NVL, two PCIe cards bridged by NVLink with 94 GB each (188 GB), for the same market$^\dagger$.

The systems: NVIDIA's own Eos, 576 DGX H100 nodes with 4{,}608 H100s, reached 121 PFLOP/s on HPL and ninth place in November 2023$^\dagger$; Microsoft Azure's Eagle, an H100 cloud system, 561 PFLOP/s and seventh place on the June 2026 list$^\dagger$; CSCS's Alps in Switzerland, built from GH200 superchips, 435 PFLOP/s$^\dagger$; the University of Bristol's Isambard-AI, 5{,}448 GH200$^\dagger$; and J\"ulich's JUPITER, about 24{,}000 GH200 in an Eviden BullSequana system, which reached 1.000 EFLOP/s on HPL in November 2025, the first European system to do so, and held fifth place in June 2026$^\dagger$. These are the Hopper machines on the TOP500; the far larger Hopper installations, the hyperscalers' training clusters of tens of thousands of H100s, were never submitted (slide 27).

\begin{terms}
\term{HGX H100}{NVIDIA's eight-GPU baseboard with four NVSwitch chips, the building block of H100 servers from NVIDIA and from server makers.}
\term{DGX H100}{NVIDIA's own server on the HGX H100 board: eight H100s, two Xeon CPUs, eight 400 Gb/s adapters, 32 PFLOP/s sparse FP8.}
\term{NVSwitch 3}{The third-generation NVLink switch chip: 64 NVLink 4 ports, with SHARP in-network reductions.}
\term{SHARP}{Scalable hierarchical aggregation and reduction protocol: performing the arithmetic of a reduction inside a network switch, so that each endpoint sends and receives once; from Mellanox InfiniBand, in NVSwitch since Hopper.}
\term{all-reduce}{The collective operation that sums (or otherwise combines) a tensor across all GPUs and leaves the result on every one; the dominant communication in data-parallel training.}
\term{NVLink Switch System}{A second tier of NVSwitch chips joining up to 256 Hopper GPUs into one NVLink domain; offered, with no public deployment known.}
\term{Grace}{NVIDIA's first server CPU (2022--23): 72 Arm Neoverse V2 cores, up to 480 GB of LPDDR5X at up to 512 GB/s, joined to a GPU by NVLink-C2C.}
\term{Neoverse V2}{Arm's licensable high-performance server core used in Grace; Vera replaces it with NVIDIA's own Olympus core.}
\term{LPDDR5X}{Low-power DRAM of the kind used in phones and laptops, soldered to the module; Grace uses it for capacity at low power.}
\term{NVLink-C2C}{The chip-to-chip form of NVLink joining a CPU and a GPU coherently; 900 GB/s on Grace Hopper, 1.8 TB/s on Vera Rubin.}
\term{superchip}{NVIDIA's name for a module carrying one CPU and one or two GPUs joined by NVLink-C2C: GH200, GB200, Vera Rubin.}
\term{GH200}{The Grace Hopper superchip: Grace plus an H100 with 96 GB HBM3 or 144 GB HBM3e.}
\term{H200}{The H100 die with 141 GB of HBM3e at 4.8 TB/s (2024); a memory upgrade sold as a product.}
\term{coherent memory}{Memory that two processors share with hardware keeping their caches consistent, so that either can read the other's data without explicit copies.}
\end{terms}

\begin{intuition}
With Hopper the node stopped being a CPU with GPUs attached and became a GPU fabric with a CPU attached: eight GPUs talk to each other at 900 GB/s through switches that can do arithmetic, and the CPU that joins them is chosen for how fast it can hand memory to a GPU. Everything NVIDIA sells after this is a bigger version of the same picture.
\end{intuition}

% ---------------------------------------------------------------------
\slide{15}{The memory hierarchy, and the ridge that keeps moving}{Register files and L1 barely changed per SM; L2 grew 8$\times$, HBM 21$\times$; the tensor peak grew 280$\times$.}

\para{What the slide shows} A table across the left half compares six GPUs, V100, A100, H100, GB200, GB300 and Rubin, on eleven rows: SMs enabled; registers per SM; L1 and shared memory per SM; tensor memory per SM; L2; HBM capacity; HBM generation; bandwidth; the fastest dense peak and the precision it is at; and the ridge point, peak divided by bandwidth, in FLOP per byte. Entries NVIDIA has not published are marked ``n.p.'' and a footnote names the microbenchmark source for the Blackwell SM count and tensor memory. The right panel draws three conclusions: per-SM storage is flat, the ridge point rose from 139 to about 1{,}800 FLOP per byte (peaking at GB300), and HBM bandwidth per stack rose ten-fold.

\para{Details} The table is the deck's argument in numbers. Read it row by row.

SMs enabled: 80, 108, 132, 148, 160, not published. The count grew 2$\times$ from V100 to GB300, and the Blackwell figures are not NVIDIA's: the 148 SMs of the B200 come from Jarmusch and Chandrasekaran's microbenchmarks (2025), which enumerate the SMs the runtime exposes, and the 160 of the B300 from NVIDIA's Blackwell Ultra blog, which gives 8 GPCs of 20 SMs across the two dies$^\dagger$. NVIDIA has published no SM count for Rubin as of November 2026.

Registers per SM: 256 KB on all three generations for which it is published, unchanged since Pascal (2016). L1 and shared memory per SM: 128 KB (V100), 192 KB (A100), 256 KB (H100), then unpublished. Tensor memory: none until Blackwell, then 256 KB per SM, a new storage array dedicated to tensor-core accumulators (slide 18), again from microbenchmarks and NVIDIA's PTX documentation rather than a whitepaper. So the storage inside an SM grew from 384 KB to 512 KB to perhaps 768 KB (if Blackwell's L1 is unchanged) over nine years, while the SM's dense FP16 tensor rate went from 512 to 1{,}024 to 2{,}048 FMA per clock and beyond. The operands of a matrix multiply are not being staged through more registers; they are being streamed through new paths: \texttt{cp.async} on Ampere, the TMA and shared-memory operands on Hopper, tensor memory on Blackwell (slides 8, 13, 18).

L2: 6, 40, 50 MB, then unpublished. The jump from Volta to Ampere, 6.7$\times$, is the one structural change in the on-chip hierarchy in the period; Hopper added 25\%, and NVIDIA stopped reporting the figure with Blackwell, so that the microbenchmark papers are again the source.

HBM: 16 or 32 GB of HBM2 at 0.9 TB/s (V100), 40 GB of HBM2 at 1.6 TB/s or 80 GB of HBM2e at 2.0 (A100), 80 GB of HBM3 at 3.35 (H100), 186 GB of HBM3e at 8.0 (GB200; 180 GB and 7.7 TB/s on the HGX B200), 288 GB of HBM3e at 8.0 (GB300, twelve-high stacks), 288 GB of HBM4 at 19.2 (Rubin, per the September 2026 specification page; 22 TB/s had been announced in January)$^\dagger$. Capacity grew 18$\times$ and bandwidth 21$\times$ over the nine years, through four memory generations and two changes in stack height. Per stack, the figures are 0.23 TB/s for HBM2 on V100 (900 GB/s over four stacks), 0.67 for HBM3 on H100 (3.35 TB/s over five), 1.0 for HBM3e on B300 (8 TB/s over eight) and 2.4 for HBM4 on Rubin (19.2 TB/s over eight), all derived from the product figures; HBM4 doubles the interface width per stack to 2{,}048 bits, which is where most of its step comes from.

The fastest dense peak: 125 TFLOP/s of FP16 (V100), 312 of BF16 (A100), 1.98 PFLOP/s of FP8 (H100), 10 PFLOP/s of FP4 (GB200), 15 (GB300), 35 (Rubin, the training figure)$^\dagger$. This is the 280$\times$ of the deck's second slide, and the precision row beneath it is the reason: each step down the ladder of slide 10 doubled the rate for the same silicon.

The ridge point is the peak divided by the bandwidth, the arithmetic intensity (FLOP per byte of DRAM traffic) at which the roofline of slide 15 in the previous deck turns from the bandwidth slope to the compute ceiling. In FLOP per byte: $125{,}000/900 = 139$ (V100), $312{,}000/1{,}555 = 200$ (A100), $1{,}979{,}000/3{,}350 = 591$ (H100), $10{,}000{,}000/8{,}000 = 1{,}250$ (GB200), $15{,}000{,}000/8{,}000 = 1{,}875$ (GB300), $35{,}000{,}000/19{,}200 = 1{,}823$ (Rubin), all derived. The ridge therefore peaks at the GB300 and falls by about 3\% on Rubin: between the two generations HBM bandwidth rose 2.4$\times$ (8 to 19.2 TB/s) while the dense FP4 rate rose 2.33$\times$ (15 to 35 PFLOP/s), the first generation since Volta in which memory bandwidth grew faster than the fastest matrix rate. A kernel must perform 1{,}800 floating-point operations for every byte it reads from HBM to reach Rubin's FP4 peak; a dense matrix multiply of $n\times n$ matrices performs $2n^3$ operations on $3n^2$ values, so with 4-bit values ($n^2\times 1.5$ bytes) its intensity is roughly $1.3n$ FLOP per byte and $n$ must exceed about 1{,}400, with the tiles reused from shared memory to make the DRAM traffic that small. Attention over a long context, layer normalization, the elementwise operations between layers, and above all the generation of one token at a time at small batch, where every weight is read once per token, have intensities of a few FLOP per byte and sit on the bandwidth slope of the roofline, a factor of hundreds below the peak. This is why the H200 was a product, why Blackwell Ultra spent its refresh on 288 GB rather than more FLOP/s, why HBM4's 19.2 TB/s is the number NVIDIA's own spec page leads with, and why an inference GPU is judged in tokens per second per watt or per dollar rather than in FLOP/s.

Two consequences for the student. First, the figures in this table are the ones that determine a kernel's performance class, and all but the peaks are the ones a datasheet is least likely to give: when NVIDIA stopped publishing L2 and register sizes with Blackwell, the microbenchmark papers (Jia et al.\ 2018 for Volta, Abdelkhalik et al.\ 2022 for Ampere, Luo et al.\ 2024 for Hopper, Jarmusch et al.\ 2025 for Blackwell) became the primary sources, and the datasheet still omits latencies, achieved bandwidth and clocks under power limits. Second, the ridge point only rises: a kernel that was compute-bound on a V100 at 200 FLOP per byte is bandwidth-bound on everything since, and the work of a performance engineer on each new generation is to find more reuse (bigger tiles, fused kernels, lower-precision operands) to climb the slope again.

\begin{formulas}
Ridge point and the intensity of a dense matrix multiply:
\[
I^\ast = \frac{P_{\text{peak}}}{B},\qquad
I_{\text{GEMM}} \approx \frac{2n^3}{3n^2 b} = \frac{2n}{3b}\ \text{FLOP per byte},
\]
with $b$ bytes per element ($b=0.5$ for FP4, 1 for FP8, 2 for FP16), assuming each matrix is read from DRAM once. For Rubin, $I^\ast = 35\times10^{15}/19.2\times10^{12} \approx 1{,}823$; setting $I_{\text{GEMM}} = I^\ast$ with $b=0.5$ gives $n \approx 1{,}370$.
\end{formulas}

\begin{terms}
\term{tensor memory}{TMEM: a 256 KB per-SM storage array on Blackwell that holds tensor-core accumulators outside the register file.}
\term{HBM3e}{Extended HBM3 with faster pins; 8 TB/s over eight stacks on Blackwell; twelve-high stacks give 288 GB on B300.}
\term{HBM4}{The fourth HBM generation, with a 2,048-bit interface per stack; 288 GB at 19.2 TB/s over eight stacks on Rubin.}
\term{stack height}{The number of DRAM dies in an HBM stack (eight-high, twelve-high); capacity per stack scales with it, bandwidth does not.}
\term{ridge point}{Peak FLOP/s divided by bandwidth: the arithmetic intensity at which a roofline turns from bandwidth-bound to compute-bound.}
\term{arithmetic intensity}{Floating-point operations performed per byte moved from DRAM; the x-axis of the roofline.}
\term{roofline}{The performance bound $\min(P_{\text{peak}},\, B\cdot I)$ as a function of intensity $I$.}
\term{bandwidth-bound}{A kernel whose intensity is below the ridge point, so that its speed is set by memory bandwidth; token generation at small batch is the standard example.}
\term{microbenchmark}{A small program timing one hardware mechanism (a cache level, an instruction, a link) to measure what the datasheet omits.}
\term{n.p.}{Not published: a figure NVIDIA has not disclosed; the notes use it rather than an estimate.}
\end{terms}

\begin{intuition}
Every generation moves the ridge to the right, so yesterday's compute-bound kernel is today's bandwidth-bound one. The register file did not grow because the operands stopped visiting it; the memory did grow, and it still cannot keep up with the matrix units, which is why the memory line on a spec sheet is now the one to read first.
\end{intuition}

% =====================================================================
\setstage{mlp}
\section{Ada Lovelace (slide 16)}

\slide{16}{Ada Lovelace (2022): the graphics die, and the L40S}{Hopper's twin on the same process spends its area on ray tracing, a 96 MB L2 and video.}

\para{What the slide shows} A table on the left sets the two compute dies against the two graphics dies of the same generations: GA100 (A100) against GA102 (GeForce RTX 3090), and GH100 (H100 SXM) against AD102 (GeForce RTX 4090 and L40S), on fifteen rows: process, transistors, die area, SMs on the die and in the product, FP32 lanes, the FP64 rate relative to FP32, dense FP16 and FP8 tensor peaks, memory type and capacity, bandwidth, L2, NVLink, RT cores and power. A footnote gives the RTX 4090's and L40S's FP32 peaks and clocks, notes that the 4090's FP16 and FP8 tensor rates are the FP16-accumulate figures (half with FP32 accumulation), and marks the L40S SM count and L2 as derived. The right panel describes the AD102 die, explains the 96 MB L2, names DLSS 3, and describes the L40S as the data-center product of the graphics line.

\para{Details} From Ampere on, each NVIDIA generation has two large dies built for different markets (slide 9). The compute die (GA100, GH100, Blackwell) has full-rate FP64, HBM, NVLink and no display or ray-tracing hardware; the graphics die (GA102, AD102, GB202) has ray-tracing cores, display and video engines, GDDR memory, FP64 at 1/64 of the FP32 rate, and since Ada no NVLink. Ada Lovelace is the graphics generation that shipped between Hopper and Blackwell, announced on September 20, 2022 with the GeForce RTX 4090 available in October$^\dagger$. It matters to a lecture on data-center GPUs for two reasons: its largest die, AD102, is Hopper's twin on the same TSMC 4N process, so the two show what a GPU designer does with the same transistor budget for two different customers; and its professional product, the L40S, became a widely deployed inference and rendering card.

AD102 has 76.3 billion transistors on 608.5 mm$^2$, against GH100's 80 billion on 814 mm$^2$: nearly the same transistor count on three-quarters of the area, because the graphics die has no HBM interfaces (which are large) and its SMs carry no FP64 datapath$^\dagger$. The full die has 12 GPCs of 12 SMs, 144 SMs, each with 128 FP32 lanes (18{,}432 in all), four fourth-generation tensor cores and one third-generation RT core; the RTX 4090 enables 128 SMs (16{,}384 lanes, 512 tensor cores, 128 RT cores) at a boost clock of 2.52 GHz, which gives $16{,}384\times2\times2.52 = 82.6$ TFLOP/s of FP32 (derived; the whitepaper's figure)$^\dagger$. The clock is the first thing to notice against Hopper: the H100's boost is about 1.98 GHz for 700 W; the 4090 runs 27\% faster on its lanes at 450 W, because a graphics die on GDDR with no HBM and no NVLink spends its power budget on the SMs.

The SM is the Ampere GA102 SM carried forward: two 16-wide FP32 datapaths per partition, one of which doubles as INT32, so 128 FP32 lanes per SM; the fourth-generation tensor core adds FP8 as on Hopper, and the whitepaper quotes the RTX 4090 at 330 TFLOP/s dense FP16 and 661 dense FP8 with FP16 accumulation, doubled with sparsity$^\dagger$. The GeForce products halve the tensor rate when the accumulator is FP32 (165 and 330 TFLOP/s dense), a deliberate segmentation between the gaming and professional cards that the L40S does not have; the table's 330 is the FP16-accumulate figure, and a reader comparing it with the H100's 989 (FP32 accumulate, no segmentation) should keep the distinction in mind.

The 96 MB L2 is the architectural point. GA102 had 6 MB; AD102 has 96 MB on the full die, 72 MB on the RTX 4090, sixteen times its predecessor$^\dagger$. The reason is bandwidth. The 4090's memory is 24 GB of GDDR6X on a 384-bit bus at 21 Gb/s per pin, 1{,}008 GB/s, less than a third of the H100's 3.35 TB/s and about the same as the V100's HBM2 of 2017$^\dagger$; a die with four times the V100's FP32 peak on the same bandwidth would starve. A large L2 catches the reuse that a game's frame has (textures, geometry, the framebuffer tiles being worked on) and that an inference workload has (the weights of a layer being applied to a batch), so that the fraction of accesses that go to DRAM falls. Where the compute die buys its bandwidth with HBM stacks on an interposer, the graphics die buys it with SRAM on the die, as AMD had done with the ``Infinity Cache'' on its RDNA 2 GPUs in 2020. The cost is area: 96 MB of SRAM is a substantial share of the 608 mm$^2$.

Ada's graphics additions are three. Shader execution reordering (SER) lets a ray-tracing shader program ask the hardware to regroup its threads by the material or shader they will execute next, so that the divergence slide 11 of the previous deck describes, which ray tracing produces in abundance because adjacent rays hit different surfaces, is reduced before the shading code runs; NVIDIA quotes up to 2$\times$ on ray-traced shading$^\dagger$. The third-generation RT core adds an opacity micromap engine (alpha-tested geometry such as foliage resolved inside the RT core without calling a shader) and a displaced micro-mesh engine (fine geometric detail generated on the fly from a coarse mesh). DLSS 3 adds frame generation: an optical flow accelerator, a fixed-function unit that computes the motion field between two rendered frames, feeds a network on the tensor cores that synthesizes an intermediate frame, so that the display shows two frames for each one the shaders render$^\dagger$. Ada also carries two eighth-generation NVENC video encoders with AV1 support, which matters for the L40S's streaming and transcoding market.

The L40S is the data-center product of the AD102 die, introduced in 2023 after the L40 of 2022: 142 SMs (18{,}176 lanes), 48 GB of GDDR6 with error correction at 864 GB/s, 96 MB of L2, 91.6 TFLOP/s FP32, 362 TFLOP/s dense FP16 and 733 dense FP8 on the tensor cores (733 and 1{,}466 with sparsity), 350 W, a dual-slot passively cooled PCIe card with no NVLink$^\dagger$. The datasheet states the lane count, memory, FP32 and tensor peaks and power; the 142-SM figure is derived from 18{,}176 lanes at 128 per SM, and the 96 MB L2 comes from secondary sources, since the datasheet does not list either. It sits in ordinary servers that cannot host an HGX board, and it serves the workloads that do not need HBM bandwidth or GPU-to-GPU links: inference of models that fit in 48 GB, fine-tuning, rendering, virtual workstations, video. Against the H100 PCIe (350 W, 80 GB HBM2e at 2.0 TB/s, 1{,}513 dense FP8) it has 43\% of the bandwidth, half the FP8 and 60\% of the memory, and sells for a fraction of the price; for small-batch inference of a 7B- or 13B-parameter model it is bandwidth-bound on both, and the ratio of their bandwidths is the ratio of their speeds. The lesson for reading a datasheet: the two dies divide the market by memory bandwidth and by interconnect, not by FLOP/s.

What the table does not show, and the next sections do, is that the division moved. Blackwell's graphics die (GB202, GeForce RTX 5090, 2025) kept the pattern; but NVIDIA's data-center roadmap added a third kind of part, the Rubin CPX (slide 26), a Rubin die with GDDR7 rather than HBM, aimed at the prefill phase of inference where compute rather than bandwidth is the limit. The graphics die's memory system, in other words, found a data-center use of its own.

\begin{terms}
\term{Ada Lovelace}{NVIDIA's graphics GPU architecture of 2022, named after the mathematician; the dies are AD102 to AD107, the products GeForce RTX 40 and the L40 and L40S.}
\term{AD102}{The largest Ada die: 76.3 billion transistors, 608.5 square millimeters, TSMC 4N, 144 SMs, 96 MB L2.}
\term{graphics die}{The die of a generation built for gaming, rendering and inference: RT cores, display and video engines, GDDR memory, FP64 at 1/64 rate; GA102, AD102, GB202.}
\term{compute die}{The die built for training and HPC: full-rate FP64, HBM, NVLink, no RT or display hardware; GA100, GH100, the Blackwell data-center dies.}
\term{GDDR6X}{Graphics DRAM with PAM4 signaling, 21 Gb/s per pin on the RTX 4090; 1 TB/s over a 384-bit bus.}
\term{RT core}{The fixed-function ray-tracing unit in each SM of a graphics die, which traverses bounding-volume hierarchies and tests ray-triangle intersections; third generation on Ada.}
\term{shader execution reordering}{SER: an Ada mechanism that regroups ray-tracing threads by the shader they will run next, reducing divergence.}
\term{opacity micromap}{An Ada RT-core feature that resolves alpha-tested (partially transparent) triangles inside the RT core without a shader call.}
\term{displaced micro-mesh}{An Ada RT-core feature that generates fine geometric detail from a coarse mesh and a displacement map during traversal.}
\term{DLSS 3}{Deep learning super sampling, version 3: Ada's frame generation, which synthesizes an intermediate frame from two rendered frames and their optical flow on the tensor cores.}
\term{optical flow accelerator}{A fixed-function unit on Ada that computes the motion field between two frames, used by DLSS 3.}
\term{NVENC}{NVIDIA's hardware video encoder; Ada's eighth generation adds AV1 and has two per die.}
\term{L40S}{The AD102 data-center card (2023): 142 SMs, 48 GB GDDR6 at 864 GB/s, 733 TFLOP/s dense FP8, 350 W, PCIe, no NVLink.}
\term{FP16 accumulate}{A tensor-core mode that sums products in 16-bit precision; GeForce cards run FP32-accumulate tensor math at half the FP16-accumulate rate.}
\end{terms}

\begin{intuition}
Give two designers the same transistors and ask one for a gaming card and the other for a training GPU: the first spends them on ray tracing, a huge cache and clock speed, the second on FP64, HBM and links. The L40S shows that the first designer's part is a fine inference card too, as long as the model fits in 48 GB and nobody needs to talk to a second GPU.
\end{intuition}

% =====================================================================
\setstage{posc}
\section{Blackwell and Blackwell Ultra (slides 17--22)}

\slide{17}{Blackwell B200 (2024): two reticle-limited dies as one GPU}{March 2024, shipping from Q4 2024: 208B transistors on TSMC 4NP, joined by a 10 TB/s die-to-die link.}

\para{What the slide shows} A schematic of the Blackwell package fills the upper left: two large dies side by side, each drawn as a grid of forty small squares (the SMs, symbolically) above a bar labeled ``L2 (size not published)'', with a thin vertical band between them labeled ``NV-HBI 10 TB/s'' and four HBM3e stacks drawn along each outer edge, eight in all. Each die is captioned ``104B transistors, 80 SMs''. A caption under the drawing gives the counts the slide can stand behind: 160 SMs of 128 lanes across the two dies, 148 enabled on the B200 (a number from microbenchmarks, not from NVIDIA), one GPU as far as software is concerned, and 186 GB (GB200) or 180 GB (HGX B200) of HBM3e at 8 or 7.7 TB/s. Two stat tiles give the headline figures, 208 billion transistors on TSMC 4NP and 1.8 TB/s of NVLink 5 over 18 links. The right panel, ``What is new'', lists the two dies and their bond, the fifth-generation tensor cores, the links, the fixed-function engines, and the product variants together with what changed between the March 2024 announcement and the shipping parts.

\para{Details} Blackwell was announced at GTC on March 18, 2024, named for the statistician David Blackwell, and is the first NVIDIA compute GPU built from two dies$^\dagger$. The reason is the reticle: a lithography stepper exposes a field of about 858 mm$^2$ at most, and Volta (815 mm$^2$), Ampere (826) and Hopper (814) had already pushed single dies to that limit (slides 11 and 15 of the previous lecture; slide 11 here). Blackwell keeps the process family rather than shrinking it, moving from TSMC's 4N to 4NP, an improved variant of the same 4 nm-class node, so more transistors meant more silicon, and more silicon meant a second die. NVIDIA's figure is 208 billion transistors for the package, 104 billion per die$^\dagger$. The two dies are joined by NV-HBI (NVIDIA High-Bandwidth Interface), a die-to-die link that NVIDIA rates at 10 TB/s$^\dagger$; it carries the L2-to-L2 traffic and lets the two halves present one address space, one L2 domain and one set of SMs to CUDA. A kernel does not know which die it runs on. This is a different choice from AMD's MI250X (2021), whose two dies appeared to software as two GPUs, and closer to AMD's MI300X (2023), whose eight compute dies form one GPU over an interposer; the next lecture compares them.

NVIDIA did not publish the SM count, the L2 size, the register-file size or the shared-memory size for Blackwell, a departure from the whitepapers of Ampere and Hopper. What is known comes from NVIDIA's own ``Inside NVIDIA Blackwell Ultra'' blog (2025), which gives the full configuration as 160 SMs in 8 GPCs of 20 SMs, 20{,}480 lanes$^\dagger$, and from microbenchmarking: Jarmusch and Chandrasekaran (2025) report 148 SMs enabled on the B200, which at 128 FP32 lanes per SM gives 18{,}944 lanes, and a tensor memory of 256 KB per SM (slide 18). These are the numbers the slide quotes, with their provenance stated, because a lecture should not present a reverse-engineered count as a vendor specification. The L2 is drawn without a size for the same reason.

The links are published. NVLink 5 raises the per-link rate to 100 GB/s (bidirectional) and keeps 18 links per GPU, so 1.8 TB/s per GPU against Hopper's 900 GB/s$^\dagger$. The NVLink Switch chip (the fourth generation of NVSwitch) has 72 ports and 7.2 TB/s of aggregate bandwidth, 50 billion transistors on the same 4NP process$^\dagger$; two of them per switch tray and nine trays per rack give the NVL72 of slide 19. The host interface on the shipping HGX B200 is PCIe 5.0, although the March 2024 announcement said PCIe 6$^\dagger$. Three fixed-function additions are worth naming: a decompression engine rated at 800 GB/s for the compressed columnar formats used by analytics engines (LZ4, Snappy, Deflate), so a query engine can read compressed data without a CPU pass; a RAS (reliability, availability, serviceability) engine that monitors thousands of on-chip sensors to predict faults before they take a node down, a fleet-management concern once a training job spans tens of thousands of GPUs; and confidential computing with encrypted NVLink traffic, so that a model and its inputs can run in a trusted execution environment at close to full speed$^\dagger$.

The variants matter for reading any Blackwell number. B100 was announced as a 700 W part that drops into Hopper-class air-cooled servers; NVIDIA never announced a cancellation, but press reports describe it as de-emphasized in favor of the B200 and GB200 within the year and little deployed$^\dagger$. B200 is the 1{,}000 W part on the eight-GPU HGX B200 board. GB200 is the GPU as it appears on the Grace Blackwell superchip, run at 1{,}200 W with liquid cooling, and it is the GB200 whose numbers NVIDIA quotes for the NVL72 rack. The three differ in clock and therefore in peak: HGX B200 is rated at 37 TFLOP/s FP64, 75 FP32, 2.25 PFLOP/s dense FP16, 4.5 dense FP8 and 9 dense FP4 (18 with sparsity), with 180 GB of HBM3e at 7.7 TB/s; GB200 at 40 TFLOP/s FP64, 80 FP32, 2.5 PFLOP/s dense FP16, 5 dense FP8 and 10 dense FP4 (20 sparse), with 186 GB at 8 TB/s$^\dagger$. The numbers announced in March 2024 were higher: a customer deck from that month lists FP64 at 45 TFLOP/s for GB200, 40 for B200 and 30 for B100, 192 GB of HBM3e on every variant, and PCIe Gen 6$^\dagger$. The shipping figures of 180--186 GB, 37--40 TFLOP/s and PCIe 5 are what the datasheets and product pages say in 2025--26. The 192 GB figure is the raw capacity of eight 24 GB stacks; the usable capacity is lower because some of it is reserved, and NVIDIA's own pages now quote 186 GB for GB200 and 180 GB for HGX B200$^\dagger$. The slide's point is not that the announcement was wrong but that a lecture must say which number it is using and when it was published; slide 3 made the same point for Hopper (whose numbers rose) and slide 25 makes it for Rubin (whose numbers fell).

Two more dates fix the timeline. In August 2024 NVIDIA confirmed a mask change to Blackwell ``to improve production yields'', after a report that a design flaw had delayed the ramp; the company said the change had no functional effect and that production would ramp in the fourth quarter of calendar 2024$^\dagger$. Microsoft Azure announced the first GB200 systems (servers, not stated as NVL72) on October 8, 2024, Dell shipped the first GB200 NVL72 racks to CoreWeave on November 18, 2024, and volume ran through 2025$^\dagger$. Within a generation, then, Blackwell went from announcement to volume in about a year, which is the cadence slide 3 describes.

\begin{terms}
\term{reticle limit}{The largest die area a lithography stepper can expose in one shot, about 858 mm$^2$ for current scanners; Volta, Ampere and Hopper dies were within a few percent of it, so Blackwell went to two dies.}
\term{NV-HBI}{NVIDIA High-Bandwidth Interface, the die-to-die connection between Blackwell's two dies, rated at 10 TB/s; it makes the two dies one GPU to software.}
\term{TSMC 4NP}{The improved variant of TSMC's 4 nm-class process used for Blackwell and Blackwell Ultra; the same node family as Hopper's 4N, so the transistor gain came from area, not density.}
\term{HGX B200}{NVIDIA's eight-GPU baseboard with B200 GPUs at 1{,}000 W, air- or liquid-cooled, with NVLink 5 and four NVLink Switch chips; the successor of HGX H100.}
\term{GB200}{The Grace Blackwell superchip: one Grace CPU and two Blackwell GPUs on one board, the GPUs run at 1{,}200 W; also the name of the GPU in that configuration, whose peaks are slightly higher than B200's.}
\term{B100}{The 700 W Blackwell variant announced for Hopper-class servers; de-emphasized after launch and not shipped in volume.}
\term{NVLink 5}{The fifth generation of NVIDIA's GPU-to-GPU link: 100 GB/s per link, 18 links per GPU, 1.8 TB/s per GPU.}
\term{NVLink Switch}{The fourth-generation NVSwitch chip: 72 NVLink 5 ports, 7.2 TB/s, 50 billion transistors; two per switch tray in the NVL72.}
\term{decompression engine}{A fixed-function unit on Blackwell that decompresses LZ4, Snappy and Deflate data at up to 800 GB/s for database and analytics workloads.}
\term{RAS engine}{A reliability, availability and serviceability unit that monitors on-chip sensors and predicts failures so that a GPU can be taken out of a cluster before it fails a job.}
\term{confidential computing}{Running a workload in a hardware-isolated environment whose memory and links are encrypted, so that neither the host operating system nor the cloud operator can read the model or the data; first on Hopper, extended on Blackwell.}
\term{mask change}{A revision of the photomasks used to pattern a chip, made here in August 2024 to improve Blackwell yields without changing its function.}
\end{terms}

\begin{intuition}
Blackwell is Hopper's engineering taken past the edge of the reticle: when one die cannot get bigger, glue two together fast enough that software cannot tell. Everything that is not published about its insides is a reminder that the vendor now sells the package and the rack, not the SM.
\end{intuition}

% ---------------------------------------------------------------------

\slide{18}{Fifth-generation tensor cores: FP4 and tensor memory}{Four-bit values with one scale per block of 16 or 32; accumulators move out of the register file.}

\para{What the slide shows} The upper left panel draws the two 4-bit block formats. The NVFP4 row shows sixteen small squares (the E2M1 values) followed by a box labeled E4M3 (the block scale) and a longer box labeled FP32 (per tensor). The MXFP4 row shows thirty-two smaller squares followed by a box labeled E8M0, the power-of-two scale. A caption gives the cost per value, 4.5 bits for NVFP4 and 4.25 for MXFP4, and lists the sixteen E2M1 codes: 0, 0.5, 1, 1.5, 2, 3, 4, 6 and their negatives. Below is a bar chart of the GB200's tensor peaks for one GPU: 2.5 PFLOP/s FP16/BF16, 5 FP8/FP6, 10 FP4 dense, 20 FP4 with sparsity. The right panel, ``What changed inside the SM'', explains micro-tensor scaling, why a block of 16 values is the right size, the tensor memory, and where FP64 stayed.

\para{Details} Slide 10 ended the precision ladder with the observation that below 8 bits a number cannot stand alone. E2M1 has one sign bit, two exponent bits and one mantissa bit: with the exponent bias of 1 and a subnormal at the bottom, its positive values are 0, 0.5, 1, 1.5, 2, 3, 4 and 6, eight magnitudes in all. No tensor of real weights or activations fits that range on its own; the values in one layer of a transformer span several orders of magnitude, and the outliers that FP8 training already had to manage (slide 12) would be clipped or flushed to zero. The answer, which both the Open Compute Project's Microscaling (MX) specification of September 2023 and NVIDIA's NVFP4 of June 2025 adopt, is a shared scale for a small block of values: each 4-bit value is multiplied by its block's scale, so the block as a whole has the range of the scale and the values inside it need only resolve relative magnitude within the block$^\dagger$.

The two formats differ in the block size and in the scale's type. MXFP4, defined in the OCP MX specification and in the paper by Rouhani et al.\ (2023), uses blocks of 32 values and a scale in E8M0, an 8-bit exponent with no mantissa, which is to say a power of two; the cost is 32 $\times$ 4 bits plus 8 bits of scale, 136 bits for 32 values or 4.25 bits per value. NVFP4 uses blocks of 16 and a scale in E4M3, the FP8 encoding with three mantissa bits, plus a second, FP32 scale for the whole tensor; the cost is 16 $\times$ 4 plus 8 bits, 72 bits for 16 values or 4.5 bits per value, with the FP32 tensor scale amortized to nothing$^\dagger$. The smaller block follows the local magnitude of the data more closely, and the E4M3 scale can land within a factor of about $1 + 2^{-3}$ of the ideal scale where a power-of-two scale can be off by up to a factor of $\sqrt{2}$; NVIDIA's blog puts the effect as a halving of the quantization error of the block scaling itself. The tensor-level FP32 scale keeps the E4M3 block scales in their own range. NVIDIA's claim for inference is accuracy within 1\% of FP8 on its benchmark suite when models are quantized to NVFP4 with its Model Optimizer tooling$^\dagger$. For training, Abecassis et al.\ (NVIDIA, 2025) report pretraining a 12-billion-parameter model on 10 trillion tokens with NVFP4 matrix multiplies and a loss curve matching an FP8 run, using a recipe that keeps some sensitive layers in higher precision, applies random Hadamard transforms to spread outliers, and uses stochastic rounding for gradients. This is the first published FP4 pretraining at that scale and is the reason the deck treats FP4 as a training format and not only an inference one.

The hardware side is the second-generation Transformer Engine and the fifth-generation tensor core. In Hopper, the Transformer Engine was a software library that chose per-tensor scales and cast tensors to FP8 before the matrix multiply (slide 12). In Blackwell, micro-tensor scaling is done in the tensor core: the block scales travel with the operands, the multiplier expands each 4-bit or 6-bit value by its block scale as it is consumed, and products are accumulated in FP32$^\dagger$. FP8 moves to the same block-scaled form (MXFP8, with blocks of 32 and E8M0 scales), which NVIDIA used for the attention and logit matrix multiplies of its DeepSeek-V3 submission in MLPerf Training 6.0 (June 2026), alongside NVFP4 for the rest (slide 20). The rates follow the ladder: on the GB200, 2.5 PFLOP/s dense for FP16 and BF16, 5 for FP8 and FP6, 10 for FP4, and twice each with 2:4 structured sparsity$^\dagger$; the chart on the slide shows the dense figures and the one sparse figure NVIDIA headlines, 20 PFLOP/s. Per SM per clock, 10 PFLOP/s over 148 SMs at a clock near 1.9 GHz is about 17{,}800 FP4 FMA per clock per SM, or 4{,}450 per tensor core (derived, with the clock inferred; NVIDIA has not published the per-clock rate), against 1{,}024 dense FP16 FMA per clock per SM on Hopper's tensor cores.

The tensor memory is the structural novelty. On Hopper, a warpgroup's \texttt{wgmma} instruction read its A operand from registers or shared memory, its B operand from shared memory, and kept the accumulator tile in the warpgroup's registers, which meant that the register file both fed and drained the matrix unit and that large accumulator tiles consumed the registers that occupancy depends on. Blackwell adds a 256 KB tensor memory (TMEM) per SM, separate from the 256 KB register file and the shared memory, that holds the accumulators and optionally the A operand$^\dagger$. The new matrix instruction, \texttt{tcgen05.mma}, is issued by a single thread rather than a warpgroup, executes asynchronously, reads its operands from shared memory and TMEM, and writes its results to TMEM; the threads of a CTA (and, with the two-CTA mode, of a pair of CTAs across two SMs) then read the results out of TMEM with dedicated load instructions. Jarmusch and Chandrasekaran (2025) describe and measure this path; the practical effect is that a matrix pipeline no longer competes with the rest of the kernel for registers, and that the operand-delivery problem slide 15 describes is solved one level further down, with a dedicated store next to the multipliers. The same microbenchmarks measure the B200's FP64 at 44.8 TFLOP/s, consistent with the datasheet's 37--40, and confirm that there is no tensor-core FP64 path on Blackwell; FP64 runs on the CUDA cores at the rates slide 21 discusses.

\begin{terms}
\term{E2M1}{The 4-bit floating-point encoding with one sign, two exponent and one mantissa bit; its eight positive magnitudes are 0, 0.5, 1, 1.5, 2, 3, 4 and 6.}
\term{block scaling}{Storing one scale factor for a small block of low-precision values so that the block has the range of the scale and the values need only resolve relative magnitude; the mechanism that makes FP4 and FP6 usable.}
\term{microscaling, MX}{The Open Compute Project's family of block-scaled formats (MXFP8, MXFP6, MXFP4, MXINT8) with blocks of 32 values and an E8M0 (power-of-two) shared scale, specified in September 2023.}
\term{NVFP4}{NVIDIA's 4-bit block format: 16 E2M1 values share an E4M3 scale, and each tensor has one FP32 scale; 4.5 bits per value.}
\term{E8M0}{An 8-bit exponent-only encoding used as the MX block scale; it represents powers of two from $2^{-127}$ to $2^{127}$.}
\term{micro-tensor scaling}{The second-generation Transformer Engine's hardware application of block scales inside the tensor core, so that FP4, FP6 and block-scaled FP8 operands are expanded as they are consumed.}
\term{MXFP8}{The block-scaled FP8 format (32 values per E8M0 scale) that replaces Hopper's per-tensor FP8 scaling on Blackwell; used, for the attention and logit GEMMs, in NVIDIA's MLPerf Training 6.0 DeepSeek-V3 run.}
\term{tensor memory, TMEM}{A 256 KB per-SM memory on Blackwell, separate from the register file and shared memory, that holds tensor-core accumulators (and optionally an operand) so that matrix pipelines do not consume registers.}
\term{tcgen05.mma}{Blackwell's matrix-multiply instruction: issued by one thread, asynchronous, operands from shared memory and tensor memory, results into tensor memory; the successor of Hopper's warpgroup wgmma.}
\term{random Hadamard transform}{A rotation of a vector by a Hadamard matrix with random signs that spreads outlier values across all components, lowering the dynamic range a block must cover; used in the NVFP4 training recipe.}
\term{stochastic rounding}{Rounding a value up or down with probability proportional to its distance from each neighbor, so that rounding errors average to zero over many operations; used for gradients in low-precision training.}
\end{terms}

\begin{intuition}
A 4-bit number is a position on a sixteen-rung ladder; the block scale decides where the ladder stands. Sixteen neighbors share a ladder because their magnitudes are usually alike; the tensor-wide FP32 scale makes sure no ladder is placed where the E4M3 scale cannot reach. The tensor memory is the same idea applied to accumulators: give the matrix unit its own shelf so it stops borrowing the register file.
\end{intuition}

% ---------------------------------------------------------------------

\slide{19}{NVLink 5 and the GB200 NVL72: the rack as the unit}{72 GPUs and 36 Grace CPUs in one liquid-cooled rack; every GPU 1.8 TB/s from every other; 120 kW.}

\para{What the slide shows} A tall schematic of a rack on the left: nine compute trays at the top and nine at the bottom, each drawn with four GPU markers, and nine switch trays in the middle in a different color. Annotations beside the rack give the contents: each compute tray holds two Grace CPUs and four B200 GPUs; each switch tray holds two NVLink Switch chips; the 72 GPUs reach each other at 1.8 TB/s through 18 switch chips of 72 ports each, 130 TB/s in all; the spine carries over 5{,}000 copper cables with no optics inside the rack; the rack is direct liquid cooled and draws about 120 kW (Dell's figure). The right panel, ``The rack, in numbers and in purpose'', gives the rack's totals, the superchip, the reason for 72, and the alternatives and dates.

\para{Details} The GB200 NVL72 is the unit NVIDIA designed Blackwell around, and the unit its data-center revenue is counted in. Its contents, as NVIDIA's product page gives them$^\dagger$: 72 Blackwell GPUs and 36 Grace CPUs, 1{,}440 PFLOP/s of FP4 with sparsity (720 dense), 720 PFLOP/s of FP8 with sparsity, 2{,}880 TFLOP/s of FP64, 13.4 TB of HBM3e at 576 TB/s, 17 TB of LPDDR5X on the Grace CPUs (36 $\times$ 480 GB), and 130 TB/s of NVLink bandwidth (72 $\times$ 1.8 TB/s). The arithmetic is per-GPU figures times 72: 72 $\times$ 10 PFLOP/s FP4 dense is 720; 72 $\times$ 40 TFLOP/s FP64 is 2{,}880; 72 $\times$ 186 GB is 13.4 TB; 72 $\times$ 8 TB/s is 576 TB/s (derived). The press release of March 2024 said ``1.4 exaflops of AI performance'' and ``30 TB of fast memory''; the first is the sparse FP4 figure and the second the HBM plus LPDDR5X total, which is how such headlines should be decoded$^\dagger$.

The physical design is new. A compute tray is a 1U liquid-cooled module holding two GB200 superchips, that is, two Grace CPUs and four Blackwell GPUs, with each GPU's 18 NVLink 5 ports brought to the back of the tray; eighteen trays give the 72 GPUs. A switch tray holds two NVLink Switch chips of 72 ports each; nine trays give 18 chips and 1{,}296 ports, which is exactly 72 GPUs $\times$ 18 links. Every GPU has one link to every switch chip, so any GPU can send to any other through one switch hop at the full 1.8 TB/s, and the aggregate is 130 TB/s$^\dagger$. The trays connect through a backplane spine of copper cables, more than five thousand of them with a total length NVIDIA quotes as about two miles$^\dagger$; the choice of copper over optics inside the rack saves the power of the transceivers (NVIDIA cites about 20 kW per rack)$^\dagger$ at the cost of a limited reach, which is why the NVLink domain is one rack. The rack is direct liquid cooled, with cold plates on the GPUs, CPUs and switches, at a power the press variously puts at 120 to 132 kW$^\dagger$; a conventional air-cooled rack of servers is 10 to 20 kW, so the NVL72 required new facilities, which is one reason its deployment ran through 2025 rather than 2024.

Why 72, and not the eight of every DGX since 2016? The models of 2024--26 are mixture-of-experts transformers with hundreds of billions to trillions of parameters, served with long contexts. Two kinds of traffic dominate. Expert parallelism places different experts on different GPUs and routes each token to the experts it needs, an all-to-all exchange of activations at every layer; and the KV cache of a long context, the keys and values of every previous token, is read at every generated token and is spread across the GPUs that hold the model. Both patterns want every GPU to reach every other at the highest bandwidth available. With eight GPUs on NVLink and the rest of the model on other nodes over InfiniBand at 400 Gb/s per GPU (50 GB/s, a thirty-sixth of NVLink 5), the inter-node step dominates; with 72 GPUs on one NVLink fabric, a model of a trillion parameters in FP4 (about 0.5 TB of weights) fits in the rack's 13.4 TB with room for KV cache, and the all-to-all stays on NVLink. NVIDIA's claim for the NVL72 against the same number of H100s is up to 30$\times$ the inference throughput on trillion-parameter models at the same latency and up to 4$\times$ the training speed, both vendor figures that combine FP4, the larger NVLink domain and the Grace memory$^\dagger$; the direction is right even if the multiplier is a marketing one. The consequence slide 3 anticipated: the unit of design, of sale and of power planning moved from the board to the rack.

The alternatives: HGX B200 is the eight-GPU board for conventional air-cooled or liquid-cooled servers, with 180 GB per GPU, 37 TFLOP/s FP64 and 1{,}000 W per GPU; a second tier of NVLink Switch chips outside the rack can, in NVIDIA's design, join up to 576 GPUs into one NVLink domain, though NVIDIA cites no deployed 576-GPU domain and the scale-out network carries the traffic beyond a rack in practice$^\dagger$ (slide 28). The GB200 superchip itself is one Grace and two B200 GPUs on one board, with the Grace connected to each GPU over NVLink-C2C at 900 GB/s as on GH200 (slide 14). The dates: Azure announced the first GB200 systems on October 8, 2024; Dell delivered the first GB200 NVL72 racks to CoreWeave on November 18, 2024; volume ran through 2025; OpenAI and Oracle's Stargate site at Abilene, Texas, is planned around some 450{,}000 GB200 GPUs at 1.2 GW$^\dagger$ (slide 27).

\begin{terms}
\term{GB200 NVL72}{NVIDIA's rack-scale system: 18 compute trays (36 Grace CPUs, 72 Blackwell GPUs) and 9 NVLink switch trays joined by a copper spine into one 72-GPU NVLink domain; liquid cooled, about 120 kW.}
\term{compute tray}{A 1U liquid-cooled module of the NVL72 holding two GB200 superchips: two Grace CPUs and four Blackwell GPUs.}
\term{switch tray}{A module of the NVL72 holding two NVLink Switch chips; nine trays provide the 1{,}296 ports for 72 GPUs $\times$ 18 links.}
\term{spine}{The copper cable backplane that connects the compute and switch trays of the NVL72; over 5{,}000 cables, no optics.}
\term{NVLink domain}{The set of GPUs that can reach each other over NVLink at full bandwidth: 8 on an HGX board, 72 in the NVL72, 576 with a second switch tier on paper.}
\term{expert parallelism}{Placing the experts of a mixture-of-experts layer on different GPUs and routing each token's activations to the experts it selects; an all-to-all communication pattern at every layer.}
\term{mixture of experts}{A transformer whose feed-forward layers are replaced by several ``expert'' sub-networks of which only a few are used per token; it has many more parameters than it computes with, so it needs memory and all-to-all bandwidth.}
\term{KV cache}{The keys and values computed for every previous token of a context, stored so that generating the next token need not recompute them; it grows with context length and is read at every step.}
\term{direct liquid cooling}{Cooling chips with cold plates through which liquid circulates, removing heat without moving air; required above roughly 30--40 kW per rack.}
\term{HGX B200}{The eight-GPU Blackwell baseboard for servers, with 180 GB per GPU and 1{,}000 W per GPU; the non-rack-scale way to buy Blackwell.}
\end{terms}

\begin{intuition}
The NVL72 is the DGX-1's idea, eight GPUs that behave like one, scaled until the copper runs out: 72 GPUs on one switch fabric, one rack, one liquid loop. The model that justified it is a mixture of experts whose traffic is all-to-all; the engineering that made it possible is a switch chip with 72 ports and a backplane of five thousand cables.
\end{intuition}

% ---------------------------------------------------------------------

\slide{20}{Blackwell Ultra (2025): more HBM, 1.5$\times$ FP4, almost no FP64}{Same dies, re-balanced: 288 GB of HBM3e, 1.5$\times$ FP4, 2$\times$ exponentials, FP64 cut to 1.4 TFLOP/s.}

\para{What the slide shows} A horizontal bar chart compares the GB300 with the GB200 per GPU, each bar a ratio with GB200 as 1 and a dashed line at 1: the SFU exponential rate at 2$\times$; HBM capacity from 186 to 288 GB (1.55$\times$); dense FP4 from 10 to 15 PFLOP/s (1.5$\times$); power from 1{,}200 to 1{,}400 W (1.17$\times$); FP16 at 2.5 PFLOP/s, FP8 at 5 PFLOP/s and HBM bandwidth at 8 TB/s unchanged; INT8 from 5 to about 0.17 POPS (0.034$\times$); and FP64 from 40 to 1.4 TFLOP/s (0.035$\times$). The right panel gives the announcement and first-shipment dates, the changes in detail, the GB300 NVL72 rack totals, and the MLPerf Training 6.0 result.

\para{Details} Blackwell Ultra, the B300 and its superchip form GB300, was announced at GTC on March 18, 2025, and the first GB300 NVL72 racks were installed at CoreWeave, built by Dell, on July 3, 2025$^\dagger$. It is the ``Ultra'' refresh of the yearly cadence (slide 3): the same two 4NP dies and the same package, with a different configuration of what is enabled, a new memory stack and a new balance of units. NVIDIA's developer blog ``Inside NVIDIA Blackwell Ultra'' (August 2025) is the primary description$^\dagger$.

The changes, in the order of the chart. NVIDIA's blog gives the full implementation as 160 SMs, 8 GPCs of 20 SMs, 20{,}480 CUDA cores and 640 fifth-generation tensor cores, against 148 SMs measured on the B200; it adds that the SM count varies by SKU, and the number enabled on a shipping B300 has not been published$^\dagger$. The memory is 288 GB of HBM3e in eight 12-high stacks of 36 GB, against eight 8-high stacks of 24 GB; the bandwidth is unchanged at 8 TB/s because the pin rate and the number of stacks did not change, only the height. Dense FP4 rises to 15 PFLOP/s, 1.5$\times$ the GB200's 10, while the sparse FP4 figure stays at 20 PFLOP/s: NVIDIA's blog explains that the Ultra tensor core raises the dense NVFP4 rate and that the sparse path was not changed, so for the first time the ``2$\times$ with sparsity'' rule does not hold$^\dagger$. The special function units' exponential throughput doubles to 10.7 tera-exponentials per second; the exponential is the operation in the softmax of attention, and with attention layers consuming a growing share of inference time on long contexts the SFU had become a bottleneck in the attention kernel, which is why NVIDIA quotes a 2$\times$ attention-layer speedup for the Ultra$^\dagger$. The host interface moves to PCIe 6.0 (x16, 256 GB/s), and the matching NIC is ConnectX-8 at 800 Gb/s. The power rises to 1{,}400 W per GPU in the GB300 configuration.

What was removed is as important as what was added. The FP64 rate falls from 40 TFLOP/s to about 1.4 (NVIDIA's GB300 NVL72 page gives 100 TFLOP/s for the 72-GPU rack, or 1.39 per GPU; Lockwood's compilation of the datasheets gives 1.3 for GB300 and 1.2 for HGX B300)$^\dagger$. INT8 tensor throughput falls from about 5 POPS dense to about 0.17 POPS dense: the GB300 NVL72 page lists 24 POPS of INT8 with sparsity per rack, which is $24{,}000 / 72 / 2 \approx 170$ TOPS dense per GPU (the Hot Chips 37 table's 0.319 POPS is the sparse figure)$^\dagger$. Both are design choices: the FP64 units and the INT8 tensor paths were removed or reduced to make room for the denser FP4 path and the larger SFUs, within the same die area and process. The Ultra is therefore not a faster Blackwell but a differently balanced one, an inference part first, with the arithmetic that scientific computing needs (FP64) and that classical quantized inference used (INT8) traded for the arithmetic that transformer inference uses (FP4 and exponentials). Slide 21 takes up the FP64 consequence.

At rack scale the GB300 NVL72 keeps the GB200 NVL72's shape (72 GPUs, 36 Grace CPUs, 18 NVLink Switch chips, 130 TB/s of NVLink) and changes the totals: 1{,}440 PFLOP/s of FP4 with sparsity and 1{,}080 dense (72 $\times$ 15), 720 PFLOP/s of FP8 with sparsity, 100 TFLOP/s of FP64, 20.7 TB of HBM3e (72 $\times$ 288 GB; NVIDIA's page rounds it to 20 TB) at 576 TB/s, and 37 TB of ``fast memory'' counting the Grace LPDDR5X$^\dagger$ (the GTC 2025 announcement said 40 TB; the shipping page says 37). NVIDIA's marketing figure for the GB300 NVL72 is 1.5$\times$ the GB200 NVL72's ``AI performance'', which is the dense FP4 ratio$^\dagger$. The racks shipped into the same power and cooling infrastructure as GB200 NVL72, which NVIDIA presented as a feature: a data center designed for GB200 could take GB300 without rebuilding.

The measured result the slide cites comes from MLPerf Training 6.0, published by MLCommons on June 16, 2026. NVIDIA submitted results on a GB300 NVL72 cluster of 8{,}192 GPUs that trained the DeepSeek-V3 671B pretraining benchmark to its target in 2.02 minutes, and a GB200 cluster of the same size trained the Llama 3.1 405B benchmark in 7.07 minutes; NVIDIA's analysis puts the GB300 at 1.3 to 1.6$\times$ the GB200 per GPU across the benchmarks$^\dagger$. The DeepSeek-V3 submission used MXFP8, the block-scaled FP8 of slide 18, for the attention batched matrix multiplies and the logit GEMM, with the NVFP4 recipe elsewhere; NVFP4 training had first appeared in MLPerf Training v5.1 (November 2025) on Llama 3.1 405B. No Rubin results appeared in the 6.0 round$^\dagger$. This is the kind of number the deck prefers to a peak: a benchmark with a defined target accuracy, run at a scale that exercises the rack and the network, audited by a consortium, even though the configuration and the software were the vendor's.

\begin{terms}
\term{Blackwell Ultra, B300, GB300}{The 2025 refresh of Blackwell on the same dies: 160 SMs on the full die (per-SKU count not published), 288 GB of HBM3e, 1.5$\times$ dense FP4, 2$\times$ exponential throughput, FP64 and INT8 reduced, 1{,}400 W; GB300 is the GPU on the Grace Blackwell Ultra superchip.}
\term{12-high stack}{An HBM stack of twelve DRAM dies; with 3 GB dies it holds 36 GB, giving 288 GB over eight stacks at the same bandwidth as an 8-high stack of the same pin rate.}
\term{special function unit, SFU}{The SM unit that computes transcendental functions (exponential, reciprocal, square root, sine); its exponential rate limits the softmax in attention.}
\term{softmax}{The function that turns attention scores into weights, $e^{x_i}/\sum_j e^{x_j}$; one exponential per score, so attention over long contexts is exponential-bound once the matrix multiplies are fast.}
\term{PCIe 6.0}{The sixth generation of PCI Express, 64 GT/s per lane with PAM4 signaling, 256 GB/s bidirectional for 16 lanes; Blackwell Ultra's host interface.}
\term{ConnectX-8}{NVIDIA's network adapter of the Blackwell Ultra generation, 800 Gb/s per port, InfiniBand or Ethernet; one per GPU in GB300 systems.}
\term{GB300 NVL72}{The Blackwell Ultra rack: 72 B300 and 36 Grace, 1.08 EFLOP/s dense FP4, 100 TFLOP/s FP64, 20.7 TB HBM3e, 130 TB/s NVLink.}
\term{MLPerf Training}{MLCommons' benchmark suite for training speed: time to train a defined model to a defined accuracy on a defined dataset, with results audited and published twice a year; version 6.0 appeared in June 2026.}
\term{DeepSeek-V3}{A 671-billion-parameter mixture-of-experts language model published in December 2024, trained in FP8; the basis of a MLPerf Training 6.0 pretraining benchmark.}
\end{terms}

\begin{intuition}
An ``Ultra'' is the same chip with its dials turned for the customer who buys most of them: more memory per GPU for longer contexts and bigger models, more FP4 and more exponentials for attention, and nearly nothing left for FP64 and INT8. Read the bar chart as a statement of who the product is for.
\end{intuition}

% ---------------------------------------------------------------------

\slide{21}{FP64: what happened to double precision}{Up nearly 9$\times$ to H100, flat on Blackwell, nearly gone on Blackwell Ultra, back to 33 TFLOP/s on Rubin.}

\para{What the slide shows} A bar chart on a logarithmic axis gives the FP64 peak per GPU for six generations, with two bars each: FP64 on the CUDA cores (lighter) and FP64 on the tensor cores (darker). V100 7.8 and 7.8; A100 9.7 and 19.5; H100 34 and 67; GB200 40 and 40; GB300 1.4 and 1.4; Rubin 33 and 33. The legend sits below the axis. The right panel, ``Why, and what HPC does about it'', explains the trajectory, describes FP64 emulation on integer tensor cores, and lists which supercomputing centers chose which vendor.

\para{Details} Double precision is the format of most scientific simulation, and from Fermi (2010) to Hopper (2022) each NVIDIA compute die raised its FP64 rate. The previous lecture followed it to the V100's 7.8 TFLOP/s on 2{,}560 FP64 lanes at 1.53 GHz. Ampere kept 32 FP64 lanes per SM but had 108 SMs (9.7 TFLOP/s at 1.41 GHz) and, more importantly, added FP64 to the tensor cores: 64 FP64 FMA per clock per SM, 19.5 TFLOP/s for the A100, so that a DGEMM ran twice as fast as the vector units could manage and the HPL benchmark of a supercomputer became a tensor-core number$^\dagger$ (slide 7). Hopper doubled both again: 64 FP64 lanes per SM and 132 SMs give 34 TFLOP/s on the CUDA cores, and the fourth-generation tensor core gives 67 TFLOP/s of FP64 matrix throughput$^\dagger$ (slide 11). In two generations the FP64 matrix peak rose 8.6$\times$ (7.8 to 67), and the systems of slide 27 were bought on that curve: Alps, JUPITER, Eos.

Blackwell broke the curve. The GB200's FP64 is 40 TFLOP/s and the B200's 37, with no separate tensor-core FP64 rate: NVIDIA's datasheets list the same figure for ``FP64'' and ``FP64 Tensor Core''$^\dagger$, and the Blackwell microbenchmarks find no FP64 tensor path. Against Hopper's vector rate, 40 is an 18\% gain; against Hopper's tensor rate it is a 40\% loss (derived). The March 2024 announcement had said 45 TFLOP/s for GB200, which would have been a 33\% loss; the shipping parts are lower still. Blackwell Ultra then removed most of the FP64 hardware, leaving 1.2--1.4 TFLOP/s per GPU and 100 TFLOP/s for a 72-GPU rack, less than three H100s$^\dagger$. Rubin's specification page lists 33 TFLOP/s of FP64, vector and matrix the same, so the Rubin generation returns to roughly the H100's vector rate, not its tensor rate, with 2.4 PFLOP/s per Vera Rubin NVL72 rack$^\dagger$ (slide 25).

The reasons are economic and architectural together. A fused multiply-add in FP64 needs a 53 $\times$ 53-bit multiplier array, about seven times the area of an FP16 multiplier and some fifty times that of an FP4 multiplier, and on a tensor core the multipliers are most of the area. Every square millimeter spent on FP64 is a square millimeter not spent on FP4, HBM controllers or the SFUs that attention needs, and the customers who buy by the tens of thousands of GPUs buy for inference and training. NVIDIA's revenue in the quarter ending July 2026 was \$96.2 billion, nearly all of it data center$^\dagger$; supercomputing centers are a small share of that. Blackwell Ultra, as slide 20 says, is an inference part, and its FP64 cut is the clearest statement of what the die is for. In interviews during 2024 NVIDIA executives argued that FP64 demand is concentrated in a few applications and that emulation and mixed precision cover most of the rest; the HPC community's reaction, recorded in the trade press of May 2024, was less sanguine$^\dagger$.

Emulation is the technical answer. The Ozaki scheme (Ozaki, Ogita, Oishi and Rump, 2012) splits each FP64 operand into a sum of slices that are exactly representable in a lower-precision format, computes the products of the slices with error-free accumulation, and sums them back; on a tensor core with INT8 inputs and INT32 accumulation, a DGEMM becomes a set of INT8 GEMMs whose count depends on the accuracy required, typically seven to ten for full double precision. Because Blackwell's INT8 tensor rate (4.5--5 POPS) is more than a hundred times its FP64 rate, emulated DGEMM can exceed the native rate even at ten slices. The software arrived in two steps: cuBLAS 12.9 (May 2025) added FP32 emulation on BF16 tensor cores (the BF16x9 scheme) and said FP64 emulation was under development; FP64 emulation on the INT8 tensor cores shipped in cuBLAS with CUDA 13.0 Update 2 (October 2025), and NVIDIA's blog reports up to 2.3$\times$ the native FP64 DGEMM throughput on a GB200 NVL72 at full double-precision accuracy$^\dagger$. Two caveats follow from the deck's own numbers. First, Blackwell Ultra cut INT8 along with FP64 (slide 20) to about 0.17 POPS dense, so the Ultra has too little to emulate with, and the technique does not rescue it. Second, the ``200 TFLOP/s emulated DGEMM'' quoted for Rubin on Lockwood's compilation page has no NVIDIA source the deck could find and is marked unverified. Emulation also applies only to matrix multiplication: the stencils, sparse solvers and particle codes that make up much of HPC do their FP64 in the vector units, where the Blackwell rate is 37--40 and the Ultra rate 1.4.

The other answers are the CPU and the choice of vendor. Grace and Vera (slide 24) are server-class CPUs with high memory bandwidth, and the Vera Rubin NVL72 pairs 36 of them with its 72 GPUs; a code that needs FP64 can keep part of it on the CPU side. And the centers chose. Doudna, NERSC's 2026 system, is Vera Rubin, with Dell as the integrator$^\dagger$. Frontier (MI250X) and El Capitan (MI300A), the first two U.S.\ exascale machines, are AMD, as is ORNL's planned Discovery, built on the MI430X, a part AMD specified for ``true hardware FP64'' at 288 TFLOP/s per GPU$^\dagger$ (next lecture). Aurora is Intel. In the June 2026 TOP500 top five, LineShine (CPU only), El Capitan, Frontier and Aurora are not NVIDIA; JUPITER, with GH200 superchips, is$^\dagger$ (slide 27).

\begin{terms}
\term{FP64, double precision}{The 64-bit IEEE floating-point format (1 sign, 11 exponent, 52 mantissa bits), the working precision of most scientific simulation.}
\term{DGEMM}{Double-precision general matrix multiply, the BLAS routine at the heart of HPL and of dense linear algebra; the operation tensor-core FP64 and FP64 emulation accelerate.}
\term{FP64 tensor core}{The FP64 matrix path added on Ampere (64 FMA per clock per SM) and doubled on Hopper, giving 19.5 and 67 TFLOP/s; absent on Blackwell and Rubin, whose FP64 matrix rate equals the vector rate.}
\term{Ozaki scheme}{A method (Ozaki et al., 2012) for computing an accurate high-precision matrix product from products of lower-precision slices of the operands with error-free accumulation; the basis of FP64 emulation on integer tensor cores.}
\term{FP64 emulation}{Computing double-precision matrix products on INT8 (or other low-precision) tensor cores by slicing the operands; offered by cuBLAS on Blackwell from CUDA 13.0 Update 2 (October 2025), after FP32 emulation on BF16 in 12.9.}
\term{INT32 accumulation}{Summing integer products into a 32-bit integer without rounding error; the property that makes integer slices usable for exact emulation.}
\term{mixed-precision iterative refinement}{Solving a linear system in low precision and correcting the residual in high precision until the answer is accurate to double; the technique behind the HPL-MxP benchmark.}
\term{HPL}{The High-Performance Linpack benchmark, a dense FP64 LU factorization, whose result ranks the TOP500; it rewards FP64 matrix throughput.}
\term{MI430X}{AMD's HPC-oriented CDNA 5 GPU with 432 GB of HBM4 and a stated 288 TFLOP/s of FP64, due in 2027; chosen for ORNL's Discovery and France's Alice Recoque.}
\end{terms}

\begin{intuition}
For ten years FP64 rode along with the AI build-out because the same tensor core could do both; Blackwell is the point where the two markets wanted different silicon and the larger one won. Emulation gives DGEMM back on Blackwell, Rubin gives the vector rate back, and the centers that need more bought elsewhere.
\end{intuition}

% ---------------------------------------------------------------------

\slide{22}{Power: 300 W to 1,400 W per GPU, 120 kW to 600 kW per rack}{Perf per watt rose mostly by cutting precision; the rest was bought with power and liquid cooling.}

\para{What the slide shows} A bar chart of GPU power by product: V100 300 W, A100 400, H100 700, B200 1{,}000, GB200 1{,}200, GB300 1{,}400, and an empty bar for Rubin labeled ``n.p.'' Each bar carries a number above it, the dense FP16 tensor throughput per watt derived from the datasheets: 0.42 TFLOP/s per watt for V100, 0.78 for A100, 1.41 for H100, 2.25 for B200, 2.08 for GB200, 1.79 for GB300. The right panel, ``Reading the bars, and the rack'', gives the per-GPU and per-rack power history, the cooling and power-delivery changes, and the scale of the build-outs.

\para{Details} The thermal design power of NVIDIA's flagship data-center GPU rose from 300 W on the V100 SXM2 (2017) to 400 W on the A100 SXM4 (2020), 700 W on the H100 SXM5 (2022), 1{,}000 W on the HGX B200 and 1{,}200 W on the GB200 (2024--25), and 1{,}400 W on the GB300 (2025)$^\dagger$: 4.7$\times$ in eight years. Rubin's TDP is not published; figures of 1{,}800 to 2{,}300 W come from supply-chain reporting (the only sourced number is an analyst's 2{,}300 W from September 2026) and are marked unverified in the deck$^\dagger$. The previous lecture's observation that power ``doubled'' from G80 to V100 over eleven years is no longer the trend; the power per GPU is now rising faster than the process is improving, which is the definition of buying performance with watts.

The numbers above the bars put the gain in context. Dense FP16 tensor throughput per watt is 125 TFLOP/s over 300 W on V100, 0.42; 312 over 400 on A100, 0.78; 989 over 700 on H100, 1.41; 2{,}250 over 1{,}000 on B200, 2.25; 2{,}500 over 1{,}200 on GB200, 2.08; 2{,}500 over 1{,}400 on GB300, 1.79 (all derived from datasheet peaks). At a fixed precision, then, the efficiency of the chip improved 4$\times$ from Volta to Blackwell and slipped on the Ultra, whose extra watts went to memory and SFUs rather than FP16. The headline gains come from the precision ladder: at the fastest precision each product offers, V100's 0.42 TFLOP/s per watt in FP16 becomes GB300's 15 PFLOP/s over 1{,}400 W, 10.7 TFLOP/s per watt in FP4, a 25$\times$ improvement (derived), of which the factor 4 from halving the word length twice is the larger part. The slide's title says the same thing: performance per watt rose mostly by cutting precision; the rest was bought with power and cooling.

Per rack, the progression is DGX-1 (eight P100 or V100, 2016--17) at 3.2 kW maximum system power, DGX A100 at 6.5 kW, DGX H100 at 10.2 kW, and GB200 NVL72 at about 120 kW (Dell's figure; NVIDIA's pages give no rack power, and some reports say 132 kW)$^\dagger$. Air can cool perhaps 30 to 40 kW in a rack with heroic airflow; above that the heat must leave in liquid. The NVL72 is direct liquid cooled throughout, and its deployment waited on facilities that could supply chilled water at tens of kilowatts per rack, which is part of why its volume came in 2025 rather than 2024 (slide 19). The Kyber rack announced at GTC 2025 for Rubin Ultra was specified at 600 kW, with compute blades mounted vertically to pack 576 dies (144 four-die packages) into one rack; at GTC 2026 NVIDIA renamed it Kyber NVL144 under the package-counting convention and showed renderings that an analyst (Lockwood) read as half the density (two chassis instead of four), but NVIDIA made no density statement and did not restate the power$^\dagger$ (slide 26). Whether a 600 kW rack arrives, the direction is set: the rack has become a power-delivery and cooling problem first.

The Vera Rubin NVL72 shows the engineering that follows. It accepts inlet water at 45\,$^\circ$C, warm enough that a data center can reject the heat to the outside air without chillers in most climates, which lowers the facility's power overhead; it distributes 800 V DC to the trays, which halves the current and the copper of a 400 V system for the same power and removes a conversion stage; and its compute trays are cable-free and hot-swappable, with blind-mate connectors for power, liquid and NVLink$^\dagger$ (slide 25). Each of these is a response to a number on this slide.

The last paragraph of the panel is the scale. A cluster of 100{,}000 GPUs at 1.2 kW each is 120 MW of GPUs; with CPUs, switches, storage and a cooling overhead of 20--30\%, it is on the order of 150 MW (derived), the output of a mid-sized gas turbine plant and the load of a city of a hundred thousand homes. The Stargate site at Abilene is planned at 1.2 GW for some 450{,}000 GB200 GPUs$^\dagger$ (slide 27). The constraint most 2026 build-outs cite in public is power and the time to connect it, not the supply of GPUs; this is why the vendors now quote tokens per second per megawatt, and why a perf-per-watt figure derived from the datasheet, as above, is the number to carry into any comparison between generations or vendors.

\begin{terms}
\term{thermal design power, TDP}{The sustained power a part is designed to dissipate and the cooling system must remove; NVIDIA's data-center GPUs went from 300 W (V100) to 1{,}400 W (GB300).}
\term{performance per watt}{Throughput divided by power; here dense FP16 tensor TFLOP/s per watt from the datasheets, 0.42 on V100 to 2.25 on B200.}
\term{liquid cooling, direct-to-chip}{Removing heat with cold plates on the chips and a liquid loop rather than with air; necessary above roughly 30--40 kW per rack and standard on the NVL72.}
\term{Kyber}{NVIDIA's rack design announced for Rubin Ultra at GTC 2025 with vertical compute blades, 576 dies (``NVL576'') and 600 kW; renamed Kyber NVL144 at GTC 2026, where renderings appeared less dense (analyst reading) and the power was not restated.}
\term{800 V DC distribution}{Supplying rack power as 800 volt direct current instead of 400 V or 48 V, halving current and conductor size for the same power; used in the Vera Rubin NVL72.}
\term{inlet water temperature}{The temperature of the coolant supplied to a liquid-cooled rack; 45 $^\circ$C allows heat rejection without chillers and lowers facility overhead.}
\term{power usage effectiveness, PUE}{Total facility power divided by IT power; the cooling and distribution overhead, 1.1--1.3 in modern data centers.}
\term{tokens per second per megawatt}{The inference-throughput efficiency measure NVIDIA and its customers use for racks and data centers; the successor of FLOP/s per watt as the figure of merit.}
\end{terms}

\begin{intuition}
Each bar is a decision to spend more electricity for more throughput; the small numbers above the bars show that at a fixed precision the chips got only 4$\times$ better in eight years, and the rest of the 25$\times$ came from shorter numbers. The rack, the water temperature and the voltage are now part of the architecture because the power is.
\end{intuition}

% =====================================================================
\setstage{trn}
\section{Rubin and Vera (slides 23--26)}

\slide{23}{Rubin (2026): HBM4, 3 nm, and the third dual-die GPU}{GTC 2025 specs, CES 2026 production, first racks August 2026, spec sheet revised in September 2026.}

\para{What the slide shows} A package drawing in the same style as slide 17: two reticle-sized dies, each labeled ``reticle die, TSMC 3 nm'' over a grid of forty SM markers and an ``L2 (n.p.)'' bar, a die-to-die band between them, and eight HBM4 stacks along the edges. The caption records what is and is not known: 336 billion transistors, a figure from the CES 2026 keynote slides; 288 GB of HBM4 at 19.2 TB/s on the specification page as revised in September 2026, after 13 TB/s was announced in 2025 and 22 in January 2026; a PCIe 5 host interface; SM count and TDP not published. Two stat tiles: 35 PFLOP/s of dense NVFP4 for training, and 19.2 TB/s from eight stacks totaling 288 GB. The right panel, ``What NVIDIA has said'', gives the compute peaks and their ratios to GB300, what is known of the inside, the links, and the shipping status as of November 2026.

\para{Details} Rubin, named for the astronomer Vera Rubin together with the Vera CPU, is the architecture that follows Blackwell on the yearly cadence of slide 3. Its public history has four dates. At Computex in June 2024, NVIDIA put ``Rubin'' and ``Rubin Ultra'' on a roadmap for 2026 and 2027 and announced the cadence itself$^\dagger$. At GTC on March 18, 2025, it gave the first specifications: a two-die GPU with 288 GB of HBM4 at 13 TB/s, 50 PFLOP/s of NVFP4 for inference, NVLink 6 at 3.6 TB/s per GPU, and a rack it then called Vera Rubin NVL144$^\dagger$. At CES on January 5, 2026, it declared Rubin ``in full production'' and presented the generation as six chips (slide 24), with the HBM4 bandwidth now 22 TB/s and the rack renamed Vera Rubin NVL72$^\dagger$. Hot Chips 38 in August 2026 carried the first architecture talk, ``NVIDIA Rubin GPU: Driving the Era of Agentic AI'' (Mandal, Dash and Dimitrov), and a companion talk on the Vera CPU$^\dagger$. The first revenue racks shipped in August 2026, and in September NVIDIA's specification page was revised downward in two places (slide 25).

What NVIDIA has published about the chip$^\dagger$. The package carries two reticle-limited dies on TSMC's 3 nm-class process (the exact variant, N3P, is reported but not confirmed by NVIDIA) and eight HBM4 stacks of 36 GB, 288 GB in all, the same capacity as Blackwell Ultra. The transistor count of 336 billion appears on the CES keynote slides but not in the press release, so the deck marks its source; it is 1.6$\times$ Blackwell's 208 billion, a gain consistent with the node change on roughly the same silicon area (derived). The compute peaks on the specification page are, per GPU: 50 PFLOP/s of NVFP4 for inference and 35 for training, 17.5 PFLOP/s of FP8 and FP6, 4 PFLOP/s of FP16 and BF16, 2 PFLOP/s of TF32, 130 TFLOP/s of FP32, and 33 TFLOP/s of FP64, with the FP64 vector and matrix figures the same. NVIDIA does not label the two NVFP4 figures as sparse and dense; the ratio 50/35 is not the 2:1 of structured sparsity, and the training figure is the one that corresponds to the dense rate used in the ladder of slide 10, so the deck quotes 35 PFLOP/s as the dense NVFP4 peak and 50 as the inference figure, which includes the effect of Rubin's new sparsity and compression features. Against the GB300 the ratios are 35/15 $=$ 2.3$\times$ in dense FP4, 17.5/5 $=$ 3.5$\times$ in FP8, 4/2.5 $=$ 1.6$\times$ in FP16 and 33/1.4 $=$ 24$\times$ in FP64 (derived); the FP8 gain being larger than the FP4 gain says that the tensor core was rebalanced toward 8-bit training, and the FP64 figure says that the Ultra's cut was a one-generation decision (slide 21).

Of the inside, three things are stated and the rest is not. NVIDIA describes a third-generation Transformer Engine with ``adaptive compression'' of operands, which from the Hot Chips presentation appears to be a hardware path that compresses weights and activations in memory and expands them in the tensor core, reducing the bytes the HBM must deliver per operation; a more general form of 2:4 structured sparsity than Ampere's, which required exactly two nonzeros in every group of four; and tensor cores NVIDIA counts as the sixth generation$^\dagger$. The SM count, the register-file and shared-memory sizes, the L2 size and the clock have not been published, and the microbenchmark papers that filled those gaps for Blackwell had not appeared for Rubin by November 2026; the deck's package drawing therefore shows an L2 without a size and the slide's caption says so. The TDP is likewise unpublished, with 1{,}800 to 2{,}300 W reported in the press and marked unverified.

The links are NVLink 6 at 3.0 TB/s per GPU on the current specification page (3.6 TB/s was announced in 2025 and at CES 2026, and the Hot Chips talk still described the switch as providing 3.6 TB/s all-to-all per GPU), NVLink-C2C at 1.8 TB/s to the Vera CPU (slide 24), and PCIe 5 as the host interface, a step back from Blackwell Ultra's PCIe 6 that the Hot Chips presentation confirmed$^\dagger$.

The status in November 2026 is the last paragraph of the panel. NVIDIA said at CES in January that Rubin was in full production and repeated it at GTC Taipei on May 31, 2026, with ``production shipments starting this fall''; the first revenue Vera Rubin NVL72 racks shipped in August 2026, with Microsoft the first customer; the second-quarter fiscal 2027 results of August 26, 2026, reported racks running at CoreWeave, Google Cloud, Microsoft Azure, Oracle Cloud and Nebius; CoreWeave announced a production cluster for Cognition in September; and NVIDIA guided that Vera Rubin would be about a fifth of data-center revenue, roughly \$20 billion, in the quarter from August to October 2026$^\dagger$. That quarter's results were due in mid-November, after this deck was prepared. A report in April 2026 by TrendForce had described Rubin schedule pressure from HBM4 qualification, the move from ConnectX-8 to ConnectX-9, and power and cooling, and cut its estimate of Rubin's share of 2026 high-end shipments from 29 to 22 percent$^\dagger$; the August shipments and the September specification revision are consistent with a product that shipped on roughly its announced schedule with some of its targets trimmed.

\begin{terms}
\term{Rubin}{NVIDIA's 2026 GPU architecture and GPU: two reticle dies on TSMC 3 nm, 288 GB of HBM4, NVLink 6, 35 PFLOP/s dense NVFP4; named with the Vera CPU for the astronomer Vera Rubin.}
\term{HBM4}{The fourth generation of high-bandwidth memory, with a 2{,}048-bit interface per stack (twice HBM3's); on Rubin eight stacks of 36 GB give 288 GB at 19.2 TB/s, 2.4 TB/s per stack.}
\term{TSMC 3 nm}{TSMC's N3 process family; Rubin's dies are the first NVIDIA data-center GPUs on it, after two generations on 4 nm-class nodes.}
\term{NVLink 6}{The sixth-generation GPU link: 3.0 TB/s per GPU on the current Rubin specification (3.6 TB/s announced), with a switch that also performs in-network arithmetic.}
\term{third-generation Transformer Engine}{Rubin's hardware and software for low-precision training and inference, adding ``adaptive compression'' of operands to the block scaling of the second generation.}
\term{adaptive compression}{NVIDIA's name for Rubin's compression of tensor operands between memory and the tensor core, reducing bytes per operation; details beyond the name had not been published by November 2026.}
\term{specification page}{A vendor's product web page with the datasheet figures; for Rubin it is the authoritative source and was revised in September 2026, so the retrieval date must be quoted with the numbers.}
\term{full production}{The vendor's statement that a part is being manufactured in volume, as NVIDIA said of Rubin in January and May 2026; it precedes revenue shipments by months.}
\term{Hot Chips}{The annual conference at which chip makers present architectures; Rubin and Vera were presented at Hot Chips 38 in August 2026.}
\end{terms}

\begin{intuition}
Rubin is the Blackwell recipe one node and one memory generation later: two dies, now on 3 nm; eight stacks, now HBM4; the same 288 GB; a tensor core turned toward FP8 training and FP4 inference; and FP64 restored to a working level. Read its numbers with their dates attached, because three of them changed between the announcement and the shipping page.
\end{intuition}

% ---------------------------------------------------------------------

\slide{24}{Vera, and the platform of six chips}{A CPU built for the rack: 88 custom cores, 1.8 TB/s to the GPUs; the generation ships as six chips.}

\para{What the slide shows} The upper left drawing is a Vera Rubin superchip: a Vera CPU box (88 Olympus cores, 176 threads, LPDDR5X, 96 PCIe/CXL lanes, six dies on an interposer) joined by a thick double arrow labeled ``NVLink-C2C 1.8 TB/s'' to two Rubin GPU boxes (288 GB HBM4, NVLink 6 at 3 TB/s each), with a caption placing two superchips in each compute tray and eighteen trays in a Vera Rubin NVL72 rack: 36 Vera and 72 Rubin, 144 dies. Below it, a row of six boxes names the chips NVIDIA announced together at CES on January 5, 2026: Vera (CPU), Rubin (GPU), NVLink 6 (switch), ConnectX-9 (1.6 Tb/s NIC), BlueField-4 (DPU) and Spectrum-6 (co-packaged-optics Ethernet switch). A note adds the seventh, the Groq 3 LPX, and NVLink Fusion. The right panel, ``Vera against Grace, and why a CPU'', compares the two CPUs and explains why NVIDIA builds one.

\para{Details} Grace (slide 14) was NVIDIA's first server CPU, built from 72 licensed Arm Neoverse V2 cores. Vera is its second, and the first with NVIDIA's own core. The Hot Chips 38 talk ``NVIDIA's Vera CPU'' (Evans and Xekalakis, August 2026) and the CES material give the following$^\dagger$. The core, called Olympus, implements the Armv9.2 architecture with a 10-wide decode and runs two threads per core in a mode NVIDIA calls spatial multithreading, so 88 cores present 176 threads. The package holds six dies on an interposer: a monolithic compute die carrying the 88 cores and chiplets for memory and I/O. The memory is LPDDR5X, as on Grace, on SOCAMM modules (small outline compression-attached memory modules, a removable form factor NVIDIA introduced for Grace Blackwell; the SOCAMM detail for Vera is from press coverage of CES, not NVIDIA's release)$^\dagger$, with the capacity and bandwidth per CPU not stated on the specification page; the rack figure of 54 TB of LPDDR5X over 36 CPUs implies 1.5 TB per Vera (derived). The I/O is 96 lanes of PCIe or CXL, and the link to the GPUs is NVLink-C2C at 1.8 TB/s, twice Grace's 900 GB/s. A Vera serves two Rubin GPUs on a superchip; the rack total of 65 TB/s of C2C over 36 CPUs is 36 $\times$ 1.8 TB/s, which suggests the 1.8 TB/s is the CPU's total and each GPU has 0.9 TB/s to CPU memory, the same as a Grace Hopper GPU had but to a CPU with more cores and a larger memory (derived; NVIDIA does not state the split).

Why NVIDIA builds a CPU at all is the panel's main argument, and it has four parts. First, capacity: a Rubin GPU holds 288 GB; a Vera holds about 1.5 TB; the rack holds 20.7 TB of HBM4 and 54 TB of LPDDR5X (slide 25). The models and their working state (KV caches for many long contexts, embedding tables, optimizer state, checkpoints being written) exceed HBM, and LPDDR5X at a fraction of HBM's cost per byte is where they spill, reachable over C2C at a rate PCIe cannot provide. Second, placement: data preparation, tokenization, retrieval and, in the agentic workloads NVIDIA emphasized in 2026, the orchestration logic between model calls run on CPU cores, and running them beside the GPUs on a coherent link avoids a trip through a host server's PCIe. Third, ratio: an x86 server pairs two CPU sockets with eight GPUs through PCIe switches; a superchip pairs one CPU with two GPUs directly, and the rack's CPU-to-GPU ratio is fixed by design rather than by the server vendor's board. Fourth, control: the CPU's memory controllers, its link and its power are NVIDIA's to co-design with the GPU, which is the point of the next paragraph.

At CES 2026 NVIDIA presented the generation as ``six new chips, one AI supercomputer''$^\dagger$: the Vera CPU; the Rubin GPU; the NVLink 6 switch chip, which besides switching performs 130 TFLOP/s of in-network arithmetic for collective operations (slide 25); ConnectX-9, a network adapter at 1.6 Tb/s per port for scale-out, one per GPU; BlueField-4, a data-processing unit for storage, security and infrastructure services; and Spectrum-6, an Ethernet switch with co-packaged optics, the optical transceivers integrated on the switch package rather than plugged into its front panel, for the scale-out network. Each previous generation had introduced these parts separately over a year or more; presenting them as one platform is the statement that the product is the rack and the data center, and that a customer who buys the GPU is expected to buy the rest. At GTC in March 2026 a seventh part joined: the Groq 3 LPX, a rack of SRAM-based ``language processing units'' for the decode phase of inference, the result of NVIDIA's licensing and hiring agreement with Groq in December 2025, valued at about \$20 billion$^\dagger$ (slide 26). NVLink Fusion, announced in May 2025, points the other way: it licenses NVLink to partners (MediaTek, Marvell, Alchip, Astera Labs, Synopsys, Cadence, Fujitsu, Qualcomm, and Intel from September 2025) so that a custom CPU or accelerator can sit on NVIDIA's scale-up fabric, a response to the UALink consortium's open alternative$^\dagger$ (next lecture).

The comparison with Grace is in the panel's first item. Grace: 72 Neoverse V2 cores, 480 GB of LPDDR5X at up to 512 GB/s, NVLink-C2C at 900 GB/s, one GPU per CPU on GH200 and two on GB200. Vera: 88 Olympus cores with 176 threads, LPDDR5X of about 1.5 TB (derived from the rack), NVLink-C2C at 1.8 TB/s, two GPUs per CPU. The move from a licensed core to a custom one follows the path Apple, Amazon and Google took with their Arm server and client CPUs; NVIDIA's stated reasons are the per-core performance of the 10-wide design and the two-thread mode, and the ability to tune the core for the data-movement and orchestration work the GPUs hand it$^\dagger$. The lecture's reading is simpler: the CPU is now a component of the GPU platform, designed to the rack's requirements, and its specifications are published to the same (limited) extent as the GPU's.

\begin{terms}
\term{Vera}{NVIDIA's second server CPU (2026): 88 custom Olympus cores, 176 threads, LPDDR5X memory, 96 PCIe/CXL lanes, NVLink-C2C at 1.8 TB/s; one per two Rubin GPUs on a superchip.}
\term{Olympus core}{NVIDIA's first in-house CPU core, Armv9.2, 10-wide decode, two threads per core; 88 of them in Vera.}
\term{spatial multithreading}{NVIDIA's term for Vera's two-threads-per-core mode, in which the two threads are given separate execution resources rather than sharing a pipeline in time.}
\term{Armv9.2}{A version of the Arm 64-bit instruction-set architecture, with scalable vector and matrix extensions; the architecture Olympus implements.}
\term{SOCAMM}{Small outline compression-attached memory module, a removable LPDDR5X module format introduced by NVIDIA for Grace Blackwell and used for Vera.}
\term{CXL}{Compute Express Link, a coherent interconnect over PCIe signaling for attaching memory and accelerators to a CPU; Vera's 96 lanes can run PCIe or CXL.}
\term{superchip}{One CPU and one or two GPUs on a board joined by NVLink-C2C: GH200 (Grace plus one Hopper), GB200 (Grace plus two Blackwell), Vera Rubin (Vera plus two Rubin).}
\term{ConnectX-9}{NVIDIA's network adapter of the Rubin generation, 1.6 Tb/s per port; one per GPU gives 32.4 TB/s of scale-out bandwidth per NVL72 rack.}
\term{BlueField-4}{NVIDIA's data-processing unit (DPU) of the Rubin generation: a network adapter with Arm cores that runs storage, security and infrastructure services off the host CPU.}
\term{Spectrum-6}{NVIDIA's Ethernet switch of the Rubin generation, with co-packaged optics.}
\term{co-packaged optics, CPO}{Optical transceivers integrated onto the switch package instead of pluggable modules, cutting the power and the electrical reach problem of high-speed links.}
\term{Groq 3 LPX}{A rack of SRAM-based inference processors (language processing units) from NVIDIA's December 2025 agreement with Groq, positioned for the decode phase; in production from August 2026.}
\term{NVLink Fusion}{NVIDIA's 2025 program licensing NVLink to partners so custom CPUs and accelerators can join an NVLink fabric.}
\end{terms}

\begin{intuition}
Grace was NVIDIA's proof that a GPU company could ship a CPU; Vera is the CPU designed as a rack component from the start: two GPUs per socket, terabytes of cheap memory a cable-length away, and nothing between it and the GPUs but NVLink. The ``six chips'' are the same thought applied to the whole rack: every part a model's data touches is now designed together.
\end{intuition}

% ---------------------------------------------------------------------

\slide{25}{Vera Rubin NVL72 against the Blackwell racks}{Same 72-GPU shape; against GB200: 3.5$\times$ the dense FP4, 2.4$\times$ the HBM bandwidth, 1.7$\times$ the NVLink.}

\para{What the slide shows} A table with three columns, GB200 NVL72, GB300 NVL72 and VR NVL72, and rows for the first racks (November 2024, July 2025, August 2026), GPUs and CPUs (72 and 36 in each), dense FP4 (720, 1{,}080, 2{,}520 PFLOP/s), FP4 sparse or inference (1{,}440, 1{,}440, 3{,}600), dense FP8 (360, 360, 1{,}260), FP64 (2{,}880, 100, 2{,}400 TFLOP/s), HBM (13.4, 20.7, 20.7 TB), HBM bandwidth (576, 576, 1{,}400 TB/s), CPU memory (17, 17, 54 TB), total NVLink (130, 130, 216 TB/s), NVLink per GPU (1.8, 1.8, 3.0 TB/s), scale-out per GPU (400 Gb/s, 800 Gb/s, 1.6 Tb/s) and GPU power (1{,}200, 1{,}400 W, not published). A footnote gives the derivations and the source date. The right panel, ``What changed between announcement and shipment'', describes the September 2026 revision and NVIDIA's October statement, the NVL144-to-NVL72 renaming, what is new in the rack, and the lesson.

\para{Details} The Vera Rubin NVL72 keeps the shape of the Blackwell racks (18 compute trays of two superchips, 9 switch trays, 72 GPUs, 36 CPUs, one NVLink domain) and changes every component inside it. The table's figures for it are from NVIDIA's specification page as revised in September 2026$^\dagger$; the Blackwell figures are from the GB200 and GB300 NVL72 pages$^\dagger$. The arithmetic behind the rows: dense FP4 is 72 $\times$ 35 PFLOP/s $=$ 2{,}520; the inference figure is 72 $\times$ 50 $=$ 3{,}600; dense FP8 is 72 $\times$ 17.5 $=$ 1{,}260, against 360 for the Blackwell racks, which is half their 720 sparse FP8 figure (NVIDIA publishes only the sparse FP8 rack number for GB200 and GB300, so the dense value is derived); FP64 is 72 $\times$ 33 $=$ 2{,}376, which NVIDIA rounds to 2.4 PFLOP/s; HBM is 72 $\times$ 288 GB $=$ 20.7 TB at 72 $\times$ 19.2 TB/s $=$ 1{,}382 TB/s, which NVIDIA rounds to 1{,}400; CPU memory is 36 $\times$ 1.5 TB $=$ 54 TB; NVLink is 72 $\times$ 3.0 TB/s $=$ 216 TB/s; and the specification page lists 32.4 TB/s of networking bandwidth for the rack (the 2025 figure was 28.8 TB/s, which is 72 $\times$ 1.6 Tb/s counted in both directions; how the current figure is counted is not stated)$^\dagger$. The ratios in the subtitle follow: 2{,}520/720 $=$ 3.5$\times$ dense FP4 against GB200 (2.3$\times$ against GB300), 1{,}400/576 $=$ 2.4$\times$ HBM bandwidth, 216/130 $=$ 1.7$\times$ NVLink (derived).

Two rows deserve comment. FP64 at 2{,}400 TFLOP/s per rack is below the GB200 NVL72's 2{,}880 and far above the GB300's 100, so a center buying Vera Rubin gets back most, not all, of the double precision the Ultra removed (slide 21). And the HBM bandwidth row is the one that changed: at 22 TB/s per GPU the rack would have been 1{,}580 TB/s, 2.7$\times$ Blackwell; at 19.2 it is 2.4$\times$.

The revision is the panel's first item. Between September 1 and 11, 2026 (dated by archived copies), NVIDIA edited the Vera Rubin NVL72 page: the HBM4 bandwidth per GPU went from 22 to 19.2 TB/s, so the rack figure from 1{,}580 to 1{,}400 TB/s; NVLink 6 went from 3.6 to 3.0 TB/s per GPU, so the rack from 260 to 216 TB/s; and the HGX Rubin NVL8 eight-GPU board's NVFP4 figure went from 280 to 270 PFLOP/s. The compute peaks of the NVL72 rack did not change, but the eight-GPU HGX Rubin NVL8's dense peaks were lowered at the same time, NVFP4 from 280 to 270 PFLOP/s, FP8 from 140 to 130 and HBM bandwidth from 176 to 130 TB/s, a 3--7\% cut on compute and a quarter on memory$^\dagger$. The edit was reported by the press on October 26, and on October 29 NVIDIA issued a statement that the earlier numbers had been ``preliminary architectural targets'' and that the revised ones ``reflect what the system sustains at qualified, volume production''$^\dagger$. The deck quotes the statement because it is the most direct account a vendor has given of why announced and shipping numbers differ, and because it names the mechanism: a memory interface's sustained rate and a serializer's qualified rate are determined late, by yield and by signal integrity at volume, while a FLOP/s figure is fixed by the count of multipliers on the die and the clock, which were set at tape-out. That is the panel's last item, the lesson: HBM and SerDes rates move, FLOP/s rarely do, and a reader of a spec sheet should expect the bandwidth rows to be the ones with a date.

The name is the second item. At GTC 2025 the rack was ``Vera Rubin NVL144'', the 144 counting the dies: 72 packages of two dies each, the same silicon that CES 2026 called ``Vera Rubin NVL72'', counting packages as the Blackwell racks had. Both names describe the same rack; the deck uses NVL72 and notes the history so that a reader who meets ``NVL144'' in a 2025 source does not take it for a different system. The 2025 figures for the NVL144 (3.6 EFLOP/s of FP4 inference, 1.2 EFLOP/s of FP8 training, 13 TB/s of HBM4 per GPU, 260 TB/s of NVLink, 75 TB of ``fast memory'', 28.8 TB/s of ConnectX-9 bandwidth) are the same rack at an earlier stage of its specification; the inference figure held, the HBM figure went up and then down, and the CPU memory figure (75 TB total implied about 54 TB of LPDDR5X plus 20.7 TB of HBM) is unchanged$^\dagger$.

What is new in the rack beyond the chips is the third item. The NVLink 6 switch performs 130 TFLOP/s of arithmetic inside the switch, so that all-reduce and other collective operations are computed as data passes through rather than on the GPUs; this is the SHARP idea of Hopper's NVSwitch 3 (slide 14) at a scale and rate that makes it a stated specification$^\dagger$. ConnectX-9 gives each GPU 1.6 Tb/s to the scale-out network, four times the GB200's ConnectX-7 (slide 28). The power and cooling engineering (800 V DC distribution, 45\,$^\circ$C inlet water, cable-free hot-swappable trays with blind-mate connectors) was described on slide 22. The Hot Chips 2026 presentation also named a reliability feature, RIST, alongside the RAS engine inherited from Blackwell; the coverage describes it as in-system testing of deployed silicon, and no further detail had been published by November 2026$^\dagger$.

A note on what the table does not say. The GPU power for Rubin is not published, so a performance-per-watt comparison between the racks cannot be made from vendor data; the TrendForce report of April 2026 named power and cooling among the reasons for schedule pressure, and the rack's 800 V and 45\,$^\circ$C specifications suggest a power per rack above the Blackwell racks' 120--140 kW, but the deck does not put a number on it$^\dagger$.

\begin{terms}
\term{Vera Rubin NVL72}{NVIDIA's 2026 rack: 72 Rubin GPUs (144 dies) and 36 Vera CPUs, 2.52 EFLOP/s dense NVFP4 (3.6 inference), 20.7 TB of HBM4 at 1{,}400 TB/s, 216 TB/s of NVLink 6; called NVL144 at GTC 2025.}
\term{NVL144}{The GTC 2025 name of the Vera Rubin rack, counting dies rather than packages; dropped at CES 2026 in favor of NVL72.}
\term{HGX Rubin NVL8}{The eight-GPU Rubin baseboard, 270 PFLOP/s of NVFP4 after the September 2026 revision (280 before).}
\term{preliminary architectural target}{NVIDIA's phrase (October 2026) for a specification published at announcement before volume qualification; the Rubin HBM4 and NVLink 6 figures were such targets.}
\term{qualified, volume production}{The stage at which a part's sustained rates have been verified across manufactured units; the basis of the revised Rubin figures.}
\term{SerDes}{Serializer/deserializer, the circuit that drives a high-speed link such as NVLink; its qualified rate at volume sets a link's bandwidth and is finalized late in a product's development.}
\term{in-network compute}{Performing arithmetic (reductions, aggregation) inside a switch as data passes through; 130 TFLOP/s in the NVLink 6 switch, the successor of SHARP on NVSwitch 3.}
\term{RIST}{A reliability feature of Rubin named at Hot Chips 2026 and described in coverage as in-system testing of deployed silicon; details beyond the name were not published by November 2026.}
\term{blind-mate connector}{A connector that engages automatically when a module is slid into a chassis, so a tray can be inserted without cables being attached by hand; used for power, liquid and NVLink in the Vera Rubin trays.}
\end{terms}

\begin{intuition}
Same rack, new parts: 3.5$\times$ the FP4, 2.4$\times$ the memory bandwidth, 1.7$\times$ the fabric, and a lesson printed in the revision history. The numbers that moved were the ones set by yield (memory and link rates); the ones that held were set by the die. Quote a rack by its spec-sheet date.
\end{intuition}

% ---------------------------------------------------------------------

\slide{26}{After Rubin: CPX, Rubin Ultra, Feynman}{Announced for 2026--28: one product dropped, one redesigned, one replaced by a licensed design.}

\para{What the slide shows} A five-row table across the top: product, when announced, what it is, and its status in November 2026. Rubin CPX (September 2025; a Rubin die with 128 GB of GDDR7 and 30 PFLOP/s of NVFP4 for the prefill phase, in an ``NVL144 CPX'' rack of 8 EFLOP/s; absent from the GTC 2026 roadmap, with a revival reported but unconfirmed). Groq 3 LPX (December 2025 and March 2026; racks of SRAM-based language processing units for decode, from the Groq licensing deal; in full production per NVIDIA in August 2026). Rubin Ultra and the Kyber NVL144 rack (March 2025; four reticle dies per package, so 576 dies per rack and ``NVL576'' under the 2025 die-counting name; 1 TB of HBM4e, 100 PFLOP/s of FP4 per package, 15 EFLOP/s per 600 kW rack; due in the second half of 2027, with the GTC 2026 renderings of Kyber read by analysts as less dense and its power not restated). Feynman (2025 and 2026; the next architecture, with the Rosa CPU, NVLink 8 and an NVL1152 rack with co-packaged optics; 2028, process undisclosed). NVLink Fusion (May 2025; NVLink IP for partners' chips; design-ins announced, no shipping product cited). Two panels below: ``The design driver: disaggregated inference'' and ``How to read a roadmap, after 2024--26''.

\para{Details} The slide collects what NVIDIA has announced beyond Rubin and records what happened to each announcement in the twenty months since GTC 2025, because the pattern is itself the lesson.

Rubin CPX was announced on September 9, 2025, as ``a new class of GPU'' for massive-context inference$^\dagger$: a monolithic Rubin die with 128 GB of GDDR7 instead of HBM, 30 PFLOP/s of NVFP4, and hardware for video decoding and encoding, intended for the prefill phase of inference, in which a long prompt is pushed through the model once and the work is bound by arithmetic rather than by memory bandwidth, so that the expensive HBM of the main Rubin GPU would be reserved for the decode phase. NVIDIA described a Vera Rubin NVL144 CPX rack combining 144 CPX GPUs with a standard rack for 8 EFLOP/s of NVFP4 and 100 TB of fast memory, for the end of 2026. At GTC in March 2026 the CPX was absent from the roadmap, and the decode-side role it was to complement was given to the Groq 3 LPX; in September 2026 the analyst Ming-Chi Kuo reported that the CPX had been revived with 168 GB of HBM4 at 2{,}300 W for production in early 2027, a report the deck marks unconfirmed$^\dagger$.

The Groq 3 LPX is the product of a different kind of decision. In December 2025 NVIDIA entered a licensing and hiring agreement with Groq, the maker of SRAM-based inference processors, valued in reports at about \$20 billion$^\dagger$. Groq's ``language processing unit'' keeps model weights in on-chip SRAM across many chips rather than in HBM, which gives very high bandwidth per byte stored at the cost of capacity, and suits the decode phase, where each generated token requires reading the weights once with little arithmetic. At GTC 2026 NVIDIA presented racks of these processors as the Groq 3 LPX, the seventh chip of the platform (slide 24), and in August 2026 said they were in full production$^\dagger$. The two decisions together, CPX shelved and LPX adopted, show NVIDIA attacking the prefill-decode split from the decode side with licensed silicon rather than from the prefill side with its own.

Rubin Ultra is the 2027 refresh. At GTC 2025 NVIDIA described a package of four reticle-sized dies with 1 TB of HBM4e (sixteen stacks), 100 PFLOP/s of FP4 per package, and a rack called Kyber holding 144 such packages, 576 dies, in vertically mounted blades, for 15 EFLOP/s of FP4 inference and 5 EFLOP/s of FP8 training, 365 TB of fast memory, 4.6 PB/s of HBM bandwidth, 1.5 PB/s of NVLink 7, 115 TB/s of ConnectX-9, and 600 kW, in the second half of 2027$^\dagger$. At GTC 2026 NVIDIA listed the same rack as ``Kyber NVL144'', counting the 144 packages as it had begun to do for Vera Rubin NVL72 at CES; the renderings shown there have two chassis of blades where the 2025 design had four, which the analyst Glenn Lockwood read as half the density, but NVIDIA made no statement about density and did not restate the 600 kW$^\dagger$. The per-package figure of 100 PFLOP/s against Rubin's 50 (inference) is the four-die package's doubling; the rest of the 15 EFLOP/s is 144 packages against 72.

Feynman, named for the physicist, is the architecture after Rubin, placed in 2028 on the GTC 2025 roadmap and described further at GTC 2026 with a new CPU (Rosa), NVLink 8, and an NVL1152 rack using co-packaged optics for the scale-up fabric, the first time NVIDIA has proposed optics inside the NVLink domain; the process is undisclosed$^\dagger$. NVLink Fusion, from May 2025, is the program that licenses NVLink to partners' silicon; design wins were announced but no shipping product cited by November 2026$^\dagger$ (slide 24).

The design driver behind CPX and LPX is the first lower panel. Inference has two phases with opposite characters. Prefill runs the prompt through the model once; its matrix multiplies have a large batch dimension (the prompt's tokens) and are bound by arithmetic, so it wants FLOP/s and tolerates GDDR. Decode generates one token at a time; each step reads every weight and the whole KV cache for a handful of tokens of arithmetic, so it is bound by memory bandwidth and wants HBM or SRAM. NVIDIA's Dynamo framework, announced in March 2025, schedules the two phases on different pools of GPUs and moves the KV cache between them$^\dagger$; once the phases are separated in software, hardware specialized for each becomes possible, and CPX (prefill, GDDR7) and LPX (decode, SRAM) are the two ends of that split, with the HBM GPU serving both in the middle. The same split explains Blackwell Ultra's balance (slide 20) and the memory-first design of the HBM4 generation.

The second lower panel is the method. Across 2024--26, announced peaks moved 10--25\% before shipment, upward for Hopper and downward for Blackwell and Rubin (slides 3, 17, 25); one announced product disappeared from the roadmap within six months; a rack was redesigned a year after it was shown; a third-party design was adopted. None of this is unusual for a roadmap, but a lecture or a procurement must treat it accordingly: quote the date of every figure and prefer the specification page to the keynote; distinguish announced from shipping from measured, where measured means an audited benchmark (MLPerf) or a public list (TOP500), both of which lag the product by a year or more; and read a date such as ``H2 2027'' as a plan. The deck's own practice follows: Blackwell and Rubin figures are from datasheets and specification pages with dates, items NVIDIA has not published are marked as such, and the roadmap on this slide is a table of claims with their status rather than a forecast.

\begin{terms}
\term{Rubin CPX}{A Rubin-generation GPU announced in September 2025 with 128 GB of GDDR7 and 30 PFLOP/s of NVFP4 for the prefill phase of inference; absent from the March 2026 roadmap, revival reported but unconfirmed.}
\term{GDDR7}{The graphics DRAM generation after GDDR6X, soldered beside the GPU rather than stacked on an interposer; cheaper and lower-bandwidth than HBM, chosen for CPX.}
\term{prefill}{The first phase of inference, in which the prompt's tokens are processed through the model at once; arithmetic-bound.}
\term{decode}{The second phase of inference, generating one token per step by reading all weights and the KV cache; memory-bandwidth-bound.}
\term{disaggregated inference}{Running prefill and decode on separate pools of hardware, moving the KV cache between them; implemented by NVIDIA's Dynamo framework (2025).}
\term{Dynamo}{NVIDIA's open-source distributed inference framework (March 2025) for disaggregated serving, KV-cache routing and scheduling across GPU pools.}
\term{language processing unit, LPU}{Groq's inference processor, which holds model weights in on-chip SRAM across many chips for very high bandwidth; the basis of the Groq 3 LPX racks.}
\term{Rubin Ultra}{The 2027 refresh of Rubin: four reticle dies per package, 1 TB of HBM4e, 100 PFLOP/s of FP4 per package; the Kyber rack (NVL576 by dies in 2025, NVL144 by packages in 2026) was specified at 15 EFLOP/s and 600 kW; GTC 2026 renderings appeared less dense and the power was not restated.}
\term{HBM4e}{The extended-speed variant of HBM4 planned for Rubin Ultra, as HBM3e was of HBM3.}
\term{Feynman}{The NVIDIA architecture after Rubin, placed in 2028, with the Rosa CPU, NVLink 8 and an NVL1152 rack using co-packaged optics; process not disclosed.}
\term{roadmap}{A vendor's dated plan of future products; its figures are targets that have moved 10--25\% before shipment in recent generations.}
\end{terms}

\begin{intuition}
A roadmap is the vendor thinking aloud: CPX was a thought about prefill that was set aside, the LPX a thought about decode that was bought in, Kyber a thought about density that was halved. Keep the thoughts, since they say where the pressure is, but date every number and wait for the spec sheet.
\end{intuition}
% =====================================================================
\setstage{outc}
\section{Systems and scale (slides 27--28)}

\slide{27}{Where they run: NVIDIA systems, 2020--2026}{Centers bought A100, H100 and GH200 machines; the largest Blackwell installations do not run HPL.}

\para{What the slide shows} Two tables and a panel. The left table lists eleven NVIDIA-based systems by name, site and year, GPU type and count, and HPL result or rank: Selene (NVIDIA, 2020, A100, about 63 PFLOP/s, fifth in November 2020); Perlmutter (NERSC, 2021, A100, fifth in June 2021); Leonardo (CINECA, 2022, A100, fourth in November 2022); Eos (NVIDIA, 2023, 4,608 H100, 121 PFLOP/s, ninth in November 2023); Eagle (Microsoft Azure, H100, 561 PFLOP/s, seventh in June 2026); Alps (CSCS, 2024, GH200, 435 PFLOP/s, tenth in June 2026); Isambard-AI (Bristol, 2025, 5,448 GH200, 217 PFLOP/s, eleventh in June 2025); JUPITER (J\"ulich, 2025, about 24,000 GH200, 1.000 EFLOP/s, fifth in June 2026); Doudna (NERSC, 2026, Vera Rubin built by Dell, due late 2026); Solstice and Equinox (Argonne, announced October 2025, 100,000 and 10,000 Blackwell GPUs, AI systems, Equinox planned for the first half of 2026); and Stargate Abilene (OpenAI and Oracle, about 450,000 GB200 planned, 1.2 GW). A footnote marks the Selene, Perlmutter and Leonardo figures as coming from secondary reports and notes that Doudna's early-access system Cech (72 Grace CPUs, 144 Blackwell GPUs) arrived in January 2026. The right table is the top ten of the June 2026 TOP500 list with Rmax in PFLOP/s and the accelerator in each machine: LineShine (China, CPU only, 2,198), El Capitan (LLNL, AMD MI300A, 1,809), Frontier (ORNL, AMD MI250X, 1,353), Aurora (Argonne, Intel Max 1550, 1,012), JUPITER (FZJ, NVIDIA GH200, 1,000), HPC7 (Eni, AMD MI300A, 572), Eagle (Azure, NVIDIA H100, 561), HPC6 (Eni, AMD MI250X, 478), Fugaku (RIKEN, CPU only, 442) and Alps (CSCS, NVIDIA GH200, 435). The panel, ``Two worlds'', observes that no Blackwell or Rubin system appears in that top ten.

\para{Details} The left table traces the NVIDIA generations through the systems that national centers and the company itself built from them, and the pattern is one generation of GPU per system cycle. Selene was NVIDIA's own machine, built in 2020 from DGX A100 nodes in the SuperPOD reference design; the figure of about 63 PFLOP/s on HPL and the fifth place in November 2020 come from the TOP500 list as reported at the time, and the deck marks them as secondary because the notes' authors did not re-verify them against the list entry$^\dagger$. Perlmutter at NERSC (Lawrence Berkeley National Laboratory) was the first large A100 system at a U.S.\ national center; its GPU partition of A100 nodes placed fifth in June 2021$^\dagger$. Leonardo at CINECA in Bologna, one of the EuroHPC pre-exascale systems, used A100s in its booster partition and placed fourth in November 2022; by June 2026 it had slipped to twelfth$^\dagger$. Eos was NVIDIA's own Hopper system, 576 DGX H100 nodes for 4,608 GPUs, and its 121 PFLOP/s on HPL placed it ninth in November 2023$^\dagger$; it has since left the top ten. Eagle is Microsoft Azure's H100 cloud system, which at 561 PFLOP/s is the only hyperscaler machine in the June 2026 top ten and, with 14,400 H100s in its submitted configuration, the largest H100 HPL result$^\dagger$. Alps at CSCS in Lugano was the first large GH200 Grace Hopper system; its 435 PFLOP/s result dates from 2024 and placed tenth in June 2026$^\dagger$. Isambard-AI at the University of Bristol, with 5,448 GH200 superchips (168 in phase 1 and 5,280 in phase 2), is the United Kingdom's largest AI research system; it debuted at eleventh place on the June 2025 TOP500 with 216.5 PFLOP/s on HPL$^\dagger$. JUPITER at Forschungszentrum J\"ulich, built by Eviden on roughly 24,000 GH200 superchips in its Booster module, was the first European system to exceed one exaflop on HPL: 1.000 EFLOP/s in November 2025, when it ranked fourth, and the same figure in June 2026, when it ranked fifth$^\dagger$. Its listed power draw is 15.8 MW, the lowest of the five exascale systems on the list (derived from the list's power column).

The last three rows are systems of 2026 and beyond, and the table changes character with them. Doudna, NERSC's next system (also called NERSC-10), is a Dell machine on Vera Rubin; Berkeley Lab took delivery of its early-access system, Cech, in January 2026, a Grace Blackwell system of 72 Grace CPUs and 144 Blackwell GPUs rated at 5.76 PFLOP/s FP64 and 1.44 EFLOP/s NVFP4 by the laboratory, and the full system is scheduled for late 2026$^\dagger$. Solstice and Equinox are the two systems NVIDIA, Oracle and the Department of Energy announced for Argonne on October 28, 2025: Solstice with 100,000 Blackwell GPUs and Equinox with 10,000, the latter planned for the first half of 2026 (its delivery status had not been reported when these notes were prepared), described together as 2,200 exaflops of AI performance without a stated precision$^\dagger$. Stargate Abilene is the OpenAI and Oracle site in Texas planned for about 450,000 GB200 GPUs across eight buildings at 1.2 GW, of which two buildings were operating in September 2025 with the rest due by mid-2026$^\dagger$. None of these three is an HPL machine in the sense of the first eight rows: Doudna will run HPL because NERSC's users run FP64 science, but Solstice, Equinox and Abilene are sized in low-precision AI exaflops and gigawatts, and whether any of them submits to the TOP500 is a choice their owners have not announced.

The right table is the June 2026 TOP500 top ten and it makes three points. First, the top of the list is no longer an NVIDIA story: of the five exascale systems (those above 1 EFLOP/s on HPL), one is a Chinese CPU-only machine, LineShine at the National Supercomputing Center in Shenzhen, the first Chinese system to rank first since 2017; two are AMD (El Capitan on MI300A APUs at Livermore, Frontier on MI250X at Oak Ridge); one is Intel (Aurora on Data Center GPU Max 1550 at Argonne); and one is NVIDIA (JUPITER). Second, NVIDIA's entries in the top ten are GH200 and H100 machines, the generations of 2022--24; no Blackwell system of any size appears, eighteen months after GB200 racks began shipping. Third, the systems that do appear have power figures that put the FP64 result in context: El Capitan 29.7 MW, Frontier 24.6 MW, Aurora 38.7 MW, JUPITER 15.8 MW, so that JUPITER delivers 1 EFLOP/s at roughly 63 GFLOP/s per watt on HPL (derived), the best of the exascale machines on that metric.

The ``Two worlds'' panel states the reason for the absence of Blackwell. The TOP500 ranks systems by HPL, a dense FP64 linear solve, and submission is voluntary. The GB200 and GB300 racks deployed in 2025--26 were bought in the tens of thousands by cloud providers and AI laboratories for training and inference in FP8 and FP4; their owners neither run HPL nor submit results, and the machines are sized in exaflops of FP4 (a GB200 NVL72 rack is 1.44 EFLOP/s sparse FP4 and 2.9 PFLOP/s FP64). Slide 21 gives the other half of the explanation: Blackwell's FP64 is 37--40 TFLOP/s per GPU and Blackwell Ultra's about 1.4, so even a center that wanted an HPL number from a Blackwell Ultra rack would get a modest one (100 TFLOP/s per NVL72 rack, from slide 25). The November 2026 list, released at SC26 in the second half of November, may change the picture if a Blackwell or Rubin system is submitted; as of the date of these notes it had not been published.

\begin{terms}
\term{TOP500}{The twice-yearly list (June at ISC, November at SC) ranking submitted systems by their HPL result; voluntary, and therefore a record of what owners chose to measure.}
\term{HPL}{High-Performance Linpack, a dense FP64 LU factorization and solve; its achieved rate, Rmax, is the TOP500's ranking metric.}
\term{Rmax and Rpeak}{The measured HPL rate and the theoretical FP64 peak of a system; their ratio is the HPL efficiency.}
\term{exascale system}{A system exceeding $10^{18}$ FP64 FLOP/s on HPL; five existed on the June 2026 list.}
\term{DGX SuperPOD}{NVIDIA's reference design for a cluster of DGX nodes on a rail-optimized InfiniBand fabric; Selene and Eos were built to it.}
\term{Selene}{NVIDIA's in-house A100 supercomputer (2020), the first large DGX A100 SuperPOD.}
\term{Eos}{NVIDIA's in-house H100 supercomputer (2023), 4,608 GPUs, 121 PFLOP/s on HPL.}
\term{Eagle}{Microsoft Azure's H100 system, 561 PFLOP/s on HPL, the largest cloud system on the June 2026 list.}
\term{JUPITER}{The J\"ulich exascale system on about 24,000 GH200 superchips; Europe's first machine above 1 EFLOP/s on HPL (November 2025).}
\term{Alps}{The CSCS system in Lugano built on GH200 superchips; 435 PFLOP/s on HPL.}
\term{Isambard-AI}{The University of Bristol's AI research system of 5,448 GH200 superchips; 216.5 PFLOP/s, eleventh on the June 2025 TOP500.}
\term{Doudna}{NERSC's next system (NERSC-10), a Dell machine on Vera Rubin due in late 2026, preceded by the Cech early-access system.}
\term{Solstice and Equinox}{The two Blackwell AI systems announced for Argonne by NVIDIA, Oracle and DOE in October 2025, of 100,000 and 10,000 GPUs.}
\term{Stargate Abilene}{The OpenAI and Oracle data-center site in Texas planned for about 450,000 GB200 GPUs at 1.2 GW.}
\term{LineShine}{The CPU-only system at the National Supercomputing Center in Shenzhen that led the June 2026 TOP500 at 2.198 EFLOP/s.}
\term{El Capitan}{Livermore's exascale system on AMD MI300A APUs, 1.809 EFLOP/s, second in June 2026 after leading in 2024--25.}
\term{Frontier}{Oak Ridge's exascale system on AMD MI250X, the first machine above 1 EFLOP/s (June 2022); 1.353 EFLOP/s.}
\term{Aurora}{Argonne's exascale system on Intel Data Center GPU Max 1550 (Ponte Vecchio), 1.012 EFLOP/s at 38.7 MW.}
\end{terms}

\begin{intuition}
The TOP500 measures what its submitters measure. For ten years that was the same thing the AI buyers wanted; since Blackwell it is not, and the list has become a record of the HPC centers while the largest GPU installations in the world sit outside it. Read the two tables together: one says where NVIDIA's FP64 machines are, the other says that its largest customers no longer build FP64 machines.
\end{intuition}

% ---------------------------------------------------------------------
\slide{28}{Scale-up and scale-out: the NVLink domain grew from 8 to 72}{Per-GPU link bandwidth rose 19$\times$ since the DGX-1; the number of GPUs on one fabric rose 9$\times$, with 576 on paper.}

\para{What the slide shows} A scatter plot on a logarithmic vertical axis of the number of GPUs in one NVLink domain against year, with solid markers for systems that shipped in volume and hollow markers for those announced or without a known deployment: DGX-1 (2016, 8), DGX-2 (2018, 16), DGX A100 (2020, 8), DGX H100 (2022, 8) with the NVLink Switch System (2022, 256) as a hollow point above it, GB200 NVL72 (late 2024, 72), Vera Rubin NVL72 (2026, 72), and two hollow points for the Kyber rack, labeled ``Kyber NVL144 (576 dies)'' (2027), and the Feynman-generation NVL1152 (2028). The right panel, ``Two tiers of network'', gives the per-GPU scale-up bandwidth by generation (160 GB/s on P100, 300 on V100, 600 on A100, 900 on H100, 1,800 on Blackwell, 3,000 on Rubin), the per-GPU scale-out NIC speed (200 Gb/s with A100, 400 with H100, 800 with GB300, 1,600 with Rubin), and the ratio between them, about 15$\times$.

\para{Details} A GPU cluster has two networks, and the distinction between them is the organizing idea of this slide. The scale-up network is NVLink: a memory-semantic fabric over which a GPU loads and stores directly into another GPU's memory, with latency and bandwidth close to local HBM in the early generations and within a small factor of it since. The scale-out network is InfiniBand or Ethernet: a packet network reached through a network interface card, with a software stack (verbs, RDMA, NCCL) in the path, lower bandwidth per GPU, and in principle unlimited size. Every GPU in a modern system is attached to both, and the question that each generation answered differently is how many GPUs the scale-up tier should span.

The chart's history has four phases. The DGX-1 of 2016 connected eight P100s in a hybrid cube mesh of point-to-point NVLink 1.0 links, 160 GB/s per GPU, with no switch; a GPU could reach only its direct neighbors at full rate. The DGX-2 of 2018 introduced NVSwitch and connected 16 V100s through twelve switch chips so that every GPU reached every other at the full 300 GB/s. The DGX A100 (2020) and DGX H100 (2022) returned to eight GPUs, now behind six and then four NVSwitch chips, because eight was the number a 6U or 8U air-cooled chassis could hold at 400 and 700 W per GPU, and because the hyperscalers wanted a board (HGX) they could put in their own servers. Hopper's NVSwitch 3 and the NVLink Switch System could in principle join 256 GPUs across 32 nodes through a second tier of external switches, and NVIDIA's DGX SuperPOD documentation described such a configuration, but the deck plots it hollow because no public deployment of a 256-GPU domain is known and it never became the volume product$^\dagger$. The fourth phase is the rack: GB200 NVL72 (first racks November 2024) put 72 GPUs behind 18 NVLink Switch chips in one liquid-cooled rack, and Vera Rubin NVL72 (August 2026) kept the number, so the solid markers end at 72. The two hollow points at the right are announcements: the Kyber rack of Rubin Ultra, described at GTC 2025 as NVL576 (144 packages of four dies) for the second half of 2027 and listed as Kyber NVL144 at GTC 2026 under the package-counting name, where renderings that analysts read as less dense appeared without a restated power figure, and the Feynman-generation NVL1152 with co-packaged optics described for 2028$^\dagger$. The deck treats both as plans, following the lesson of slide 26.

The per-GPU bandwidths in the panel are the second axis of the story. NVLink bandwidth per GPU doubled or nearly doubled at every generation: 160 GB/s (P100, four links of 40), 300 (V100, six of 50), 600 (A100, twelve of 50), 900 (H100, eighteen of 50), 1,800 (Blackwell, eighteen of 100) and 3,000 (Rubin, as revised in September 2026 from the 3,600 announced). The ratio 3,000/160 is the 19$\times$ in the subtitle (derived). Over the same period the domain size went from 8 to 72, a factor of 9, so the aggregate NVLink bandwidth of a domain rose by about 170$\times$ (derived): from 8 GPUs at 160 GB/s to 72 at 3 TB/s, or 216 TB/s per Vera Rubin rack as NVIDIA states it. The switch itself gained a function along the way: NVSwitch 3 (Hopper) added SHARP, the Scalable Hierarchical Aggregation and Reduction Protocol inherited from Mellanox InfiniBand, so that an all-reduce's sums are computed inside the switch as data passes through, halving the traffic and removing a round from the collective; the NVLink 6 switch in the Vera Rubin rack carries that forward with 130 TFLOP/s of in-network arithmetic, a figure NVIDIA gave at Hot Chips 2026$^\dagger$.

Scale-out bandwidth rose more slowly. One NIC per GPU is the convention since the DGX A100 (eight ConnectX-6 at 200 Gb/s), and the per-GPU rate went to 400 Gb/s with ConnectX-7 on H100 and GB200 systems, 800 Gb/s with ConnectX-8 on GB300, and 1.6 Tb/s with ConnectX-9 on Vera Rubin (32.4 TB/s per rack, from slide 25). These are bits per second: 1.6 Tb/s is 200 GB/s, and the panel's ratio compares 3.0 TB/s of NVLink with 200 GB/s of NIC, about 15$\times$ (derived). The fabric behind the NICs is Quantum InfiniBand (HDR at 200 Gb/s, NDR at 400, XDR at 800) or Spectrum-X Ethernet, arranged as a rail-optimized fat tree: GPU $i$ of every node attaches to leaf switch $i$, so that the collectives of data-parallel training, which pair the same GPU index across nodes, stay within a leaf. The 100,000-GPU clusters of 2024--26 are scale-out constructions, thousands of NVL72 racks or HGX nodes on such a fabric.

Which parallelism goes on which tier follows from the bandwidths. Tensor parallelism splits a layer's weight matrices across GPUs and exchanges activations after every layer, several times per token; it needs the scale-up tier and was limited to the eight GPUs of a node until the rack. Expert parallelism in a mixture-of-experts model routes each token to a few of many experts spread across GPUs, an all-to-all exchange per layer, and is the workload that made 72 the design point (slide 19). Pipeline parallelism passes activations between stages once per micro-batch and data parallelism all-reduces gradients once per step; both tolerate the scale-out tier. The practical rule for a training job is therefore: tensor and expert parallelism inside the NVLink domain, pipeline and data parallelism across it, and the domain size sets the largest model that can be served or trained without the slower tier in its inner loop. This is why the next lecture's comparison of fabrics matters: AMD's Helios rack carries 72 GPUs on UALink over Ethernet at 3.6 TB/s per GPU, the UALink consortium's 1.0 specification (April 2025) targets up to 1,024 accelerators in one domain, and NVLink Fusion (May 2025) opens NVLink to other vendors' chips; the scale-up tier has become the contested layer.

NCCL, NVIDIA's collective communication library, is the software that spans both tiers. It discovers the topology (NVLink, NVSwitch, PCIe, NICs), builds rings and trees across it, and chooses algorithms per message size; a user of PyTorch or JAX sees only the bandwidth the collectives achieve. That is the slide's closing point: the two-tier hierarchy is invisible in the programming model and entirely visible in the performance model.

\begin{formulas}
Scale-up bandwidth per GPU: $N_{\text{links}} \times b_{\text{link}}$; Blackwell: $18 \times 100\ \text{GB/s} = 1.8$ TB/s; Rubin: 3.0 TB/s on the spec sheet (link count and rate not broken out). Domain aggregate: $N_{\text{GPU}} \times b_{\text{GPU}}$; Vera Rubin NVL72: $72 \times 3.0 = 216$ TB/s. Growth of the aggregate since the DGX-1: $(72 \times 3{,}000)/(8 \times 160) \approx 170\times$ (derived). Scale-out per GPU: $1.6\ \text{Tb/s} / 8 = 200$ GB/s; ratio to scale-up $3{,}000/200 = 15$ (derived).
\end{formulas}

\begin{terms}
\term{scale-up network}{The fabric that joins GPUs into one memory-semantic domain (NVLink and NVSwitch); high bandwidth per GPU, limited in size.}
\term{scale-out network}{The packet network (InfiniBand or Ethernet) that joins nodes or racks through NICs; lower bandwidth per GPU, unlimited in size.}
\term{NVLink domain}{The set of GPUs that can load and store into each other's memory over NVLink at full rate: 8 in an HGX node, 72 in an NVL72 rack.}
\term{hybrid cube mesh}{The DGX-1's switchless NVLink topology, in which each P100 had direct links to some of the other seven.}
\term{NVSwitch}{NVIDIA's NVLink crossbar chip; generations 1 (DGX-2), 2 (DGX A100), 3 (DGX H100, with SHARP), the NVLink Switch (Blackwell) and the NVLink 6 switch (Rubin).}
\term{NVLink Switch System}{Hopper's two-tier NVLink configuration joining up to 256 GPUs across nodes; documented, with no public deployment known.}
\term{SHARP}{Scalable Hierarchical Aggregation and Reduction Protocol: reductions computed inside the switch as data passes, inherited from Mellanox InfiniBand and added to NVSwitch 3.}
\term{in-network compute}{Arithmetic performed by a network switch on data in flight; 130 TFLOP/s in the NVLink 6 switch.}
\term{ConnectX}{NVIDIA's (formerly Mellanox's) NIC family: ConnectX-6 at 200 Gb/s, -7 at 400, -8 at 800, -9 at 1,600.}
\term{Quantum InfiniBand}{NVIDIA's InfiniBand switch line: HDR (200 Gb/s), NDR (400), XDR (800).}
\term{Spectrum-X}{NVIDIA's Ethernet switch and NIC combination tuned for AI traffic, the Ethernet alternative to InfiniBand for scale-out.}
\term{rail-optimized topology}{A fat tree in which GPU $i$ of every node attaches to the same leaf switch, so that same-index collectives stay within one switch.}
\term{tensor parallelism}{Splitting a layer's weights across GPUs with an activation exchange per layer; requires the scale-up tier.}
\term{expert parallelism}{Placing a mixture-of-experts model's experts on different GPUs and routing tokens between them all-to-all; the workload behind the 72-GPU rack.}
\term{pipeline parallelism}{Assigning consecutive layers to different GPUs and passing activations between stages; tolerates the scale-out tier.}
\term{data parallelism}{Replicating the model and splitting the batch, with a gradient all-reduce per step; the classic scale-out workload.}
\term{NCCL}{NVIDIA Collective Communications Library: topology-aware all-reduce, all-gather, all-to-all and related collectives over NVLink, PCIe and NICs.}
\term{UALink}{Ultra Accelerator Link, the consortium scale-up standard (1.0 specification April 2025) used by AMD's Helios rack over Ethernet.}
\term{NVLink Fusion}{NVIDIA's program (May 2025) licensing NVLink for partners' CPUs and accelerators.}
\term{Kyber}{The rack design announced for Rubin Ultra as NVL576 (dies) in 2025 and listed as Kyber NVL144 (packages) in 2026; the 2026 renderings appeared less dense (analyst reading) and the 600 kW was not restated.}
\end{terms}

\begin{intuition}
Think of the two tiers as a city: NVLink is the building, where everyone shares one floor plan and walks to each other's desks; the scale-out fabric is the road network between buildings. The last decade made the buildings bigger (8 rooms to 72) and the hallways faster (160 GB/s to 3 TB/s) much faster than it widened the roads, and a model is designed by deciding which conversations must happen inside one building.
\end{intuition}
% =====================================================================
\setstage{vc}
\section{Software (slide 29)}

\slide{29}{Software: what each generation asked of CUDA}{Every hardware feature since Volta arrived with a library that hides it; the programmer's view is a version number.}

\para{What the slide shows} A table with one row per generation and three columns, generation, CUDA version with its release date, and what arrived with it: Turing with CUDA 10.0 (September 2018: RTX and OptiX 6 (2019) ray tracing, TensorRT 5 with Turing INT8 and INT4, WMMA for INT8 and INT4, Nsight Compute); Ampere with CUDA 11.0 (June 2020, release candidate in May: TF32 as the default in cuDNN and the frameworks and opt-in in cuBLAS, \texttt{cp.async}, MIG, cuSPARSELt for 2:4 sparsity); Hopper with CUDA 11.8 (October 2022) and 12.0 (December 2022: clusters, DSMEM, TMA and \texttt{wgmma} in PTX, Transformer Engine, CUTLASS 3 and CuTe in 2023, FlashAttention-3 in 2024); Ada with CUDA 11.8 and 12 (FP8 in TensorRT, the DLSS 3 SDK, AV1 in NVENC); Blackwell with CUDA 12.8 (January 2025: \texttt{tcgen05} and tensor memory, NVFP4 in TensorRT-LLM, cuBLAS FP32 emulation in 12.9 and FP64 emulation in 13.0 Update 2 of October 2025, Dynamo in March 2025, cuTile and native CUDA Python at GTC 2025); a row for CUDA 13.0 (August 2025), which drops Maxwell, Pascal and Volta and requires the 580 driver; and Rubin with CUDA 13.x in 2026 (third-generation Transformer Engine, NVFP4 training recipes, NVLink 6 collectives). A footnote says the dates are release dates (general availability, or the release candidate where noted) and that the version-to-feature mapping is simplified. The right panel, ``The pattern, and the gap it leaves'', makes three points: features arrive wrapped in libraries; large matrix multiplies reach about 70--75\% of the dense peak, attention 60--75\%, end-to-end training 35--45\%; and above the libraries sit PyTorch, Triton, JAX, NCCL, the serving systems and Dynamo, with SYCL, HIP and Kokkos for HPC codes.

\para{Details} The table is a reading of the CUDA release history against the hardware history of slides 4--26, and its thesis is in the subtitle: each new unit in the silicon was made usable by a library, so that most programmers met the hardware as a version requirement and a performance change rather than as an instruction. The rows are worth taking in turn, because the mechanism differs each time.

\emph{Turing and CUDA 10.} Ray tracing reached developers through the RTX APIs (DirectX Raytracing, Vulkan ray tracing) and NVIDIA's OptiX 6.0 (February 2019) for offline rendering; a CUDA programmer never addresses an RT core. The integer tensor cores of slide 4 appeared in two places: the WMMA (warp matrix multiply-accumulate) API gained INT8 and INT4 fragment types, and TensorRT 5 added Turing support to the INT8 calibration path that TensorRT had carried since version 2 and 3 (2017), which measures the dynamic range of each layer's activations on sample data and chooses per-tensor INT8 scales; that is how a model trained in FP32 ran on a T4 at 130 TOPS with the small accuracy loss the slide-5 panel mentions$^\dagger$. Nsight Compute 1.0, the kernel profiler that replaced nvprof, shipped with CUDA 10.0 in the same month.

\emph{Ampere and CUDA 11.} The most consequential software decision of the period was invisible: cuDNN 8 made TF32 the default for FP32 convolutions on Ampere, and the deep-learning frameworks (PyTorch 1.7, TensorFlow 2.4) and NVIDIA's NGC containers turned it on by default for matrix multiplies as well, so that FP32 training code, unchanged, ran its matrix work on the tensor cores at 156 TFLOP/s instead of the 19.5 of the FP32 lanes (slide 7), with the operands rounded to a 10-bit mantissa; cuBLAS itself kept FP32 as its default math mode and applies TF32 only when the caller sets \texttt{CUBLAS\_TF32\_TENSOR\_OP\_MATH} or the \texttt{NVIDIA\_TF32\_OVERRIDE} environment variable$^\dagger$. PyTorch later exposed a switch (\texttt{allow\_tf32}) and in 2022 (version 1.12) changed its own default for matrix multiplies back to full FP32, keeping TF32 for convolutions; the episode is the clearest example of a precision decision made by a library on the user's behalf$^\dagger$. The asynchronous copy of slide 8 arrived as the PTX \texttt{cp.async} instruction and the CUDA \texttt{memcpy\_async} API; MIG arrived as a driver and \texttt{nvidia-smi} feature; 2:4 sparsity arrived as cuSPARSELt, a library that takes a pruned weight matrix and runs the sparse tensor-core path. CUDA 11.0 itself appeared as a release candidate on May 28, 2020 and in general availability in June, with Update 1 in July$^\dagger$.

\emph{Hopper and CUDA 12.} Support for the Hopper target (\texttt{sm\_90}) arrived in CUDA 11.8 (October 3, 2022), and CUDA 12.0 (December 6, 2022) completed it$^\dagger$. Hopper's new programming model (slide 13) is the one case where the hardware features were exposed directly in PTX and the CUDA language: thread block clusters and distributed shared memory are CUDA constructs (cluster launch attributes, \texttt{cluster.map\_shared\_rank}), and the TMA and \texttt{wgmma} are PTX instructions. Almost nobody uses them directly. CUTLASS 3 (2023) and its CuTe layout library wrapped the TMA-plus-\texttt{wgmma} pipeline into templates, and FlashAttention-3 (Shah et al., 2024) used the same primitives to bring attention from about 35\% of the H100's FP16 peak with FlashAttention-2 to about 75\% with warp specialization, pingpong scheduling of two warpgroups, and FP8 support. The Transformer Engine library (2022) is the FP8 story of slide 12: it owns the scale factors, the format choice and the casts, and exposes drop-in layer modules for PyTorch and JAX. NCCL gained SHARP support so that all-reduce used the NVSwitch 3 reductions of slide 14.

\emph{Ada.} The graphics die's compute features came through the same libraries as Hopper's (FP8 in TensorRT for L40S inference) plus the graphics SDKs: DLSS 3 frame generation through the DLSS SDK and the Streamline framework, AV1 encoding through NVENC's video SDK$^\dagger$.

\emph{Blackwell and CUDA 12.8.} Blackwell's fifth-generation tensor cores and tensor memory (slide 18) are reached through the new \texttt{tcgen05} family of PTX instructions, which allocate TMEM, issue asynchronous matrix multiplies from a single thread and commit results; CUTLASS added Blackwell kernels that use them. FP4 arrived as a quantization option: TensorRT Model Optimizer produces NVFP4 checkpoints with the block scales of slide 18, and TensorRT-LLM runs them, so a deployment engineer meets NVFP4 as a flag and an accuracy report. cuBLAS 12.9 (May 2025) added FP32 emulation on BF16 tensor cores and announced FP64 emulation as under development; FP64 emulation on the INT8 tensor cores for DGEMM on Blackwell (slide 21) shipped in CUDA 13.0 Update 2 (October 2025), selectable by an environment variable or an API call$^\dagger$. Dynamo (March 2025) is the inference-serving framework that schedules prefill and decode on separate GPU pools (slide 26). cuTile and the native CUDA Python work announced at GTC 2025 (CUDA Python bindings themselves date from 2021; the 2025 announcement was \texttt{cuda.core} and a Python-native kernel path) are a tile-based programming model and native Python support meant to let kernels be written at the level of tiles rather than threads, which is the level at which TMA and the tensor cores operate$^\dagger$.

\emph{CUDA 13.0 and Rubin.} CUDA 13.0 (August 2025) dropped offline compilation support for Maxwell, Pascal and Volta, so the V100 of the previous lecture is no longer a target of the current toolkit, and requires the R580 driver branch$^\dagger$. For Rubin the public software story as of November 2026 is the third-generation Transformer Engine with its ``adaptive compression'' and the NVFP4 training recipes that Abecassis et al.\ (2025) validated on Blackwell (slide 18); the CUDA 13.x releases of 2026 add the Rubin architecture target and NVLink 6 collectives in NCCL. The deck marks these as 2026 items without a specific minor version because the mapping was still moving.

The panel's second point, how much of the peak is reachable, deserves its own numbers because the datasheet never gives them. For a large BF16 matrix multiply on an H100 SXM, independent measurements reported by SemiAnalysis in December 2024 found about 720 TFLOP/s against the 989 dense peak, roughly 73\%; the same study measured AMD's MI300X at about 620 against its 1,307 marketed, a point the next lecture takes up$^\dagger$. Attention kernels reach 60--75\% with FlashAttention-3's techniques and less with earlier implementations. End-to-end training of a large language model is typically quoted at a model FLOP utilization (MFU) of 35--45\%: the ratio of the FLOPs the model's forward and backward passes require to the FLOPs the hardware could deliver in the elapsed time, with everything else (attention at lower efficiency, communication, optimizer steps, data loading, pipeline bubbles) in the gap. Inference at small batch sizes is bandwidth-bound (slide 15) and its FLOP/s utilization is low by construction; it is measured in tokens per second per GPU instead. The practical reading is that the dense peak overstates a training system's useful throughput by a factor of two to three and an inference system's by far more, which is why MLPerf, not the datasheet, is the comparison the next lecture uses.

The panel's third point is the stack above the libraries. PyTorch's \texttt{torch.compile} with the Inductor backend, OpenAI's Triton language and Google's JAX with XLA generate GPU kernels from high-level code, and Triton in particular has become the way most researchers write a custom kernel without CUDA C++; NCCL does the collectives for all of them. On the serving side, vLLM, SGLang and TensorRT-LLM implement continuous batching, paged KV caches and quantized formats, and Dynamo sits above them to disaggregate prefill from decode across a cluster. For HPC codes the portability layers (SYCL, HIP, Kokkos, OpenMP offload) exist and work, as the next lecture's portability studies measure, but the AI stack was written against CUDA first and ported later, which is the substance of the ``moat'' argument: the silicon can be matched on a datasheet more easily than the dozen libraries that make the datasheet's numbers reachable.

\begin{terms}
\term{CUDA toolkit}{NVIDIA's compiler (nvcc), runtime, driver API and libraries; versioned 10 (2018), 11 (2020), 12 (2022), 13 (2025).}
\term{PTX}{NVIDIA's virtual instruction set, the level at which new hardware features (cp.async, wgmma, tcgen05) first appear.}
\term{WMMA}{Warp matrix multiply-accumulate, the CUDA C++ API for tensor-core fragments, extended to INT8 and INT4 on Turing.}
\term{TensorRT}{NVIDIA's inference optimizer and runtime; INT8 calibration dates from versions 2 and 3 (2017), version 5 added Turing support, later versions FP8 and NVFP4.}
\term{INT8 calibration}{Measuring activation ranges on sample data to choose per-tensor integer scales for a model trained in floating point.}
\term{cuBLAS and cuDNN}{NVIDIA's dense linear algebra and deep-learning primitive libraries; cuDNN made TF32 the default for FP32 convolutions on Ampere, cuBLAS offers it as an opt-in math mode.}
\term{cuSPARSELt}{The library that runs 2:4 structured-sparse matrix multiplies on the tensor cores.}
\term{Nsight Compute}{NVIDIA's kernel profiler, successor to nvprof, reporting occupancy, throughput, cache and stall metrics; version 1.0 shipped with CUDA 10.0 (September 2018).}
\term{CUTLASS and CuTe}{NVIDIA's open-source template library for matrix multiply kernels and its layout algebra; version 3 (2023) wrapped Hopper's TMA and wgmma.}
\term{FlashAttention-3}{Shah et al.'s 2024 attention kernel for Hopper using warp specialization and TMA, reaching about 75 percent of the FP16 peak.}
\term{Transformer Engine}{NVIDIA's library for FP8 (and later NVFP4) training and inference, managing scales, formats and casts per layer.}
\term{tcgen05}{The Blackwell family of PTX instructions for the fifth-generation tensor cores and tensor memory.}
\term{TensorRT Model Optimizer}{NVIDIA's quantization and pruning toolkit that produces FP8 and NVFP4 checkpoints for TensorRT-LLM.}
\term{TensorRT-LLM}{NVIDIA's large-language-model inference library built on TensorRT, supporting FP8 and NVFP4 serving.}
\term{Dynamo}{NVIDIA's open-source distributed inference framework (March 2025) that disaggregates prefill and decode across GPU pools.}
\term{cuTile}{A tile-level programming model announced at GTC 2025 for writing kernels at the granularity the tensor cores and TMA operate on.}
\term{FP64 emulation}{Computing double-precision matrix products from integer slices on INT8 tensor cores (the Ozaki scheme), added to cuBLAS in CUDA 13.0 Update 2 (October 2025) for Blackwell.}
\term{MFU}{Model FLOP utilization: the FLOPs a model's training step requires divided by the FLOPs the hardware could have delivered in the measured time; typically 35 to 45 percent.}
\term{torch.compile and Inductor}{PyTorch's compiler path that fuses operations and generates Triton kernels.}
\term{Triton}{OpenAI's Python-embedded language for writing GPU kernels at the block level, now the common way to write custom kernels without CUDA C++.}
\term{vLLM and SGLang}{Open-source LLM serving systems implementing continuous batching and paged KV caches.}
\term{portability layer}{A programming model (SYCL, HIP, Kokkos, OpenMP offload) that targets several vendors' GPUs from one source; used by HPC codes, rarely by the AI stack.}
\term{the moat}{The argument that NVIDIA's advantage lies in the library stack that makes peaks reachable rather than in the silicon; examined in the next lecture.}
\end{terms}

\begin{intuition}
Hardware features are announced by NVIDIA and delivered by libraries, usually a year later and usually as a default or a flag. When you read a new architecture whitepaper, ask which library will carry each feature and how much of the peak it will expose; the answer, historically, has been about three quarters for the matrix multiply at the center and much less for everything around it.
\end{intuition}
% =====================================================================
\setstage{ink}
\section{Wrap-up: the table, the summary, the sources (slides 30--32)}

\slide{30}{The generations in one table}{Eight GPUs, 2017--2026: dense peaks, memory, links and power as the datasheets state them.}

\para{What the slide shows} A single table with eight columns, one per GPU, and sixteen rows. The columns are V100 (Volta), T4 (Turing), A100 (Ampere), H100 (Hopper), L40S (Ada), GB200 (Blackwell, per GPU), GB300 (Blackwell Ultra, per GPU) and Rubin. The rows are the year of first shipment (2017, 2018, 2020, 2022, 2023, 2024, 2025, 2026); process (12 nm, 12 nm, N7, 4N, 4N, 4NP, 4NP, 3 nm); transistors in billions (21.1, 13.6, 54.2, 80, 76.3, 208, 208, 336); dies (one each through L40S, two for Blackwell, Blackwell Ultra and Rubin); SMs enabled (80, 40, 108, 132, 142, 148, 160 on the full die, not published); FP32 lanes (5,120; 2,560; 6,912; 16,896; 18,176; 18,944; 20,480; not published); FP64 vector and tensor TFLOP/s (7.8; none; 9.7 and 19.5; 34 and 67; none; 40 and 40; 1.4; 33 and 33); dense FP16 tensor TFLOP/s (125, 65, 312, 989, 362, 2,500, 2,500, 4,000); dense FP8 (none through A100; 1,979; 733; 5,000; 5,000; 17,500); dense FP4 (10,000; 15,000; 35,000 on the last three); dense INT8 TOPS (none; 130; 624; 1,979; 733; 5,000; about 170; not published); memory (16 or 32 GB HBM2; 16 GB GDDR6; 40 or 80 GB HBM2; 80 GB HBM3; 48 GB GDDR6; 186 GB HBM3e; 288 GB HBM3e; 288 GB HBM4); bandwidth in TB/s (0.9; 0.3; 1.6 or 2.0; 3.35; 0.86; 8.0; 8.0; 19.2); L2 (6 MB; 4 MB; 40 MB; 50 MB; 96 MB; not published for the last three); NVLink per GPU in GB/s (300; none; 600; 900; none; 1,800; 1,800; 3,000); and power in watts (300; 70; 400; 700; 350; 1,200; 1,400; not published). The footnote says that the values are vendor datasheet figures for the named products, that all tensor peaks are dense (NVIDIA's headlines are usually the 2:4-sparse figures, twice these), that the GB200 and GB300 columns are per GPU (half a superchip), that the Blackwell SM counts come from microbenchmarks, that the T4 and Rubin transistor counts come from secondary sources, and that ``n.p.'' means not published as of November 2026.

\para{Details} The table is the deck's data in one place, and the way to read it is row by row, asking of each row what changed, by how much, and why. Several rows reward a closer look.

\emph{Transistors and dies.} From 21.1 billion on V100 to 336 billion on Rubin is 16$\times$ in nine years (derived), but the path is not smooth: 2.6$\times$ from Volta to Ampere (a full node, 12 nm to 7 nm), 1.5$\times$ to Hopper (7 to 4 nm), 2.6$\times$ to Blackwell with no node change at all (4N to 4NP) because the package went from one die to two, and 1.6$\times$ to Rubin (two dies on 3 nm). The dies row is therefore the one to watch: after Hopper's 814 mm$^2$, which is within a few percent of the lithographic reticle limit, transistor growth per package came from adding dies, not shrinking them, and the die-to-die link (10 TB/s on Blackwell, slide 17) became part of the architecture. The T4 and L40S columns show the other direction: the graphics and inference dies are a third to nine tenths the size of the compute die of the same generation.

\emph{SMs and lanes.} The SM count rose from 80 to 160 across the two Blackwell dies, 2$\times$, while the lanes per SM doubled once (64 on Volta and Ampere to 128 from Hopper on), so FP32 lanes rose 4$\times$ from V100 to GB300 (5,120 to 20,480; derived). The Rubin entries are not published, and the Blackwell counts come from the microbenchmark papers of Jarmusch and Chandrasekaran (2025) and NVIDIA's Blackwell Ultra blog rather than from a whitepaper, which is why the footnote flags them. These rows change least and explain least: the FP32 peak (not shown, because it stopped being the headline) went from 15.7 TFLOP/s on V100 to 80 on GB200, 5$\times$, while the tensor rows went up 280$\times$.

\emph{The precision rows.} Reading down the FP16 column gives the matrix-unit growth at constant format: 125, 312, 989, 2,500, 2,500, 4,000 TFLOP/s dense, so 2.5$\times$ per generation to Hopper, 2.5$\times$ to Blackwell, flat for the Ultra refresh, 1.6$\times$ to Rubin (derived). Reading across a row at a fixed year gives the format ladder: on GB200, FP16 2.5 PFLOP/s, FP8 5, FP4 10, each rung doubling (slide 10). Reading the ``fastest format'' diagonal gives the 280$\times$ of slide 2. The INT8 row is the one with a surprise: it tracks FP8 (same bit width, same rate) until Blackwell Ultra, where it collapses to about 170 TOPS dense (the GB300 NVL72 page lists 24 POPS with sparsity per rack, so $24{,}000/72/2$ per GPU) because the integer datapath was cut along with most of the FP64 (slide 20); Rubin's INT8 figure is not published. The FP64 row is the subject of slide 21 and its numbers are repeated here so that the two trajectories, up for AI formats and up-flat-down-up for FP64, can be compared in one glance.

\emph{Memory and bandwidth.} Capacity per GPU rose from 16--32 GB to 288 GB, 9--18$\times$, and bandwidth from 0.9 to 19.2 TB/s, 21$\times$ (derived). The HBM generation changed at almost every step (HBM2, HBM2e, HBM3, HBM3e, HBM4), and the stack count rose from four (V100) to eight (Blackwell, Rubin); the capacity jump from GB200 to GB300 (186 to 288 GB) came from taller stacks (twelve dies instead of eight) at the same bandwidth, which is why the GB300 bandwidth entry is unchanged. The two GDDR columns (T4, L40S) show the alternative: a third of the HBM bandwidth at a fraction of the cost, compensated on Ada by the 96 MB L2 (slide 16). The L2 row itself goes from 6 to 50 MB across the single-die compute GPUs and is then not published, a reminder that NVIDIA's disclosure thinned after Hopper.

\emph{Links and power.} NVLink per GPU doubled at every full generation (300, 600, 900, 1,800, 3,000 GB/s, with the Rubin figure the September 2026 revision), and power rose from 300 to 1,400 W with Rubin's undisclosed. Dividing the FP16 row by the power row gives the perf-per-watt series of slide 22 (0.42, 0.78, 1.41, 2.08, 1.79 TFLOP/s per watt for the compute GPUs; derived), which rises until Blackwell and dips for the Ultra refresh because the refresh added memory and FP4, not FP16.

Two cautions about the table. First, every number is a vendor datasheet value for a specific product: the GB200 column is half of a GB200 superchip's figures and the HGX B200 board carries a different GPU (180 GB, 37 TFLOP/s FP64, 1,000 W); the Rubin column is the specification page as revised in September 2026 and earlier figures for the same product were higher. Second, ``dense'' is the deck's choice and not NVIDIA's: the headline figures in NVIDIA's own material (``20 PFLOP/s FP4'' for GB200, ``50 PFLOP/s NVFP4'' for Rubin) are the 2:4-sparse values, twice the table's, and a comparison that mixes the two conventions is wrong by a factor of two. The footnote exists so that the table can be quoted without that error.

\begin{formulas}
Growth factors from the table (derived): transistors $336/21.1 = 16$; FP32 lanes $20{,}480/5{,}120 = 4$; dense FP16 $4{,}000/125 = 32$; fastest format $35{,}000/125 = 280$; HBM capacity $288/16 = 18$; HBM bandwidth $19.2/0.9 = 21$; NVLink $3{,}000/300 = 10$; power $1{,}400/300 = 4.7$ (GB300). Per-watt FP16: $2{,}500/1{,}200 = 2.08$ (GB200), $2{,}500/1{,}400 = 1.79$ (GB300). Sparse headline $= 2\times$ dense: GB200 $2 \times 10 = 20$ PFLOP/s FP4; Rubin $2 \times 25 \ne 50$: Rubin's 50 PFLOP/s inference figure is sparse and its 35 PFLOP/s training figure dense, so the two are not a factor of two apart; NVIDIA has not explained the difference.
\end{formulas}

\begin{terms}
\term{reticle limit}{The largest area a lithography stepper exposes at once, about 830 mm$^2$ at 26 by 33 mm; GA100 (826) and GH100 (814) are near it, so later growth came from adding dies.}
\term{die-to-die link}{The interface joining two dies in one package so that they act as one GPU: NV-HBI at 10 TB/s on Blackwell.}
\term{superchip}{A board holding one Grace or Vera CPU and one or two GPUs; per-GPU figures for GB200 and GB300 are half the superchip's.}
\term{HGX}{NVIDIA's eight-GPU baseboard sold to server vendors; the HGX B200 carries a 180 GB, 1,000 W variant of the Blackwell GPU.}
\term{dense peak}{A tensor-core rate with no sparsity assumed; the deck's convention. NVIDIA's headline figures are usually the 2:4-sparse rates, twice as large.}
\term{HBM stack height}{The number of DRAM dies in a stack (8-high on GB200, 12-high on GB300); more dies give more capacity at the same bandwidth.}
\term{n.p.}{Not published: a figure NVIDIA had not disclosed as of November 2026 (Blackwell and Rubin L2 sizes, Rubin SM count and TDP).}
\end{terms}

\begin{intuition}
A table like this is the whole lecture if you know which rows to read against which: transistors against dies (growth moved from the die to the package), FP16 against FP4 (growth moved from the unit to the format), capacity against bandwidth (memory grew taller before it grew faster), FP16 against power (efficiency stalled when the refresh spent its budget elsewhere). Everything else is a footnote about who measured the number.
\end{intuition}

% ---------------------------------------------------------------------
\slide{31}{Summary: what to take from this lecture}{The SM template held; the tensor core, formats, memory and links changed; the rack became the GPU.}

\para{What the slide shows} Two panels. The left, ``Eight things to remember'', lists: the SM of slide 4 is still the SM, four partitions with a scheduler each, while the tensor core inside it grew 8$\times$ per clock and learned six formats; peaks grew by cutting precision, FP16 to FP8 to FP4, with block scaling making 4 bits usable; data movement got its own hardware (\texttt{cp.async}, TMA, DSMEM, tensor memory); memory, not FLOP/s, sets inference speed; FP64 was built up, held, nearly removed and restored, so it must be checked per product; NVLink turned the node into the rack, 8 to 72 GPUs; the CPU is NVIDIA's too (Grace, Vera), joined at memory speed; and announced is not shipped, so check the datasheet and the date. The right panel, ``Questions the next lecture takes up'', asks four: architecture or platform (AMD's MI300X and MI355X match or beat Hopper and Blackwell on datasheet memory and FLOP/s, the MI455X and Helios match the 72-GPU shape, Intel's Ponte Vecchio reached exascale, so where does NVIDIA's position come from); what FP64 costs (AMD kept 78--288 TFLOP/s on its parts while NVIDIA cut to 1.4 and came back to 33); whether the fabrics stay proprietary (NVLink against UALink and Ultra Ethernet, copper against co-packaged optics); and whether software can be caught up (ROCm, oneAPI and SYCL against CUDA's libraries, and what the portability studies measure).

\para{Details} The eight points, expanded with the slides that support them.

First, the SM template. Turing's SM (slide 4), Ampere's (slide 7) and Hopper's (slide 11) are the same drawing with different counts: four partitions, each with a warp scheduler, a 64 KB register file, 16 or 32 FP32 lanes and one or two tensor cores, sharing an L1/shared memory of 96 to 256 KB. Blackwell's SM internals are not published, but the microbenchmarks find four partitions and a 256 KB register file again, with the tensor memory added beside it (slide 18). What changed inside the template is the matrix unit: Volta's tensor core did 64 FP16 FMAs per clock, Ampere's 256, Hopper's 512 (slides 7 and 11), 8$\times$ per clock per tensor core, and it learned TF32, BF16, FP64, FP8, FP6 and FP4 along the way (slide 10). The previous lecture's constant, the SM as the unit of the architecture, held for another nine years.

Second, precision. The 280$\times$ of slide 2 decomposes into about 4$\times$ from the format (16 bits to 4) and the rest from more and faster tensor cores on more transistors. The decisive technical step was block scaling (slide 18): a shared scale per 16 or 32 values keeps the dynamic range that a 4-bit code cannot, and with it FP4 inference and, by late 2025, FP4 pretraining became practical. The formats arrived one per generation (TF32 and BF16 on Ampere, FP8 on Hopper, FP4 and FP6 on Blackwell), each with a library that made it a default or a flag (slide 29).

Third, data movement. Volta fed its tensor cores from registers loaded by ordinary load instructions. Ampere added \texttt{cp.async} to bypass registers on the way to shared memory (slide 8); Hopper added the TMA to copy whole tiles asynchronously, distributed shared memory to exchange tiles between SMs, and \texttt{wgmma} to read operands from shared memory (slide 13); Blackwell added a 256 KB tensor memory per SM so that accumulators leave the register file (slide 18). Each step moved operand delivery off the general-purpose paths, because the matrix units outgrew what the register file and load/store units could feed (slide 15).

Fourth, memory. The ridge point rose from 139 to about 1,800 FLOP per byte (slide 15), so only large dense matrix multiplies at low precision touch the roof; decoding tokens one at a time is bandwidth-bound, and the products that sold on memory alone (H200 with 141 GB at 4.8 TB/s, B300 with 288 GB) show what inference buyers pay for (slides 14 and 20).

Fifth, FP64. Ampere and Hopper put FP64 on the tensor cores (19.5 and 67 TFLOP/s), Blackwell kept 37--40 TFLOP/s on the CUDA cores with no tensor speedup, Blackwell Ultra removed most of the units (1.2--1.4), and Rubin restored 33 (slide 21). No other row of the table moves in both directions, and no other row has so clearly tracked the vendor's view of its customers; a center buying for FP64 must read the specific product's figure.

Sixth, NVLink. From the DGX-1's eight P100s at 160 GB/s each to the Vera Rubin rack's 72 GPUs at 3 TB/s each (slides 19, 25, 28): the scale-up domain is the rack, its aggregate bandwidth rose about 170$\times$, and the switch began doing arithmetic. The HGX eight-GPU board persists for air-cooled data centers, but the design and pricing unit is the rack.

Seventh, the CPU. Grace (72 Neoverse V2 cores, 900 GB/s NVLink-C2C, slide 14) and Vera (88 custom cores, 1.8 TB/s, slide 24) put NVIDIA's own CPU next to the GPU at memory speed, fixing the CPU-to-GPU ratio at one to two and removing the PCIe host link from the data path; the superchip, not the GPU, is the part that ships.

Eighth, the gap between announcement and shipment, which the deck traces in slides 3, 17, 23 and 25: Hopper's figures rose before shipment, Blackwell's fell (192 to 180 GB, 45 to 40 TFLOP/s FP64, PCIe 6 to 5), Rubin's fell twice (HBM4 13 to 22 to 19.2 TB/s, NVLink 3.6 to 3.0 TB/s), one announced product (Rubin CPX) left the roadmap within six months, and a rack (Kyber) was renamed and, by analysts' reading of its renderings, redesigned. The rule the deck adopts is to quote datasheets with dates, to separate dense from sparse, and to separate measured results (MLPerf, TOP500) from presented ones.

The four questions are the agenda of the comparison lecture. On architecture against platform: AMD's MI300X (192 GB HBM3 at 5.3 TB/s, 1,307 TFLOP/s dense FP16) exceeded the H100's memory and matched its matrix peak on the datasheet in 2023, the MI355X (288 GB, 10 PFLOP/s dense FP4) matches B300's memory and two thirds of its FP4 in 2025, the MI455X and the 72-GPU Helios rack match the NVL72 shape in 2026, and Intel's Ponte Vecchio powered Aurora to 1.012 EFLOP/s; yet the June 2026 market estimates still give NVIDIA about three quarters of data-center AI accelerator revenue$^\dagger$. Whether the difference is the software stack (slide 29), the fabric (slide 28), the supply chain, or the platform cadence (slide 3) is the question. On FP64: AMD's MI300X has 81.7 TFLOP/s FP64 vector and 163 on its matrix cores, the MI355X 78.6, and the MI430X announced for Discovery and Alice Recoque 288; NVIDIA's path through 1.4 TFLOP/s on Blackwell Ultra is the opposite bet, and the two lectures together let a student ask which vendor read the HPC market correctly. On fabrics: NVLink is proprietary and now licensed through NVLink Fusion; UALink is a consortium standard that AMD's Helios runs over Ethernet; Ultra Ethernet targets the scale-out tier; and NVIDIA's own Feynman roadmap replaces the copper spine with co-packaged optics. On software: ROCm's release history, oneAPI and SYCL, and the measured portability studies are the material for answering whether CUDA's lead is permanent.

\begin{terms}
\term{SM template}{The arrangement of four scheduler-partitions sharing an L1/shared memory that every NVIDIA SM since Volta has followed, with the counts inside changing.}
\term{block scaling}{Storing one shared scale factor per small block of low-precision values (16 for NVFP4, 32 for MXFP4) to recover dynamic range; the enabler of 4-bit formats.}
\term{operand delivery}{The path by which matrix-unit inputs reach the unit: registers on Volta, shared memory via TMA and wgmma on Hopper, shared memory and tensor memory on Blackwell.}
\term{design unit}{The level at which NVIDIA defines, sells and specifies a product: the card (through Volta), the HGX board (Ampere, Hopper), the rack (Blackwell, Rubin).}
\term{announced versus shipped}{The deck's rule that a specification is provisional until a datasheet for a shipping product states it, because announced peaks moved 10 to 25 percent in 2024--26.}
\term{Helios}{AMD's 72-GPU MI455X rack of 2026, the first non-NVIDIA rack-scale system in the NVL72 shape; examined in the next lecture.}
\term{Ponte Vecchio}{Intel's Data Center GPU Max 1550, the accelerator in Aurora; examined in the next lecture.}
\term{Ultra Ethernet}{The consortium specification (1.0, June 2025) for an AI- and HPC-tuned Ethernet transport for the scale-out tier.}
\term{co-packaged optics}{Optical transceivers integrated in the switch or GPU package in place of pluggable modules and copper; announced for NVIDIA's Spectrum-6 switches and the Feynman-generation racks.}
\end{terms}

\begin{intuition}
If the eight points compress further, they compress to one: NVIDIA kept the SM and changed everything that feeds it, from the number format to the rack. The next lecture asks whether that is an engineering advantage or a platform advantage, which is another way of asking whether a competitor who builds the same rack from different silicon gets the same result.
\end{intuition}

% ---------------------------------------------------------------------
\slide{32}{Sources}{Whitepapers and datasheets, papers and talks, systems and reporting; $^\dagger$ marks vendor material.}

\para{What the slide shows} Two columns of citations in small type under four headings. ``Architecture whitepapers and datasheets'' (all vendor material): the Turing, Ampere A100, Hopper H100 and Ada whitepapers; the T4, H100, GH200 and L40S datasheets; the Blackwell technical brief and the March 2024 customer deck; NVIDIA's ``Inside Blackwell Ultra'' blog; the GB200 NVL72, GB300 NVL72 and Vera Rubin NVL72 product pages; Lenovo's HGX B200 product guide; and the CUDA 12.0 and 13.0 release notes. ``Papers and talks'': the Hot Chips presentations on A100 (2020), Blackwell (2024), Blackwell Ultra (2025) and Rubin and Vera (2026); Choquette's IEEE Micro articles on A100 and H100; the microbenchmark papers on Turing (Jia et al.\ 2019), Ampere (Abdelkhalik et al.\ 2022), Hopper (Luo et al.\ 2024) and Blackwell (Jarmusch et al.\ 2025); the format papers (Micikevicius et al.\ on FP8, Rouhani et al.\ and the OCP specification on microscaling, NVIDIA's NVFP4 blog and Abecassis et al.\ on NVFP4 pretraining); FlashAttention-3; and Dynamo. ``Systems, benchmarks, status'': the TOP500 lists from November 2020 to June 2026 and the June 2026 Green500; MLPerf Training and Inference rounds 5.0, 5.1 and 6.0; the J\"ulich, Berkeley Lab, DOE, CSCS and Bristol system announcements; NVIDIA's press releases from Blackwell (March 2024) through the Q2 FY2027 results (August 2026); and CoreWeave's Vera Rubin announcement. ``Reporting and analysis'': Lockwood's GTC 2026 recap and processor pages; ServeTheHome, Tom's Hardware and StorageReview on CES, GTC and Hot Chips 2026; The Register on Blackwell's FP64 and mask change; Hardware Busters, Wccftech and TechPowerUp on the September 2026 Vera Rubin revision; SemiAnalysis on achieved throughput; TrendForce on the 2026 product mix. The slide ends by pointing to the paper kit for the full list with links.

\para{Details} The four headings are a hierarchy of evidence, and the deck's dagger convention follows it. Architecture whitepapers describe what was built: SM counts, units per partition, cache sizes, the organization of the die. For Turing, Ampere, Hopper and Ada these exist and the notes cite them without a dagger when they give a structural fact; for Blackwell NVIDIA published a technical brief and blog posts rather than a whitepaper, and for Rubin only keynote slides, product pages and a Hot Chips talk, which is why the Blackwell and Rubin rows of the table on slide 30 carry more ``not published'' entries and more daggers than the earlier ones. Datasheets describe what shipped: peak rates at the shipping clock, memory, power, form factor; they carry daggers because they are product documents, and they are the authority for any number that the whitepaper gave provisionally (the H100's 30 TFLOP/s FP64 in the GTC 2022 whitepaper became 34 on the datasheet). Peer-reviewed and conference material (IEEE Micro, Hot Chips, arXiv microbenchmarks) describes what was measured or explained by the designers and by independent groups; the microbenchmark papers in particular supply the latencies, achieved rates and undisclosed structures (Blackwell's 148 SMs and tensor memory) that no vendor document states. Press releases, product pages and reporting describe status: what shipped when, to whom, and at what revised specification; they carry daggers and dates.

The reporting category is larger in this deck than in the previous one because the deck's last third describes products of 2025--26 whose documentation is incomplete. Three items deserve explicit sourcing notes. The September 2026 revision of the Vera Rubin specifications (slide 25) was first noticed by Hardware Busters and TechPowerUp comparing archived copies of NVIDIA's product page, and NVIDIA's statement of October 29, 2026, is quoted from Wccftech's report; the deck uses the revised page as the authority and the reporting for the dates. The Rubin transistor count of 336 billion appears on CES 2026 keynote slides reported by ServeTheHome and not in the press release, so the deck marks it as secondary. The Blackwell SM count (148 enabled of 160) comes from Jarmusch and Chandrasekaran's microbenchmarks and NVIDIA's Blackwell Ultra blog, not from a whitepaper; the Blackwell Ultra column's 160 is the full die, since NVIDIA has not published the per-SKU count; the L40S's 142 SMs and 96 MB L2 are derived from its lane count and from secondary sources, not from its datasheet. Several figures the deck gives as ``about'' or ``reported'' rest on single secondary sources and are flagged where they appear: Rubin's TDP (1.8--2.3 kW reported), the Selene, Perlmutter and Leonardo system figures, the Groq licensing terms, and the Rubin CPX revival.

Appendix~D of these notes is the reading list, organized by section of the deck with a one-line reason for each entry and an identifier or link; it matches the paper kit distributed with the notes, which includes a fetch script for the open-access items and a BibTeX file. Where an entry is behind a paywall (the IEEE Micro articles, the Hot Chips proceedings through IEEE Xplore), the kit notes it and names an open mirror where one exists (the Hot Chips slide decks are posted by the conference). Students checking a number in these notes should start with the entry in Appendix~D for the slide's section, then the deck's own source line.

\begin{intuition}
Four kinds of document, four kinds of trust: a whitepaper tells you how the machine is built, a datasheet what was sold, a paper what was measured, a press release what was claimed on a date. The dagger marks the third and fourth kinds; the reading list tells you which of the four stands behind each slide.
\end{intuition}
% =====================================================================
\clearpage
\appendix
\setstage{ink}
\section{Notation}

\begin{center}\footnotesize\setlength{\tabcolsep}{4pt}
\begin{tabular}{@{}lp{4.6cm}p{6.3cm}@{}}
\toprule
symbol or term & meaning & notes\\
\midrule
SM & streaming multiprocessor & four scheduler-partitions sharing an L1/shared memory; 80 on V100, 108 on A100, 132 on H100, 148 on GB200, 160 on GB300\\
lane, CUDA core & one FP32 arithmetic pipeline & 64 per SM on Volta, Turing and Ampere; 128 from Hopper on\\
tensor core & the matrix multiply-accumulate unit & 8 per SM on Volta and Turing, 4 per SM from Ampere on; dense FP16 FMA per clock per SM: 512 (Volta), 1{,}024 (Ampere), 2{,}048 (Hopper)\\
$N_{\mathrm{SM}}$, $L$, $f$ & SMs, lanes per SM, clock & FP32 peak $= 2\,N_{\mathrm{SM}}\,L\,f$ (the 2 counts a fused multiply-add as two operations)\\
$M$ & tensor FMAs per clock per SM & tensor peak $= 2\,N_{\mathrm{SM}}\,M\,f$; A100: $2\times 108\times 1{,}024\times 1.41$ GHz $= 312$ TFLOP/s\\
dense, sparse & tensor peaks without and with 2:4 structured sparsity & sparse $= 2\times$ dense; NVIDIA's headline figures are usually sparse\\
$B$ & memory bandwidth, bytes per second & 0.9 TB/s on V100, 3.35 on H100, 8.0 on GB200 and GB300, 19.2 on Rubin\\
$I$, $I^\ast$ & arithmetic intensity; ridge point & $I^\ast = P_{\text{peak}}/B$: 139 FLOP/B on V100 (FP16), 591 on H100 (FP8), 1{,}823 on Rubin (FP4)\\
FP64 to FP16 & floating-point formats, 64 to 16 bits (FP64, FP32, TF32, BF16, FP16) & sign, exponent, mantissa bits: 1/11/52, 1/8/23, 1/8/10, 1/8/7, 1/5/10\\
FP8 E4M3, E5M2 & the two 8-bit formats & 1/4/3 (max 448) and 1/5/2 (max 57{,}344); Hopper on\\
FP6, FP4 E2M1 & the 6- and 4-bit formats & FP4: 1/2/1, sixteen codes; Blackwell on, with block scaling\\
MXFP4, NVFP4 & block-scaled 4-bit formats & MX: 32 values per E8M0 scale; NVFP4: 16 values per E4M3 scale plus one FP32 scale per tensor\\
HBM2 to HBM4 & stacked high-bandwidth memory generations (HBM2, HBM2e, HBM3, HBM3e, HBM4) & per-stack bandwidth about 0.23, 0.41, 0.67, 1.0, 2.4 TB/s (derived from the products)\\
NVLink 2 to 6 & GPU-to-GPU link generations & per GPU 300, 600, 900, 1{,}800, 3{,}000 GB/s (V100, A100, H100, Blackwell, Rubin)\\
NVSwitch & NVLink crossbar chip & generations 2 (Ampere), 3 (Hopper, SHARP), NVLink Switch (Blackwell, 7.2 TB/s), NVLink 6 switch (Rubin)\\
NVL72 & a rack of 72 GPUs and 36 CPUs on one NVLink domain & GB200 (2024), GB300 (2025), Vera Rubin (2026)\\
TDP & thermal design power per GPU & 300 W (V100) to 1{,}400 W (GB300); Rubin not published\\
TMEM, TMA, DSMEM & tensor memory; tensor memory accelerator; distributed shared memory & Blackwell; Hopper; Hopper\\
n.p. & not published & a figure NVIDIA had not disclosed as of November 2026\\
$^\dagger$ & secondary source & datasheet, product page, press release, benchmark list, press\\
\bottomrule
\end{tabular}
\end{center}

\clearpage
\section{Formulas at a glance}

\begin{description}[leftmargin=0pt,labelindent=0pt,itemsep=6pt,font=\normalfont\bfseries\sffamily]\sloppy
\item[FP32 and FP64 peak.] $P = 2\,N_{\mathrm{SM}}\,L\,f$. A100: $2\times 108\times 64\times 1.41\,\mathrm{GHz} = 19.5$ TFLOP/s FP32; with 32 FP64 lanes, 9.7. H100 SXM at the clock implied by its datasheet ($\approx 1.83$ GHz): $2\times 132\times 128\times 1.83 = 61.8$, against 67 quoted (the datasheet uses a slightly higher boost clock); FP64 with 64 lanes per SM, half of that. GB200 FP32: $2\times 148\times 128\times f = 80$ TFLOP/s implies $f \approx 2.1$ GHz (derived from the microbenchmark SM count).
\item[Tensor peak.] $P_{\text{tensor}} = 2\,N_{\mathrm{SM}}\,M\,f$ with $M$ the dense FMAs per clock per SM at the format. A100 FP16: $2\times 108\times 1{,}024\times 1.41 = 312$ TFLOP/s; TF32 at half the FMA rate, 156; FP64 tensor at $M = 64$, 19.5. H100 FP16: $M = 2{,}048$, $2\times 132\times 2{,}048\times 1.83 = 989$; FP8 at $M = 4{,}096$, 1{,}979. Each halving of the operand width doubles $M$ (slide 10). Sparse $= 2\times$ dense.
\item[Format arithmetic.] A binary floating-point format with $e$ exponent bits and $m$ mantissa bits has a largest normal value of about $2^{2^{e-1}}(2 - 2^{-m})$ and a relative precision of $2^{-m}$. E4M3: largest value 448 (NVIDIA's convention reserves the top exponent for NaN), precision $1/8$; E5M2: 57{,}344 and $1/4$; E2M1: values $\pm\{0, 0.5, 1, 1.5, 2, 3, 4, 6\}$. Block scaling stores a block of $n$ values as $n$ codes plus a scale $s$; the stored value is $s\times\text{code}$, and the storage is $4 + 8/n$ bits per value: 4.5 for NVFP4 ($n = 16$, E4M3 scale), 4.25 for MXFP4 ($n = 32$, E8M0 scale).
\item[Bytes per FLOP and the ridge.] $B/P$ and $I^\ast = P/B$ at the fastest dense format: V100 $125/0.9 = 139$ FLOP/B; A100 $312/1.555 = 200$; H100 $1{,}979/3.35 = 591$; GB200 $10{,}000/8 = 1{,}250$; GB300 $15{,}000/8 = 1{,}875$; Rubin $35{,}000/19.2 = 1{,}823$ (all derived). A kernel with intensity below $I^\ast$ is bandwidth-bound; token-by-token decoding of a model with $W$ bytes of weights at batch $b$ has intensity of about $2b/\text{(bytes per weight)}$ FLOP/B, so at $b = 1$ in FP8 it is 2 FLOP/B, three orders of magnitude below the ridge.
\item[Decode time from bandwidth.] One decoding step reads the weights once: $t \approx W/B$. A 70-billion-parameter model in FP8 ($W = 70$ GB) on one H100 ($B = 3.35$ TB/s): about 21 ms per token, 48 tokens per second at batch 1; on one GB300 (8 TB/s) about 9 ms; in NVFP4 (35 GB) on GB300, about 4.4 ms (derived, ignoring the KV cache and overheads). This is why memory bandwidth and capacity, not FLOP/s, sell inference GPUs (slide 15).
\item[Rack totals.] A rack figure is 72 times the per-GPU figure: GB200 NVL72 FP4 dense $72\times 10 = 720$ PFLOP/s (1{,}440 sparse); FP64 $72\times 40 = 2{,}880$ TFLOP/s; HBM $72\times 186\,\mathrm{GB} = 13.4$ TB; HBM bandwidth $72\times 8 = 576$ TB/s; NVLink $72\times 1.8 = 130$ TB/s. Vera Rubin NVL72: $72\times 35 = 2{,}520$ PFLOP/s dense FP4; $72\times 19.2 = 1{,}382 \approx 1{,}400$ TB/s; $72\times 3.0 = 216$ TB/s NVLink; $72\times 33 = 2{,}376 \approx 2{,}400$ TFLOP/s FP64.
\item[Scale-up against scale-out.] NVLink per GPU $= N_{\text{links}}\times b_{\text{link}}$ (bidirectional): Blackwell $18\times 100 = 1{,}800$ GB/s. NIC per GPU in bytes: ConnectX-9 $1.6\ \mathrm{Tb/s}/8 = 200$ GB/s. Ratio on Vera Rubin $3{,}000/200 = 15$. Domain aggregate $= N_{\text{GPU}}\times b_{\text{GPU}}$; growth since the DGX-1 $(72\times 3{,}000)/(8\times 160) \approx 170\times$.
\item[Ring all-reduce.] For $n$ GPUs and $M$ bytes, each GPU sends and receives $2(n-1)M/n$ bytes, so $t \approx 2(n-1)M/(n\,b)$ at per-GPU bandwidth $b$ (ignoring latency). Gradients of a 70B-parameter model in BF16 ($M = 140$ GB) across 72 GPUs at 3.0 TB/s: $2\times 71/72\times 140/3{,}000 \approx 92$ ms; in-switch reduction (SHARP) roughly halves the traffic.
\item[Performance per watt.] $P/\text{TDP}$ at a fixed format: dense FP16 TFLOP/s per watt 0.42 (V100), 0.78 (A100), 1.41 (H100), 2.25 (B200), 2.08 (GB200), 1.79 (GB300). At the fastest format: V100 FP16 0.42 against GB300 FP4 $15{,}000/1{,}400 = 10.7$, a factor of 25 of which about 4 is the format change (derived).
\item[Cluster power.] $N_{\text{GPU}}\times(\text{TDP} + \text{share of CPU, switches, NICs})\times\text{PUE}$: 100{,}000 GPUs at 1.2 kW plus about 25\% for the rest of the rack and a PUE of 1.2--1.3 gives roughly 150--200 MW (derived). Rack power: NVL72 $72\times 1.2\,\mathrm{kW}$ for the GPUs alone is 86 kW; with CPUs, switches, NICs and fans about 120 kW.
\item[Announced versus shipped.] Percentage change $= (\text{shipped} - \text{announced})/\text{announced}$: Blackwell FP64 $(40 - 45)/45 = -11\%$, memory $(180 - 192)/192 = -6\%$; Rubin HBM4 $(19.2 - 22)/22 = -13\%$, NVLink 6 $(3.0 - 3.6)/3.6 = -17\%$; Hopper FP64 $(34 - 30)/30 = +13\%$.
\item[Emulated FP64 (Ozaki scheme).] Split each FP64 mantissa into $k$ integer slices, form the $k^2$ (or fewer, by dropping low-order products) INT8 matrix products on the tensor cores with exact integer accumulation, and recombine with the exponents; cost scales with the number of slice products kept, and accuracy with the slices kept, so a DGEMM-equivalent result costs several INT8 GEMMs. It requires INT8 tensor cores, which Blackwell Ultra largely removed.
\end{description}

\clearpage
\section{Exercises}

\begin{enumerate}[leftmargin=1.6em,itemsep=6pt]
\item \textbf{Peak rates.} From the table on slide 30, recompute the dense FP16 tensor peak of the A100 and H100 from $2\,N_{\mathrm{SM}}\,M\,f$ using $M = 1{,}024$ and $2{,}048$ FMAs per clock per SM and the boost clocks 1.41 and 1.83 GHz. Which quoted figure is reproduced exactly, and what clock would make the other exact? \emph{Answer: $2\times 108\times 1{,}024\times 1.41 = 312$ TFLOP/s exactly; $2\times 132\times 2{,}048\times 1.83 = 989$ TFLOP/s exactly as well; the H100's FP32 figure of 67 would need 1.98 GHz with 128 lanes, which is the SXM5's maximum boost clock, so the two peaks on the datasheet are not computed at the same clock.}
\item \textbf{Dense and sparse.} An NVIDIA product page says a GPU delivers ``20 PFLOPS FP4''. A competitor's datasheet lists ``10 PFLOPS FP4''. Can you conclude which is faster? What two facts do you need? \emph{Answer: no; you need to know whether each figure is dense or 2:4-sparse (NVIDIA's headline is sparse, so GB200 is 10 PFLOP/s dense) and whether both use the same block-scaled format. If the competitor's figure is dense, the two are equal.}
\item \textbf{Format ladder.} On a GB200, dense FP16 is 2.5 PFLOP/s, FP8 5 and FP4 10. A training run uses BF16 for 70\% of its FLOPs and FP8 for 30\%. What is the effective peak? Repeat for 30\% BF16 and 70\% FP8. \emph{Answer: time per FLOP is the weighted mean of $1/P$: $0.7/2.5 + 0.3/5 = 0.34$ ns per PFLOP, effective 2.9 PFLOP/s; $0.3/2.5 + 0.7/5 = 0.26$, effective 3.8 PFLOP/s. Amdahl's law applies: the slow format dominates until it is nearly eliminated.}
\item \textbf{E4M3 range.} The largest E4M3 value is 448 and the smallest normal is $2^{-6} = 0.0156$ (subnormals reach $2^{-9}$). A layer's activations have a maximum magnitude of $3.2\times 10^{3}$ and a typical magnitude of 0.4. Choose the per-tensor scale the Transformer Engine would apply before the cast and state what happens to a value of $10^{-3}$. \emph{Answer: scale $s = 448/3{,}200 \approx 0.14$ so the maximum maps to 448; the typical value maps to 0.056, representable with 3-bit relative precision (about 12\% steps); $10^{-3}$ maps to $1.4\times 10^{-4}$, below the subnormal floor of $2\times 10^{-3}$, and flushes to zero. This is why outliers in activations hurt FP8 and why block scaling (per 16 or 32 values) helps.}
\item \textbf{NVFP4 storage.} A 70-billion-parameter model is stored in NVFP4 (16-value blocks with an E4M3 scale and one FP32 scale per tensor). How many bytes does it occupy, and how does that compare with FP8 and BF16? \emph{Answer: $70\times 10^9\times 4.5/8 = 39.4$ GB (plus a negligible FP32 scale per tensor), against 70 GB in FP8 and 140 GB in BF16. It fits in one GB200's 186 GB in any of the three; in one H100's 80 GB only in FP8 or FP4.}
\item \textbf{Ridge points.} Compute $I^\ast$ for the V100 (FP16), H100 (FP8), GB300 (FP4) and Rubin (FP4) from the table on slide 30, and explain why Rubin's ridge is lower than GB300's despite 2.3$\times$ the FLOP/s. \emph{Answer: 139, 591, 1{,}875 and 1{,}823 FLOP/B; Rubin's HBM4 bandwidth rose 2.4$\times$ (8 to 19.2 TB/s), slightly more than its dense FP4 peak (15 to 35 PFLOP/s, 2.33$\times$), so the ridge fell slightly for the first time in the series.}
\item \textbf{Decoding speed.} A 405-billion-parameter model is served in FP8 on a single GB200 NVL72 rack with the weights split evenly across 72 GPUs. Ignoring the KV cache and communication, estimate the minimum time per output token at batch 1 and the resulting tokens per second. What changes in NVFP4? \emph{Answer: each GPU reads $405/72 = 5.6$ GB per token at 8 TB/s, 0.7 ms, about 1{,}400 tokens per second if the all-reduce were free; in NVFP4 the weights halve to 2.8 GB per GPU and the floor becomes 0.35 ms. In practice communication between the 72 GPUs, not bandwidth, sets the floor at this scale, which is why NVLink bandwidth and in-switch reductions matter.}
\item \textbf{The 72-GPU rack.} Compute the aggregate NVLink bandwidth, HBM capacity, HBM bandwidth and dense FP4 peak of a GB300 NVL72 rack from the per-GPU figures (1.8 TB/s, 288 GB, 8 TB/s, 15 PFLOP/s), and compare with NVIDIA's stated rack figures. \emph{Answer: 130 TB/s, 20.7 TB (NVIDIA says 20 TB), 576 TB/s, 1{,}080 PFLOP/s (NVIDIA: 1.1 EFLOP/s dense, 1.44 sparse); the figures agree to rounding.}
\item \textbf{Announced against shipped.} Vera Rubin's HBM4 bandwidth was announced at 13 TB/s (March 2025), raised to 22 (January 2026) and revised to 19.2 (September 2026). Compute the percentage change at each step and the net change from the first announcement. Do the same for NVLink 6 (3.6 to 3.0 TB/s). \emph{Answer: $+69\%$, $-13\%$, net $+48\%$; NVLink $-17\%$. The compute peaks did not change, so the ridge point moved with the bandwidth: $35{,}000/22 = 1{,}591$ at the January figure against 1{,}823 after the revision.}
\item \textbf{FP64 across generations.} Using 7.8 (V100), 19.5 (A100 tensor), 67 (H100 tensor), 40 (GB200), 1.4 (GB300) and 33 (Rubin) TFLOP/s FP64, compute each rack-of-72 FP64 peak and the number of GB300 racks needed to equal one GB200 rack in FP64. What HPL result would you expect from a 100-rack GB300 installation at 70\% efficiency? \emph{Answer: per rack 562, 1{,}404, 4{,}824, 2{,}880, 101 and 2{,}376 TFLOP/s; about 29 GB300 racks per GB200 rack; 100 GB300 racks give 10 PFLOP/s peak and about 7 PFLOP/s on HPL, well down the June 2026 list for an installation of 7{,}200 GPUs.}
\item \textbf{Emulation.} An Ozaki-scheme DGEMM uses three INT8 slices per operand and keeps six of the nine slice products. If a GB200's INT8 tensor rate is 5 POPS dense, what is the implied FP64-equivalent rate, ignoring the splitting and recombination overhead, and how does it compare with the 40 TFLOP/s native figure? Why does the scheme not help on a GB300? \emph{Answer: $5{,}000/6 \approx 830$ TFLOP/s before overheads, more than 20 times native; the overheads (splitting, exponent alignment, recombination, memory traffic for the slices) are substantial, and NVIDIA's measured figure for cuBLAS's emulated DGEMM is up to 2.3$\times$ native on a GB200 NVL72 (CUDA 13.0 Update 2, October 2025). GB300 cut INT8 to about 0.17 POPS dense (24 POPS sparse per 72-GPU rack, so $24{,}000/72/2$), which at six slice products gives at most $170/6 \approx 28$ TFLOP/s before overheads: more than its 1.4 TFLOP/s native but far below a GB200, and in practice the recombination overhead eats most of it.}
\item \textbf{Performance per watt.} Compute dense FP16 TFLOP/s per watt for V100 (125 TF, 300 W), H100 (989, 700), GB200 (2,500, 1,200) and GB300 (2,500, 1,400), and dense FP4 per watt for GB200 and GB300. How much of the GB300's advantage over the V100 at the fastest format comes from the format change alone? \emph{Answer: 0.42, 1.41, 2.08, 1.79 TFLOP/s per watt at FP16; 8.3 and 10.7 at FP4; the ratio $10.7/0.42 = 25$, of which the format change (FP16 to FP4) accounts for 4$\times$ at equal silicon and the remaining 6$\times$ is architecture and process. At a fixed format the gain is 4$\times$ in nine years.}
\item \textbf{Cluster power.} A training cluster of 100{,}000 GB200 GPUs runs at 1.2 kW per GPU, with CPUs, switches and NICs adding 25\% and a power usage effectiveness of 1.25. Estimate the facility power, the energy of a 30-day training run in GWh, and the number of NVL72 racks. \emph{Answer: $100{,}000\times 1.2\times 1.25\times 1.25 = 188$ MW; $188\,\mathrm{MW}\times 720\,\mathrm{h} = 135$ GWh; $100{,}000/72 = 1{,}389$ racks. Stargate Abilene's 1.2 GW corresponds to roughly six such clusters.}
\item \textbf{Scale-up against scale-out.} An all-to-all exchange for expert parallelism moves 50 MB per GPU per layer among 72 GPUs. Estimate the time per layer if the exchange runs over NVLink 6 (3 TB/s per GPU, bidirectional) and if it ran over ConnectX-9 NICs (1.6 Tb/s). With 60 layers and a target of 10 ms per token, which fits? \emph{Answer: 50 MB at 1.5 TB/s per direction is about 33 $\mu$s per layer, 2 ms per token on NVLink; at 100 GB/s per direction it is 0.5 ms per layer, 30 ms per token on the NIC, three times the budget. Expert parallelism must stay inside the NVLink domain, which is the argument for the 72-GPU rack.}
\item \textbf{Reading the TOP500.} From the June 2026 top ten on slide 27, compute HPL GFLOP/s per watt for El Capitan (1,809 PF, 29.7 MW), Frontier (1,353 PF, 24.6 MW), Aurora (1,012 PF, 38.7 MW) and JUPITER (1,000 PF, 15.8 MW). Which accelerator generation is most efficient on HPL, and why does this not settle the GPU comparison of the next lecture? \emph{Answer: 61, 55, 26 and 63 GFLOP/s per watt; the GH200 (JUPITER) edges the MI300A (El Capitan). HPL is dense FP64 at full occupancy, the one workload closest to the roof; efficiency on bandwidth-bound or low-precision workloads can order the machines differently, and none of the Blackwell-era racks is on the list at all.}
\item \textbf{Reading a spec sheet in 2026.} Take the Vera Rubin NVL72 specification page and extract, per GPU: dense FP4 training peak, FP8 peak, FP64, HBM capacity and bandwidth, NVLink bandwidth, and the NIC rate. Compute the ridge point at FP4 and FP8, the bytes per FLOP, and the rack totals, then list which figures changed between the January and September 2026 versions of the page and by how much. \emph{Answer (from the deck): 35 PFLOP/s, 17.5, 33 TFLOP/s, 288 GB, 19.2 TB/s, 3.0 TB/s, 1.6 Tb/s; $I^\ast = 1{,}823$ (FP4) and 911 (FP8); $0.00055$ bytes per FLOP at FP4; rack 2{,}520 PFLOP/s, 1{,}400 TB/s, 216 TB/s; HBM4 22 to 19.2 TB/s ($-13\%$), NVLink 3.6 to 3.0 ($-17\%$), rack HBM 1{,}580 to 1{,}400 TB/s, rack NVLink 260 to 216 TB/s; compute peaks unchanged.}
\end{enumerate}
%=====
\clearpage
\setstage{ink}
\section{Reading list}
\label{app:reading}
The papers, whitepapers, datasheets and reports behind each part of the deck, in the order they are discussed, with one line on why each is worth reading. Entries with an arXiv identifier are included in the paper archive of the kit; entries without one are listed with a URL. Items marked \emph{secondary source} are press coverage, analyst notes or commentary; vendor material is named as such.

\subsection*{1. Overview: nine years after Volta (slides 1-3)}
\begin{itemize}[leftmargin=1.2em,itemsep=3pt]\sloppy
\item \textbf{Dally, Keckler and Kirk (2021).} \textit{Evolution of the graphics processing unit (GPU)}. \url{https://doi.org/10.1109/MM.2021.3113475}. IEEE Micro 41(6). The architects' retrospective from G80 to Ampere; the starting point for everything after Volta.
\item \textbf{Mann (The Register) (2024).} \textit{Chipmakers' Computex roadmaps: NVIDIA's yearly cadence (June 2024)}. \url{https://www.theregister.com/2024/06/05/chipmakers_computex_roadmaps}. Press, secondary source. The June 2, 2024 keynote in which NVIDIA laid out the annual cadence (Blackwell Ultra 2025, Rubin 2026, Rubin Ultra 2027), first signaled to investors in October 2023. Slide 3.
\item \textbf{Lockwood (2026).} \textit{GTC 2026 recap (blog, March 2026)}. \url{https://blog.glennklockwood.com/2026/03/gtc-2026-recap.html}. Independent analyst notes, open. The clearest account of what GTC 2026 changed: Rubin CPX gone, Kyber redesigned, Feynman and the Rosa CPU, co-packaged optics. Slides 3 and 26.
\item \textbf{Lockwood (2025).} \textit{B200 (processor garden page)}. \url{https://glennklockwood.com/garden/processors/B200}. Independent, open. A sourced table of every Blackwell variant's announced and shipping specifications; the reference for the announced-versus-shipped story.
\item \textbf{Lockwood (2026).} \textit{B300 (processor garden page)}. \url{https://glennklockwood.com/garden/processors/B300}. Independent, open. GB300 versus HGX B300 specifications, including the FP64 cut.
\item \textbf{Lockwood (2026).} \textit{Rubin R200 (processor garden page)}. \url{https://glennklockwood.com/garden/processors/R200}. Independent, open. The provenance of every Rubin number, with the NVL144 to NVL72 history.
\item \textbf{Micikevicius et al. (2017).} \textit{Mixed precision training}. arXiv:1710.03740. ICLR 2018. The FP16 training recipe (master weights, loss scaling) that every later format rung builds on; the Volta-era baseline for slide 10.
\end{itemize}
\subsection*{2. Turing (slides 4-5)}
\begin{itemize}[leftmargin=1.2em,itemsep=3pt]\sloppy
\item \textbf{NVIDIA (2018).} \textit{NVIDIA Turing GPU architecture whitepaper}. \url{https://images.nvidia.com/aem-dam/Solutions/design-visualization/technologies/turing-architecture/NVIDIA-Turing-Architecture-Whitepaper.pdf}. Vendor whitepaper. TU102 Table 1 (18.6B transistors, 754 mm2), the SM with RT cores and INT8/INT4 tensor cores, mesh shading, variable-rate shading, DLSS. Slide 4.
\item \textbf{NVIDIA (2018).} \textit{Tesla T4 datasheet}. \url{https://nvidia.com/content/dam/en-zz/Solutions/Data-Center/tesla-t4/t4-tensor-core-datasheet-951643.pdf}. Vendor datasheet. 2,560 CUDA cores, 320 tensor cores, 65 TFLOP/s FP16, 130 TOPS INT8, 260 TOPS INT4, 16 GB GDDR6 at 300 GB/s, 70 W. Slide 5.
\item \textbf{Jia, Maggioni, Smith and Scarpazza (2019).} \textit{Dissecting the NVidia Turing T4 GPU via microbenchmarking}. arXiv:1903.07486. The Volta microbenchmark methodology applied to the T4: cache latencies, the integer tensor modes, clock behavior under the 70 W limit. Slide 5.
\item \textbf{Jia, Maggioni, Staiger and Scarpazza (2018).} \textit{Dissecting the NVIDIA Volta GPU architecture via microbenchmarking}. arXiv:1804.06826. The baseline measurements (shared 19 cycles, L1 28, L2 about 193, global 375 to 1,029) that the later papers compare against.
\item \textbf{NVIDIA (2020).} \textit{NVIDIA DLSS 2.0: a big leap in AI rendering (GTC 2020 talk and developer page)}. \url{https://developer.nvidia.com/blog/nvidia-dlss-2-0-a-big-leap-in-ai-rendering/}. Vendor blog. The temporal network that made tensor cores a standard graphics stage. Slide 5.
\end{itemize}
\subsection*{3. Ampere (slides 6-10)}
\begin{itemize}[leftmargin=1.2em,itemsep=3pt]\sloppy
\item \textbf{NVIDIA (2020).} \textit{NVIDIA A100 Tensor Core GPU architecture (Ampere whitepaper)}. \url{https://images.nvidia.com/aem-dam/en-zz/Solutions/data-center/nvidia-ampere-architecture-whitepaper.pdf}. Vendor whitepaper. GA100 Table 4 (54.2B transistors, 826 mm2, 108 of 128 SMs), the third-generation tensor core, TF32, BF16, 2:4 sparsity, MIG, asynchronous copy, NVLink 3, DGX A100 appendix. Slides 6-9.
\item \textbf{Choquette and Gandhi (2020).} \textit{NVIDIA A100 GPU: performance and innovation (Hot Chips 32)}. \url{https://hc32.hotchips.org/assets/program/conference/day1/HotChips2020_GPU_NVIDIA_Choquette_v01.pdf}. Conference slides, open. The design talk: the SM, the tensor core data path, the L2 and the NVLink 3 fabric. Slide 7.
\item \textbf{Choquette, Gandhi, Giroux, Stam and Krashinsky (2021).} \textit{NVIDIA A100 Tensor Core GPU: performance and innovation}. \url{https://doi.org/10.1109/MM.2021.3061394}. IEEE Micro 41(2). The peer-reviewed A100 paper: tensor-core generations, sparsity, MIG, asynchronous data movement.
\item \textbf{Abdelkhalik, Arafa, Santhi and Badawy (2022).} \textit{Demystifying the Nvidia Ampere architecture through microbenchmarking and instruction-level analysis}. arXiv:2208.11174. Measured latencies and throughputs on the A100, including the tensor-core instructions and the asynchronous copy path. Slides 7-8 and 15.
\item \textbf{Tom's Hardware (2020).} \textit{Nvidia RTX 30-series Ampere architecture deep dive}. \url{https://www.tomshardware.com/features/nvidia-ampere-architecture-deep-dive}. Press, secondary source. GA102's dual-FP32 datapath, Samsung 8N, 28.3B transistors on 628 mm2: the graphics-die side of slide 9.
\item \textbf{Mishra, Latorre, Pool, Stosic, Stosic, Venkatesh, Yu and Micikevicius (2021).} \textit{Accelerating sparse deep neural networks}. arXiv:2104.08378. The 2:4 structured-sparsity recipe (prune, fine-tune, 2x on Ampere tensor cores) behind every sparse peak since 2020. Slide 8.
\item \textbf{Kalamkar et al. (Intel) (2019).} \textit{A study of BFLOAT16 for deep learning training}. arXiv:1905.12322. Why 8 exponent bits and 7 mantissa bits suffice for training without loss scaling; the case for the BF16 rung on slide 10.
\item \textbf{NVIDIA (2020).} \textit{NVIDIA Multi-Instance GPU user guide}. \url{https://docs.nvidia.com/datacenter/tesla/mig-user-guide/index.html}. Vendor documentation. The seven-instance partitioning of the A100 (and later GPUs): what an instance owns and how it is configured. Slide 8.
\item \textbf{NVIDIA (2020).} \textit{HGX A100 datasheet}. \url{https://www.nvidia.com/content/dam/en-zz/Solutions/Data-Center/HGX/pdf/nvidia-hgx-a100-datasheet.pdf}. Vendor datasheet. The 8-GPU board: 600 GB/s NVLink 3 per GPU, 4.8 TB/s aggregate through six NVSwitches. Slide 9.
\end{itemize}
\subsection*{4. Hopper (slides 11-15)}
\begin{itemize}[leftmargin=1.2em,itemsep=3pt]\sloppy
\item \textbf{NVIDIA (2022).} \textit{NVIDIA H100 Tensor Core GPU architecture (GTC 2022 whitepaper)}. \url{https://www.ece.lsu.edu/koppel/gp/refs/gtc22-whitepaper-hopper.pdf}. Vendor whitepaper (mirror; original at resources.nvidia.com). GH100 Table 3, thread block clusters, DSMEM, TMA, asynchronous transaction barriers, Transformer Engine, DPX, NVLink 4; note that its peaks are preliminary (FP64 30, later 34). Slides 11-13.
\item \textbf{NVIDIA (2023).} \textit{H100 Tensor Core GPU datasheet}. \url{https://openzeka.com/wp-content/uploads/2023/12/h100-datasheet-1.pdf}. Vendor datasheet (mirror). The shipping peaks with the sparsity footnote: TF32 494/989, FP16 989/1,979, FP8 1,979/3,958 dense/sparse; SXM, PCIe and NVL variants. Slide 12.
\item \textbf{Choquette (2023).} \textit{NVIDIA Hopper H100 GPU: scaling performance}. \url{https://ieeexplore.ieee.org/document/10070122}. IEEE Micro 43(3). The peer-reviewed Hopper paper: the SM, the fourth-generation tensor core, FP8, clusters and the NVLink Switch System.
\item \textbf{Andersch et al. (NVIDIA) (2022).} \textit{NVIDIA Hopper architecture in-depth (developer blog)}. \url{https://developer.nvidia.com/blog/nvidia-hopper-architecture-in-depth/}. Vendor blog. The most readable account of TMA, DSMEM, wgmma and the Transformer Engine. Slide 13.
\item \textbf{Micikevicius, Stosic, Burgess, Cornea, Dubey, Grisenthwaite, Ha, Heinecke, Judd, Kamalu, Mellempudi, Oberman, Shoeybi, Siu and Wu (2022).} \textit{FP8 formats for deep learning}. arXiv:2209.05433. NVIDIA, Arm and Intel's joint proposal of E4M3 and E5M2 and the per-tensor scaling recipe, with training results matching BF16. Slide 12.
\item \textbf{Luo, Fan, Li, Du, Wang and Chu (2024).} \textit{Benchmarking and dissecting the Nvidia Hopper GPU architecture}. arXiv:2402.13499. Microbenchmarks of the H100: DPX, DSMEM latency, FP8 tensor throughput, wgmma, clusters. Slides 13 and 15.
\item \textbf{NVIDIA (2023).} \textit{NVIDIA supercharges Hopper, the world's leading AI computing platform (press release, November 13, 2023)}. \url{https://nvidianews.nvidia.com/news/nvidia-supercharges-hopper-the-worlds-leading-ai-computing-platform}. Vendor press release. The H200: 141 GB of HBM3e at 4.8 TB/s, availability Q2 2024. Slide 14.
\item \textbf{NVIDIA (2023).} \textit{GH200 Grace Hopper Superchip datasheet}. \url{https://www.amax.com/content/files/2023/12/NVIDIA_GH200_Grace_Hopper_Superchip_Datasheet.pdf}. Vendor datasheet (mirror). Grace (72 Neoverse V2 cores, 480 GB LPDDR5X at up to 512 GB/s), NVLink-C2C at 900 GB/s, 96 GB HBM3 or 144 GB HBM3e. Slide 14.
\item \textbf{NVIDIA (2024).} \textit{Grace CPU performance tuning guide}. \url{https://docs.nvidia.com/grace-performance-tuning-guide.pdf}. Vendor documentation. The Grace microarchitecture and memory system for readers who want the CPU side of the superchip.
\item \textbf{NVIDIA (2023).} \textit{AWS and NVIDIA announce strategic collaboration (press release, November 2023)}. \url{https://nvidianews.nvidia.com/news/aws-nvidia-strategic-collaboration-for-generative-ai}. Vendor press release. The GH200 NVL32 rack and Project Ceiba: the first 32-GPU NVLink domains with Grace Hopper.
\item \textbf{Shah, Bikshandi, Zhang, Thakkar, Ramani and Dao (2024).} \textit{FlashAttention-3: fast and accurate attention with asynchrony and low-precision}. arXiv:2407.08691. How TMA, wgmma and FP8 are used in practice: attention at about 75 percent of the H100's FP16 peak. Slides 13 and 29.
\item \textbf{NVIDIA (2022).} \textit{NVIDIA NVLink and NVSwitch (product page, third generation)}. \url{https://www.nvidia.com/en-us/data-center/nvlink/}. Vendor page. NVLink 4 at 900 GB/s, the 64-port NVSwitch with SHARP, and the 256-GPU NVLink Switch System. Slide 14.
\end{itemize}
\subsection*{5. Ada Lovelace (slide 16)}
\begin{itemize}[leftmargin=1.2em,itemsep=3pt]\sloppy
\item \textbf{NVIDIA (2022).} \textit{NVIDIA Ada GPU architecture whitepaper}. \url{https://images.nvidia.com/aem-dam/Solutions/geforce/ada/nvidia-ada-gpu-architecture.pdf}. Vendor whitepaper. AD102 (76.3B transistors, 608.5 mm2, 144 SMs, 96 MB L2), shader execution reordering, third-generation RT cores, DLSS 3, fourth-generation tensor cores with FP8. Slide 16.
\item \textbf{Supermicro and NVIDIA (2023).} \textit{NVIDIA L40S systems datasheet}. \url{https://www.supermicro.com/datasheet/datasheet_NVIDIA_L40S_Systems.pdf}. Vendor datasheet. L40S peaks (91.6 TFLOP/s FP32, 733 dense FP8, 48 GB at 864 GB/s, 350 W). Slide 16.
\end{itemize}
\subsection*{6. Blackwell and Blackwell Ultra (slides 17-22)}
\begin{itemize}[leftmargin=1.2em,itemsep=3pt]\sloppy
\item \textbf{NVIDIA (2024).} \textit{NVIDIA Blackwell platform arrives to power a new era of computing (press release, March 18, 2024)}. \url{https://nvidianews.nvidia.com/news/nvidia-blackwell-platform-arrives-to-power-a-new-era-of-computing}. Vendor press release. 208B transistors, the 10 TB/s NV-HBI, fifth-generation tensor cores, NVLink 5, the GB200 NVL72. Slide 17.
\item \textbf{NVIDIA GmbH (2024).} \textit{Blackwell customer deck (March 2024)}. \url{https://mcomputers.cz/wp-content/uploads/2024/04/NVIDIA-Blackwell-Praha-PDF.pdf}. Vendor slides (mirror). The original B100/B200/GB200 specifications (FP64 30/40/45 TFLOP/s, 192 GB, 8 TB/s, PCIe 6) that later shipped lower; the primary source for slide 17's announced-versus-shipped comparison.
\item \textbf{NVIDIA (2024).} \textit{NVIDIA Blackwell architecture technical brief}. \url{https://resources.nvidia.com/en-us-blackwell-architecture}. Vendor document (registration). The official architecture brief: tensor-core formats, decompression and RAS engines, confidential computing.
\item \textbf{Tirumala and Wong (2024).} \textit{NVIDIA Blackwell platform: advancing generative AI and accelerated computing (Hot Chips 36)}. \url{https://doi.org/10.1109/HCS61935.2024.10665247}. Conference talk, IEEE Xplore. The official Blackwell presentation: the two-die design, NV-HBI, NVLink 5 and the rack.
\item \textbf{NVIDIA (2025).} \textit{GB200 NVL72 (product page)}. \url{https://www.nvidia.com/en-us/data-center/gb200-nvl72/}. Vendor page. The shipping rack specifications with dense and sparse footnotes: 1,440/720 PFLOP/s FP4, 2,880 TFLOP/s FP64, 13.4 TB HBM3e, 130 TB/s NVLink. Slide 19.
\item \textbf{Lenovo (2026).} \textit{ThinkSystem NVIDIA B200 180GB 1000W GPU (product guide)}. \url{https://lenovopress.lenovo.com/lp2226-thinksystem-nvidia-b200-180gb-1000w-gpu}. Vendor documentation. The HGX B200 as shipped: 180 GB, 7.7 TB/s, 37 TFLOP/s FP64, PCIe 5. Slide 17.
\item \textbf{heise online (2024).} \textit{Blackwell: Nvidia unveils its next generation of AI accelerators (page 2)}. \url{https://heise.de/en/news/Blackwell-Nvidia-unveils-its-next-generation-of-AI-accelerators-9782399.html?seite=2}. Press, secondary source. The NVLink Switch chip (7.2 TB/s, 50B transistors) and the B100/B200/GB200 power figures. Slide 19.
\item \textbf{Mann (The Register) (2024).} \textit{Nvidia's future in scientific computing hinges on a melding of AI and HPC (May 2024)}. \url{https://www.theregister.com/2024/05/22/nvidias_future_in_scientific_computing/}. Press, secondary source. The Blackwell FP64 debate as it happened, with NVIDIA's position. Slide 21.
\item \textbf{The Register (2024).} \textit{Nvidia confirms Blackwell mask change (August 2024)}. \url{https://www.theregister.com/2024/08/29/nvidia_blackwell_manufacturing/}. Press, secondary source. The Q2 FY25 call: a mask change to improve yield, no functional change, ramp in Q4 2024. Slide 17.
\item \textbf{Neowin (2024).} \textit{Microsoft Azure is the first hyperscaler with NVIDIA's Blackwell GB200 (October 2024)}. \url{https://www.neowin.net/amp/microsoft-azure-is-the-first-hyperscaler-to-have-nvidias-blackwell-system-with-gb200/}. Press, secondary source. The first GB200 deployment, October 8, 2024. Slide 19.
\item \textbf{Jarmusch, Graddon and Chandrasekaran (2025).} \textit{Dissecting the NVIDIA Blackwell architecture with microbenchmarks}. arXiv:2507.10789. Consumer Blackwell (RTX 5080) against the H100: the first independent measurements of the fifth-generation tensor core.
\item \textbf{Jarmusch and Chandrasekaran (2025).} \textit{Microbenchmarking NVIDIA's Blackwell architecture: tensor memory, tcgen05 and the B200}. arXiv:2512.02189. The B200 measured: 148 SMs, 256 KB tensor memory per SM, tcgen05 issue, 44.8 TFLOP/s FP64. The source for the SM counts the datasheets omit. Slides 15, 17, 18.
\item \textbf{NVIDIA (2025).} \textit{Inside NVIDIA Blackwell Ultra: the chip powering the AI factory era (developer blog, August 2025)}. \url{https://developer.nvidia.com/blog/inside-nvidia-blackwell-ultra-the-chip-powering-the-ai-factory-era/}. Vendor blog. 160 SMs, 288 GB HBM3e, 15/20 PFLOP/s FP4, the doubled SFU exponential rate, PCIe 6. Slide 20.
\item \textbf{Tom's Hardware (2025).} \textit{Nvidia shares Blackwell Ultra's secrets at Hot Chips 2025}. \url{https://www.tomshardware.com/pc-components/gpus/nvidia-shares-blackwell-ultras-secrets-nvfp4-boost-detailed-and-pcie-6-0-support}. Press, secondary source. The Hot Chips 37 table with the B300's 1.39 TFLOP/s FP64. Slides 20-21.
\item \textbf{NVIDIA (2025).} \textit{GB300 NVL72 (product page)}. \url{https://www.nvidia.com/en-au/data-center/gb300-nvl72}. Vendor page. 1,440/1,080 PFLOP/s FP4, 100 TFLOP/s FP64 per rack, 20 TB HBM3e, 37 TB fast memory. Slide 20.
\item \textbf{NVIDIA (2025).} \textit{Blackwell Ultra datasheet}. \url{https://resources.nvidia.com/en-us-blackwell-architecture/blackwell-ultra-datasheet}. Vendor datasheet (registration). GB300 and HGX B300 per-GPU figures.
\item \textbf{Tom's Hardware (2025).} \textit{NVIDIA's newest top-tier AI supercomputers deployed for the first time at CoreWeave (July 2025)}. \url{https://www.tomshardware.com/tech-industry/artificial-intelligence/nvidias-newest-top-tier-ai-supercomputers-deployed-for-the-first-time-grace-blackwell-ultra-superchip-systems-deployed-at-coreweave}. Press, secondary source. The first GB300 NVL72 racks, July 3, 2025. Slide 20.
\item \textbf{NVIDIA (2025).} \textit{Introducing NVFP4 for efficient and accurate low-precision inference (developer blog, June 2025)}. \url{https://developer.nvidia.com/blog/introducing-nvfp4-for-efficient-and-accurate-low-precision-inference/}. Vendor blog. The format: E2M1 values, 16-element blocks with an E4M3 scale and an FP32 per-tensor scale; the accuracy claims. Slide 18.
\item \textbf{Abecassis et al. (NVIDIA) (2025).} \textit{Pretraining large language models with NVFP4}. arXiv:2509.25149. A 12B-parameter model trained on 10 trillion tokens in NVFP4 matching the FP8 loss curve; the recipe (block scaling, stochastic rounding, mixed layers). Slide 18.
\item \textbf{Rouhani et al. (2023).} \textit{Microscaling data formats for deep learning}. arXiv:2310.10537. The OCP MX formats (MXFP8, MXFP6, MXFP4) from Microsoft, AMD, Intel, Meta, NVIDIA and others: 32-element blocks with an E8M0 scale. Slides 10 and 18.
\item \textbf{Open Compute Project (2023).} \textit{OCP microscaling formats (MX) specification version 1.0}. \url{https://www.opencompute.org/documents/ocp-microscaling-formats-mx-v1-0-spec-final-pdf}. Standard text. The normative definition of the MX formats.
\item \textbf{Ozaki, Ogita, Oishi and Rump (2012).} \textit{Error-free transformations of matrix multiplication by using fast routines of matrix multiplication and its applications}. \url{https://doi.org/10.1007/s11075-011-9478-1}. Numerical Algorithms 59. The Ozaki scheme: splitting floating-point matrices into exactly representable slices so that lower-precision (or integer) GEMMs reproduce high-precision results; the basis of FP64 emulation on tensor cores. Slide 21.
\item \textbf{Ootomo, Ozaki and Yokota (2024).} \textit{DGEMM on integer matrix multiplication unit}. \url{https://doi.org/10.1177/10943420241239588}. IJHPCA 38(4); arXiv 2306.11975. FP64 matrix multiply on INT8 tensor cores with the Ozaki scheme, faster than native FP64 on consumer GPUs: the method cuBLAS adopted. Slide 21.
\item \textbf{NVIDIA (2025).} \textit{Unlocking tensor core performance with floating-point emulation in cuBLAS (developer blog, October 2025)}. \url{https://developer.nvidia.com/blog/unlocking-tensor-core-performance-with-floating-point-emulation-in-cublas/}. Vendor blog. FP64 emulation on Blackwell's INT8 tensor cores (Ozaki scheme) shipped in CUDA 13.0 Update 2: how it is enabled and the claim of up to 2.3x native DGEMM on GB200 NVL72. Slides 21 and 29.
\item \textbf{NVIDIA (2025).} \textit{Boosting matrix multiplication speed and flexibility with NVIDIA cuBLAS 12.9 (developer blog, May 2025)}. \url{https://developer.nvidia.com/blog/boosting-matrix-multiplication-speed-and-flexibility-with-nvidia-cublas-12-9/}. Vendor blog. The 12.9 release: FP32 emulation on BF16 tensor cores, block-scaled FP8/FP4 GEMMs, and the note that FP64 emulation was still under development. Slide 29.
\item \textbf{The Register (2024).} \textit{Nvidia's DGX GB200 NVL72 (March 2024)}. \url{https://www.theregister.com/2024/03/21/nvidia_dgx_gb200_nvk72}. Press, secondary source. The rack's 120 kW and 1,200 W per GPU. Slide 22.
\end{itemize}
\subsection*{7. Rubin and Vera (slides 23-26)}
\begin{itemize}[leftmargin=1.2em,itemsep=3pt]\sloppy
\item \textbf{Tom's Hardware (2025).} \textit{Nvidia announces Rubin GPUs in 2026, Rubin Ultra in 2027, Feynman after (GTC, March 2025)}. \url{https://www.tomshardware.com/pc-components/gpus/nvidia-announces-rubin-gpus-in-2026-rubin-ultra-in-2027-feynam-after}. Press, secondary source. The GTC 2025 numbers: Vera Rubin NVL144 (3.6 EF FP4, 13 TB/s HBM4, 260 TB/s), Rubin Ultra NVL576 (15 EF, 600 kW), Feynman. Slides 23, 25, 26.
\item \textbf{NVIDIA (2025).} \textit{NVIDIA unveils Rubin CPX (press release, September 9, 2025)}. \url{https://nvidianews.nvidia.com/news/nvidia-unveils-rubin-cpx-a-new-class-of-gpu-designed-for-massive-context-inference}. Vendor press release. The prefill GPU as announced: 128 GB GDDR7, 30 PFLOP/s NVFP4, the NVL144 CPX rack. Slide 26.
\item \textbf{RCR Tech (2026).} \textit{Nvidia reportedly revives Rubin CPX (September 2026)}. \url{https://rcrtech.com/semiconductor-news/nvidia-reportedly-revives-rubin-cpx/}. Press, secondary source, unconfirmed. The report of a CPX revival with HBM4 for 2027. Slide 26.
\item \textbf{NVIDIA (2026).} \textit{NVIDIA kicks off the next generation of AI with Rubin: six new chips, one incredible AI supercomputer (press release, January 5, 2026)}. \url{https://investor.nvidia.com/news/press-release-details/2026/NVIDIA-Kicks-Off-the-Next-Generation-of-AI-With-Rubin--Six-New-Chips-One-Incredible-AI-Supercomputer/default.aspx}. Vendor press release. The CES announcement: Vera, Rubin, the NVLink 6 switch, ConnectX-9, BlueField-4, Spectrum-6; the return to the NVL72 name. Slides 23-24.
\item \textbf{ServeTheHome (2026).} \textit{NVIDIA launches next-generation Rubin AI compute platform at CES 2026}. \url{https://www.servethehome.com/nvidia-launches-next-generation-rubin-ai-compute-platform-at-ces-2026/}. Press, secondary source. The keynote slides with the 336B transistor count, the 22 TB/s HBM4 figure of the time, and Vera details. Slide 23.
\item \textbf{Tom's Hardware (2026).} \textit{Nvidia's seven-chip Vera Rubin platform turns the data center into an AI factory (GTC, March 2026)}. \url{https://www.tomshardware.com/tech-industry/artificial-intelligence/nvidias-seven-chip-vera-rubin-platforms-turns-the-data-center-into-an-ai-factory}. Press, secondary source. GTC 2026: the platform with the Groq 3 LPX added, pod-level figures. Slide 24.
\item \textbf{StorageReview (2026).} \textit{NVIDIA GTC 2026: Rubin GPUs, Groq LPUs, Vera CPUs}. \url{https://www.storagereview.com/news/nvidia-gtc-2026-rubin-gpus-groq-lpus-vera-cpus-and-what-nvidia-is-building-for-trillion-parameter-inference}. Press, secondary source. The Groq 3 LPX and the Vera rack as presented at GTC 2026. Slides 24 and 26.
\item \textbf{NVIDIA (2026).} \textit{Vera Rubin NVL72 (product page, revised September-October 2026)}. \url{https://www.nvidia.com/en-us/data-center/vera-rubin-nvl72}. Vendor page. The current specifications: 19.2 TB/s HBM4 per GPU, 3 TB/s NVLink 6, 216 TB/s per rack, 2.4 PFLOP/s FP64, 33 TFLOP/s FP64 per GPU. Slides 23 and 25.
\item \textbf{Hardware Busters (2026).} \textit{Nvidia quietly lowers Vera Rubin HBM4 bandwidth and the HGX NVL8 loses a quarter of it (October 2026)}. \url{https://hwbusters.com/news/nvidia-quietly-lowers-vera-rubin-hbm4-bandwidth-and-the-hgx-nvl8-loses-a-quarter-of-it/}. Press, secondary source. The old-versus-new table of the September 2026 revision, dated with the Wayback Machine. Slide 25.
\item \textbf{Wccftech (2026).} \textit{NVIDIA's official statement on the Vera Rubin spec change (October 29, 2026)}. \url{https://wccftech.com/nvidia-official-statement-on-vera-rubin-spec-change/}. Press, secondary source. NVIDIA's statement that the earlier figures were preliminary architectural targets. Slide 25.
\item \textbf{TechPowerUp (2026).} \textit{NVIDIA cuts HBM4 speed and NVLink bandwidth targets for Vera Rubin}. \url{https://www.techpowerup.com/354133/nvidia-cuts-hbm4-speed-and-nvlink-bandwidth-targets-for-vera-rubin}. Press, secondary source. The pin-rate arithmetic behind the revision.
\item \textbf{Hot Chips (2026).} \textit{Hot Chips 38 program (August 2026)}. \url{https://hc2026.hotchips.org/program/}. Conference program. The Rubin GPU talk (Mandal, Dash, Dimitrov) and the Vera CPU talk (Evans, Xekalakis). Slides 23-24.
\item \textbf{ServeTheHome (2026).} \textit{NVIDIA Vera Rubin NVL72 rack at Hot Chips 2026}. \url{https://www.servethehome.com/nvidia-vera-rubin-nvl72-rack-at-hot-chips-2026/}. Press, secondary source. The rack as presented: NVLink 6 switch with 130 TFLOP/s of in-network compute, 800 V DC, cable-free trays. Slide 25.
\item \textbf{ServeTheHome (2026).} \textit{NVIDIA Vera CPU at Hot Chips 2026}. \url{https://www.servethehome.com/nvidia-vera-cpu-at-hot-chips-2026/}. Press, secondary source. The Olympus core, spatial multithreading, the six-die package. Slide 24.
\item \textbf{NVIDIA (2026).} \textit{Vera Rubin ramps into full production (press release, May 31, 2026)}. \url{https://nvidianews.nvidia.com/news/vera-rubin-full-production-agentic-ai-factory}. Vendor press release. The GTC Taipei statement: production shipments starting in the fall. Slide 23.
\item \textbf{NVIDIA (2026).} \textit{Q2 fiscal 2027 financial results (8-K, August 26, 2026)}. \url{https://www.sec.gov/Archives/edgar/data/0001045810/000104581026000073/q2fy27pr.htm}. Regulatory filing. Vera Rubin racks at CoreWeave, Google Cloud, Azure, OCI and Nebius; Groq 3 LPX in production. Slide 23.
\item \textbf{Tom's Hardware (2026).} \textit{Nvidia expects to sell \$20 billion of Vera Rubin hardware this quarter}. \url{https://www.tomshardware.com/tech-industry/artificial-intelligence/nvidia-expects-to-sell-usd20-billion-worth-of-vera-rubin-hardware-this-quarter-would-account-for-20-percent-of-data-center-revenue-its-fastest-ramp-in-company-history}. Press, secondary source. First revenue racks in August 2026 and the Q3 FY27 guidance. Slide 23.
\item \textbf{CoreWeave (2026).} \textit{CoreWeave delivers NVIDIA Vera Rubin NVL72 performance at production scale, starting with Cognition (press release, September 2026)}. \url{https://investors.coreweave.com/news/news-details/2026/CoreWeave-Delivers-NVIDIA-Vera-Rubin-NVL72-Performance-at-Production-Scale-Starting-With-Cognition/default.aspx}. Vendor press release. The first production Vera Rubin cluster and CoreWeave's performance claims. Slide 23.
\item \textbf{NVIDIA (2025).} \textit{NVIDIA unveils NVLink Fusion (press release, May 2025)}. \url{https://investor.nvidia.com/news/press-release-details/2025/NVIDIA-Unveils-NVLink-Fusion-for-Industry-to-Build-Semi-Custom-AI-Infrastructure-With-NVIDIA-Partner-Ecosystem/default.aspx}. Vendor press release. NVLink opened to partners' CPUs and accelerators; the list of partners. Slide 26.
\item \textbf{Constellation Research (2025).} \textit{Nvidia's Groq deal: acquisition, acquihire or licensing? (December 2025)}. \url{https://www.constellationr.com/insights/news/nvidias-groq-deal-acquisition-acquihire-or-creative-licensing-deal}. Analyst note, secondary source. The structure of the December 2025 Groq transaction behind the Groq 3 LPX. Slide 26.
\item \textbf{NVIDIA (2025).} \textit{Introducing NVIDIA Dynamo, a low-latency distributed inference framework (developer blog, March 2025)}. \url{https://developer.nvidia.com/blog/introducing-nvidia-dynamo-a-low-latency-distributed-inference-framework-for-scaling-reasoning-ai-models}. Vendor blog. Disaggregated prefill and decode, KV-cache routing: the software premise of CPX and the LPX. Slides 26 and 29.
\end{itemize}
\subsection*{8. Systems and scale (slides 27-28)}
\begin{itemize}[leftmargin=1.2em,itemsep=3pt]\sloppy
\item \textbf{TOP500 (2026).} \textit{LineShine debuts at No. 1 as TOP500 enters a new global exascale era (June 2026)}. \url{https://www.top500.org/news/lineshine-debuts-no-1-top500-enters-new-global-exascale-era/}. Benchmark list, open. The June 2026 top ten with accelerators and power. Slide 27.
\item \textbf{TOP500 (2025).} \textit{El Capitan retains No. 1 as JUPITER becomes Europe's first exascale system (November 2025)}. \url{https://top500.org/news/el-capitan-retains-1-as-jupiter-becomes-europes-first-exascale-system-in-the-66th-top500-list/}. Benchmark list, open. JUPITER's 1.000 EFLOP/s entry. Slide 27.
\item \textbf{Forschungszentrum Juelich (2025).} \textit{Europe's first supercomputer reaches 1 exaflop/s (press release, November 2025)}. \url{https://www.fz-juelich.de/en/news/archive/press-release/2025/europes-first-supercomputer-reaches-1-exaflop-s}. Laboratory press release. JUPITER: about 24,000 GH200, the Booster configuration. Slide 27.
\item \textbf{Lawrence Berkeley National Laboratory (2026).} \textit{Berkeley Lab takes major step toward Doudna with delivery of early access system Cech (March 2026)}. \url{https://newscenter.lbl.gov/2026/03/17/berkeley-lab-takes-major-step-toward-doudna-with-delivery-of-early-access-system-cech/}. Laboratory news. Doudna (Dell, Vera Rubin, late 2026) and the Cech system (72 Grace, 144 Blackwell). Slide 27.
\item \textbf{NERSC (2025).} \textit{Doudna system (NERSC-10)}. \url{https://www.nersc.gov/what-we-do/resources/doudna-system}. Laboratory page. The official description of NERSC's Vera Rubin system.
\item \textbf{NVIDIA (2025).} \textit{Blue Lion supercomputer at LRZ will run on NVIDIA Vera Rubin (blog)}. \url{https://blogs.nvidia.com/blog/blue-lion-vera-rubin/}. Vendor blog. The LRZ (HPE) Vera Rubin system.
\item \textbf{NVIDIA (2025).} \textit{NVIDIA, Oracle and the US Department of Energy to build the DOE's largest AI supercomputer (press release, October 28, 2025)}. \url{https://nvidianews.nvidia.com/news/nvidia-oracle-us-department-of-energy-ai-supercomputer-scientific-discovery}. Vendor press release. Solstice (100,000 Blackwell GPUs) and Equinox (10,000) at Argonne. Slide 27.
\item \textbf{US Department of Energy (2025).} \textit{Energy Department announces new partnership with NVIDIA and Oracle to build largest DOE AI supercomputer}. \url{https://www.energy.gov/articles/energy-department-announces-new-partnership-nvidia-and-oracle-build-largest-doe-ai}. Government press release. The DOE side of the Solstice and Equinox announcement.
\item \textbf{Data Center Dynamics (2025).} \textit{OpenAI and Oracle to deploy 450,000 GB200 GPUs at Stargate Abilene}. \url{https://www.datacenterdynamics.com/en/news/openai-and-oracle-to-deploy-450000-gb200-gpus-at-stargate-abilene-data-center/}. Press, secondary source. The scale and power (1.2 GW) of the largest announced GB200 site. Slides 22 and 27.
\item \textbf{NVIDIA (2026).} \textit{DGX SuperPOD with DGX Vera Rubin NVL72 (blog)}. \url{https://blogs.nvidia.com/blog/dgx-superpod-rubin/}. Vendor blog. The eight-rack, 576-GPU Vera Rubin pod.
\item \textbf{NVIDIA (2025).} \textit{NVIDIA Blackwell Ultra sets records in MLPerf Training v5.1 (NVIDIA blog, November 2025)}. \url{https://blogs.nvidia.com/blog/mlperf-training-benchmark-blackwell-ultra/}. Vendor blog. The first NVFP4 training submission in MLPerf (Llama 3.1 405B on GB300), the chronology behind slide 12's adoption line. Slides 12 and 20.
\item \textbf{NVIDIA (2026).} \textit{NVIDIA Blackwell tops MLPerf Training 6.0 (developer blog, June 2026)}. \url{https://developer.nvidia.com/blog/nvidia-blackwell-tops-mlperf-training-6-0-with-industry-leading-scale-and-performance/}. Vendor blog. The 8,192-GPU GB300 results (DeepSeek-V3 671B in 2.02 minutes) and the GB300-to-GB200 ratios. Slide 20.
\item \textbf{MLCommons (2026).} \textit{MLPerf Training v6.0 supplemental results discussion (June 2026)}. \url{https://mlcommons.org/wp-content/uploads/2026/06/Final-MLPerf-Training-v6.0-Supplemental-Discussion-UNDER-EMBARGO-UNTIL-6_16_26-8_00-AM-PT.pdf}. Benchmark consortium document. The official analysis of the round.
\item \textbf{MLCommons (2026).} \textit{MLCommons releases new MLPerf Inference v6.0 benchmark results (April 2026)}. \url{https://www.streetinsider.com/Press+Releases/MLCommons+Releases+New+MLPerf+Inference+v6.0+Benchmark+Results/26256741.html}. Press release (mirror). The round's new benchmarks and largest systems.
\item \textbf{TrendForce (2026).} \textit{Rubin faces delays; Blackwell to exceed 70 percent of NVIDIA's 2026 high-end shipments (April 2026)}. \url{https://www.trendforce.com/presscenter/news/20260408-13003.html}. Analyst press release, secondary source. The 2026 product mix and the reasons given for Rubin's slower ramp.
\end{itemize}
\subsection*{9. Software (slide 29)}
\begin{itemize}[leftmargin=1.2em,itemsep=3pt]\sloppy
\item \textbf{NVIDIA (2025).} \textit{CUDA Toolkit 13.0 release notes}. \url{https://docs.nvidia.com/cuda/archive/13.0.0/cuda-toolkit-release-notes/index.html}. Vendor documentation. Architecture deprecations (Maxwell, Pascal, Volta) and the driver requirement. Slide 29.
\item \textbf{NVIDIA (2022).} \textit{CUDA Toolkit 12.0 release notes}. \url{https://docs.nvidia.com/cuda/archive/12.0.0/cuda-toolkit-release-notes/index.html}. Vendor documentation. Hopper and Ada enablement, clusters and the new PTX. Slide 29.
\item \textbf{NVIDIA (2022).} \textit{Transformer Engine (GitHub repository and documentation)}. \url{https://github.com/NVIDIA/TransformerEngine}. Vendor software, open source. The FP8 (and later MXFP8 and NVFP4) training library: recipes, delayed scaling, layer wrappers. Slides 12 and 29.
\item \textbf{Thakkar et al. (NVIDIA) (2023).} \textit{CUTLASS 3 and CuTe (GitHub repository and documentation)}. \url{https://github.com/NVIDIA/cutlass}. Vendor software, open source. The template library that exposes TMA, wgmma, clusters and tcgen05 to kernel writers. Slides 13 and 29.
\item \textbf{Patel, Nishball and Knuhtsen (SemiAnalysis) (2024).} \textit{MI300X vs H100 vs H200 benchmark part 1: training (December 2024)}. \url{https://newsletter.semianalysis.com/p/mi300x-vs-h100-vs-h200-benchmark-part-1-training}. Analyst report, secondary source. Achieved GEMM throughput against peak (about 720 of 989 TFLOP/s on the H100) and the software argument. Slide 29.
\item \textbf{NVIDIA (2025).} \textit{CUDA Python and cuTile (GTC 2025 announcements, developer blog)}. \url{https://developer.nvidia.com/blog/cuda-python-and-cutile/}. Vendor blog. The tile-based programming model and Python-native CUDA announced at GTC 2025. Slide 29.
\end{itemize}
\subsection*{10. Wrap-up, retrospectives and the spec-sheet record (slides 30-32)}
\begin{itemize}[leftmargin=1.2em,itemsep=3pt]\sloppy
\item \textbf{NVIDIA (2022).} \textit{NVIDIA H100 Tensor Core GPU architecture (official download page)}. \url{https://resources.nvidia.com/en-us-tensor-core/gtc22-whitepaper-hopper}. Vendor whitepaper (registration). The official copy of the Hopper whitepaper for citation.
\item \textbf{NVIDIA (2023).} \textit{H100 datasheet (official download page)}. \url{https://resources.nvidia.com/en-us-grace-cpu/h100-datasheet-24306-1}. Vendor datasheet (registration). The official copy of the shipping H100 specifications.
\item \textbf{Wikipedia (2026).} \textit{Rubin (microarchitecture)}. \url{https://en.wikipedia.org/wiki/Rubin_(microarchitecture)}. Tertiary source. A pointer page with dated citations; useful for finding primary sources, not for quoting.
\item \textbf{NVIDIA (2026).} \textit{CES 2026 keynote (video)}. \url{https://www.youtube.com/live/0NBILspM4c4}. Primary presentation. The Rubin slides as shown, including the transistor count that appears in no press release.
\item \textbf{HPCwire / AIwire (2026).} \textit{Huang shares NVIDIA roadmap: NVL1152, scale-up CPO (March 2026)}. \url{https://www.hpcwire.com/aiwire/2026/03/18/huang-shares-nvidia-roadmap-showing-more-chips-nvl1152-scale-up-cpo/}. Press, secondary source. The GTC 2026 roadmap through Feynman. Slide 26.
\item \textbf{SemiHub (2026).} \textit{Hot Chips 2026 recap: key presentations}. \url{https://semihub.io/en/blog/hot-chips-2026-recap.html}. Press, secondary source. Rubin and Vera beside AMD's MI400 in the same session; the comparison the next lecture takes up.
\end{itemize}
%=====
\clearpage
\printindex
\end{document}
